% Lesson 32 — Grammar: 1st, 2nd, and 3rd conditionals — Anglais Tle SH — Cours signé OnBuch+
% Programme officiel MINESEC. Compiler avec : tectonic EN32-grammar-1st-2nd-and-3rd-conditionals.tex
\newcommand{\DOCMATIERE}{Anglais}
\newcommand{\DOCNIVEAU}{Tle SH}
\newcommand{\DOCMODULE}{Chapter 3 : Environment, Well-being and Health}
\newcommand{\DOCLECON}{EN32}
\newcommand{\DOCTITRE}{Lesson 32 — Grammar: 1st, 2nd, and 3rd conditionals}
% OnBuch+ — préambule « cours signés » ENRICHI (compilé avec tectonic / XeLaTeX)
% Système visuel « Pop » d'OnBuch+ : fond crème (jamais blanc), bordures encre
% épaisses, ombres « dures » décalées, titres Archivo Black en capitales.
% Version MATHÉMATIQUES : graphes (pgfplots/tikz), boîtes pédagogiques
% (définition, propriété, méthode, attention, le savais-tu, exercice résolu,
% exercice, TP, bilan), page de garde et signature OnBuch+ ancrée en bas.
%
% Macros attendues dans main.tex AVANT \input{preamble} :
%   \DOCMATIERE  \DOCNIVEAU  \DOCMODULE  \DOCLECON (numéro)  \DOCTITRE
\documentclass[11pt]{article}

\usepackage{fontspec}
\usepackage[english,french]{babel} % français = langue principale ; l'anglais via \eng{..} / textebox
\usepackage[a4paper,margin=2cm,top=3.4cm,bottom=2.8cm,headheight=1.4cm,footskip=1.3cm]{geometry}
\usepackage{amsmath,amssymb,amsfonts}
\usepackage{mathtools} % \xmapsto, \coloneqq... souvent utilisés par les modèles
\usepackage{cancel} % simplifications barrées (bilans de forces)
% \abs{z} (module, valeur absolue) et \norm{u} (norme), utilisés par les modèles
\providecommand{\abs}{}\let\abs\relax\DeclarePairedDelimiter{\abs}{\lvert}{\rvert}
\providecommand{\norm}{}\let\norm\relax\DeclarePairedDelimiter{\norm}{\lVert}{\rVert}
\usepackage[table]{xcolor} % [table] : fournit aussi \rowcolors (alternance de lignes)
\usepackage{pagecolor}
\usepackage{tikz}
\usetikzlibrary{arrows.meta,positioning,calc,shapes.geometric,shapes.misc,shapes.arrows,decorations.pathmorphing,decorations.pathreplacing,decorations.markings,patterns,shadows,mindmap,trees,fit,backgrounds,matrix,intersections,angles,quotes}
\usepackage{pgfplots}
\usepackage[version=4]{mhchem} % équations nucléaires \ce{^{235}_{92}U}
\usepackage[european,siunitx]{circuitikz} % schémas électriques
\pgfplotsset{compat=1.18}
\usepgfplotslibrary{fillbetween,groupplots} % aires entre courbes (calcul intégral)
\usepackage{enumitem}
\usepackage{fancyhdr}
% [most] charge toutes les bibliothèques tcolorbox (listings, minted, raster,
% documentation...) et gonfle inutilement la mémoire TeX sur un document long
% (~300 boîtes) : on ne charge que ce qui est réellement utilisé.
\usepackage[skins,breakable]{tcolorbox}
\usepackage{graphicx}
\usepackage{colortbl}
\usepackage{booktabs}
\usepackage{tabularx}
\usepackage{diagbox} % case d'en-tête diagonale des tableaux à double entrée
% Colonnes extensibles centrée / gauche / droite (C, L, R), souvent utilisées par les modèles
\newcolumntype{C}{>{\centering\arraybackslash}X}
\newcolumntype{L}{>{\raggedright\arraybackslash}X}
\newcolumntype{R}{>{\raggedleft\arraybackslash}X}
\usepackage{array}
\usepackage{multicol}
\usepackage{multirow}
\usepackage{siunitx}
% parse-numbers=false : accepte aussi \SI{2,5 \times 10^{-3}}{...}
\sisetup{locale=FR,output-decimal-marker={,},group-minimum-digits=4,parse-numbers=false}
\DeclareSIUnit\ohm{\ensuremath{\symup{\Omega}}} % U+2126 absent de la police
\DeclareSIUnit\division{div} % réglages d'oscilloscope (ms/div, V/div)
\DeclareSIUnit\tr{tr} % tours (vitesse de rotation, ex. \SI{1500}{\tr\per\minute})
\DeclareSIUnit\molar{M} % concentration molaire (ex. \SI{3}{\molar}, \SI{1}{\micro\molar})
\DeclareSIUnit\an{an} % année (le modèle confond parfois avec \year, un compteur TeX) — ex. \SI{5}{\kilo\meter\per\an}
\DeclareSIUnit\FCFA{FCFA} % franc CFA (ex. \SI{500}{\FCFA\per\kilo\gram})
\DeclareSIUnit\degreeBrix{\textdegree Brix} % teneur en sucre d'un sirop/jus (réfractométrie)
\DeclareSIUnit\month{mois} % le modèle utilise parfois \month, un compteur TeX, comme unité de durée
\usepackage{qrcode}
% \degree : commande fréquemment utilisée par les modèles pour un angle ou une
% température en dehors de \SI{}{} (non fournie par les paquets chargés).
\providecommand{\degree}{^{\circ}}
% \qed (fin de démonstration), utilisé par les modèles sans amsthm
\providecommand{\qed}{\hfill$\square$}
% \barre : double barre des valeurs interdites dans un tableau de variations (tkz-tab)
\providecommand{\barre}{\ensuremath{\|}}
% environnement proof (amsthm non chargé) utilisé par les modèles
\ifdefined\proof\else\newenvironment{proof}[1][Démonstration]{\par\noindent\textit{#1.}\ }{\qed\par}\fi
\usepackage{lastpage}
\usepackage{hyperref}
\hypersetup{hidelinks}

% ============================================================
% Palette « Pop » OnBuch+ — mêmes hex que la classe P (plus_ui.dart)
% ============================================================
\definecolor{poporange}{HTML}{F59321}
\definecolor{popdark}{HTML}{E07A0C}
\definecolor{popcream}{HTML}{FFF6E8}
\definecolor{popink}{HTML}{1C1714}
\definecolor{poppurple}{HTML}{7A5AE0}
\definecolor{popgreen}{HTML}{2E9E6A}
\definecolor{poppink}{HTML}{E0457B}
\definecolor{popblue}{HTML}{2F7DE1}
\definecolor{popgold}{HTML}{C98A12}
\definecolor{popmuted}{HTML}{8A7F76}
\definecolor{popline}{HTML}{EADFD2}
\definecolor{popgray}{HTML}{8A7F76}
\colorlet{popgrey}{popgray}
% Alias tolérants (le modèle nomme parfois les couleurs différemment)
\colorlet{popwhite}{white}
\colorlet{popblack}{popink}
\colorlet{popred}{poppink}
\colorlet{popyellow}{popgold}
% Teintes claires (fonds de tableaux/encadrés) — le modèle nomme parfois
% les couleurs avec un suffixe L (ex. popblueL) sans les avoir définies.
\colorlet{poporangeL}{poporange!20}
\colorlet{popdarkL}{popdark!20}
\colorlet{poppurpleL}{poppurple!20}
\colorlet{popgreenL}{popgreen!20}
\colorlet{poppinkL}{poppink!20}
\colorlet{popblueL}{popblue!20}
\colorlet{popgoldL}{popgold!20}
\colorlet{popredL}{popred!20}
\colorlet{popgrayL}{popgray!20}
% Teintes claires pour fonds de tableaux / zones
\colorlet{popblueL}{popblue!10!white}
\colorlet{poporangeL}{poporange!14!white}
\colorlet{popgreenL}{popgreen!12!white}
\colorlet{poppurpleL}{poppurple!12!white}
\colorlet{poppinkL}{poppink!12!white}
\colorlet{popgoldL}{popgold!14!white}

\pagecolor{popcream}

% ============================================================
% Typographie
% ============================================================
\defaultfontfeatures{Ligatures=TeX}
\setmainfont{PlusJakartaSans-Medium.ttf}[Path=../../../fonts/,BoldFont=PlusJakartaSans-Bold.ttf]
% Maths en sans-serif (Fira Math) avec chiffres et lettres droites de Plus
% Jakarta, pour que les formules se fondent dans le texte.
\usepackage[math-style=ISO,bold-style=ISO]{unicode-math}
\setmathfont{FiraMath-Regular.otf}[Path=../../../fonts/]
\setmathfont{PlusJakartaSans-Medium.ttf}[Path=../../../fonts/,range={up/num,up/{Latin,latin},\mathup}]
\usepackage{icomma} % virgule décimale sans espace en mode math
% Caractères mathématiques Unicode absents des polices de texte (Plus Jakarta,
% Archivo Black) : rendus par leur équivalent mathématique, en texte comme en
% formule — sinon ils s'affichent en carré vide (ex. « Arithmétique dans ℤ »).
\usepackage{newunicodechar}
\newunicodechar{ℝ}{\ensuremath{\mathbb{R}}}
\newunicodechar{ℤ}{\ensuremath{\mathbb{Z}}}
\newunicodechar{ℂ}{\ensuremath{\mathbb{C}}}
\newunicodechar{ℕ}{\ensuremath{\mathbb{N}}}
\newunicodechar{ℚ}{\ensuremath{\mathbb{Q}}}
\newunicodechar{𝔻}{\ensuremath{\mathbb{D}}}
\newunicodechar{α}{\ensuremath{\alpha}}
\newunicodechar{β}{\ensuremath{\beta}}
\newunicodechar{γ}{\ensuremath{\gamma}}
\newunicodechar{δ}{\ensuremath{\delta}}
\newunicodechar{ε}{\ensuremath{\varepsilon}}
\newunicodechar{θ}{\ensuremath{\theta}}
\newunicodechar{λ}{\ensuremath{\lambda}}
\newunicodechar{μ}{\ensuremath{\mu}}
\newunicodechar{σ}{\ensuremath{\sigma}}
\newunicodechar{τ}{\ensuremath{\tau}}
\newunicodechar{φ}{\ensuremath{\varphi}}
\newunicodechar{ψ}{\ensuremath{\psi}}
\newunicodechar{ω}{\ensuremath{\omega}}
\newunicodechar{Σ}{\ensuremath{\Sigma}}
% majuscules produites par \MakeUppercase dans les titres (α-aminés -> Α)
\newunicodechar{Α}{\ensuremath{\alpha}}
\newunicodechar{Β}{\ensuremath{\beta}}
\newunicodechar{Γ}{\ensuremath{\Gamma}}
\newunicodechar{Δ}{\ensuremath{\Delta}}
\newunicodechar{Ε}{\ensuremath{\varepsilon}}
\newunicodechar{Θ}{\ensuremath{\Theta}}
\newunicodechar{Λ}{\ensuremath{\Lambda}}
\newunicodechar{Μ}{\ensuremath{\mu}}
\newunicodechar{Τ}{\ensuremath{\tau}}
\newunicodechar{Φ}{\ensuremath{\Phi}}
\newunicodechar{Ψ}{\ensuremath{\Psi}}
\newunicodechar{Ω}{\ensuremath{\Omega}}
\newunicodechar{↦}{\ensuremath{\mapsto}}
\newunicodechar{∈}{\ensuremath{\in}}
\newunicodechar{∉}{\ensuremath{\notin}}
\newunicodechar{∧}{\ensuremath{\wedge}}
\newunicodechar{⇔}{\ensuremath{\Leftrightarrow}}
\newunicodechar{⟺}{\ensuremath{\Longleftrightarrow}}
\newunicodechar{⇒}{\ensuremath{\Rightarrow}}
\newunicodechar{⟹}{\ensuremath{\Longrightarrow}}
\newunicodechar{∘}{\ensuremath{\circ}}
\newunicodechar{∪}{\ensuremath{\cup}}
\newunicodechar{∩}{\ensuremath{\cap}}
\newunicodechar{≡}{\ensuremath{\equiv}}
\newunicodechar{⊂}{\ensuremath{\subset}}
\newunicodechar{⊥}{\ensuremath{\perp}}
\newunicodechar{∃}{\ensuremath{\exists}}
\newunicodechar{∀}{\ensuremath{\forall}}
\newunicodechar{∛}{\ensuremath{\sqrt[3]{\,}}}
\newunicodechar{∜}{\ensuremath{\sqrt[4]{\,}}}
\newunicodechar{⃗}{\ensuremath{{}^{\to}}}
\newunicodechar{ⁿ}{\ifmmode^{n}\else\textsuperscript{\ensuremath{n}}\fi}
\newunicodechar{ˣ}{\ifmmode^{x}\else\textsuperscript{\ensuremath{x}}\fi}
\newunicodechar{ᵇ}{\ifmmode^{b}\else\textsuperscript{\ensuremath{b}}\fi}
\newunicodechar{ʳ}{\ifmmode^{r}\else\textsuperscript{\ensuremath{r}}\fi}
\newunicodechar{ᵘ}{\ifmmode^{u}\else\textsuperscript{\ensuremath{u}}\fi}
\newunicodechar{ʸ}{\ifmmode^{y}\else\textsuperscript{\ensuremath{y}}\fi}
\newunicodechar{ᵃ}{\ifmmode^{a}\else\textsuperscript{\ensuremath{a}}\fi}
\newunicodechar{ᵉ}{\ifmmode^{e}\else\textsuperscript{\ensuremath{e}}\fi}
\newunicodechar{ᵏ}{\ifmmode^{k}\else\textsuperscript{\ensuremath{k}}\fi}
\newunicodechar{ᵖ}{\ifmmode^{p}\else\textsuperscript{\ensuremath{p}}\fi}
\newunicodechar{ᴺ}{\ifmmode^{N}\else\textsuperscript{\ensuremath{N}}\fi}
\newunicodechar{ᶜ}{\ifmmode^{c}\else\textsuperscript{\ensuremath{c}}\fi}
\newunicodechar{ʰ}{\ifmmode^{h}\else\textsuperscript{\ensuremath{h}}\fi}
\newunicodechar{ᵠ}{\ifmmode^{q}\else\textsuperscript{\ensuremath{q}}\fi}
\newunicodechar{ₙ}{\ifmmode_{n}\else\textsubscript{\ensuremath{n}}\fi}
\newunicodechar{ₐ}{\ifmmode_{a}\else\textsubscript{\ensuremath{a}}\fi}
\newunicodechar{ᵢ}{\ifmmode_{i}\else\textsubscript{\ensuremath{i}}\fi}
\newunicodechar{ᵣ}{\ifmmode_{r}\else\textsubscript{\ensuremath{r}}\fi}
\newunicodechar{ₖ}{\ifmmode_{k}\else\textsubscript{\ensuremath{k}}\fi}
\newunicodechar{ᵦ}{\ifmmode_{\beta}\else\textsubscript{\ensuremath{\beta}}\fi}
\newunicodechar{ₓ}{\ifmmode_{x}\else\textsubscript{\ensuremath{x}}\fi}
\newfontfamily\popdisplay{ArchivoBlack-Regular.ttf}[Path=../../../fonts/]
\newfontfamily\poplogo{SpaceGrotesk-Bold.ttf}[Path=../../../fonts/]
\newfontfamily\popsemi{PlusJakartaSans-SemiBold.ttf}[Path=../../../fonts/]
\newfontfamily\popxbold{PlusJakartaSans-ExtraBold.ttf}[Path=../../../fonts/]
\newfontfamily\popxboldls{PlusJakartaSans-ExtraBold.ttf}[Path=../../../fonts/,LetterSpace=3.0]

\setlist[enumerate]{leftmargin=1.7em,itemsep=5pt,topsep=4pt}
\setlist[itemize]{leftmargin=1.7em,itemsep=4pt,topsep=4pt,label={\color{poporange}\textbullet}}
\setlength{\parindent}{0pt}
\setlength{\parskip}{5pt}
\renewcommand{\arraystretch}{1.25}

% pgfplots : style Pop par défaut
\pgfplotsset{
  every axis/.append style={axis line style={popink,line width=1pt},tick style={popink},
    grid style={popline,line width=0.5pt},label style={font=\small},tick label style={font=\footnotesize}},
  popcurve/.style={poporange,line width=1.8pt,no markers,smooth},
}
% Même style disponible dans un \draw TikZ simple (hors pgfplots)
\tikzset{popcurve/.style={poporange,line width=1.8pt}}
\tikzset{popgrid/.style={popline,very thin}}
\colorlet{popcurve}{poporange} % le nom du style est parfois utilisé comme couleur

% ============================================================
% Pill / squiggle
% ============================================================
\newcommand{\poppill}[3]{%
  \tikz[baseline=-0.55ex]\node[draw=popink,line width=1.2pt,rounded corners=2.6pt,%
    fill=#2,inner xsep=6.5pt,inner ysep=3pt]{%
    \popxboldls\fontsize{7}{7}\selectfont\color{#3}\MakeUppercase{#1}};%
}
\newcommand{\popsquiggle}[1]{%
  \tikz[baseline=-0.4ex]{
    \draw[#1,line width=2.6pt,line cap=round]
      plot [smooth] coordinates {(0,0.09) (0.34,0.01) (0.68,0.21) (1.02,0.01) (1.36,0.10)};
  }%
}
% Mot-clé surligné (marqueur orange sous le texte)
\newcommand{\cle}[1]{{\popxbold\color{popdark}#1}}

% ============================================================
% En-tête / pied
% ============================================================
\pagestyle{fancy}
\fancyhf{}
\renewcommand{\headrulewidth}{0pt}
\renewcommand{\footrulewidth}{0pt}
\lhead{\raisebox{-0.15ex}{\poplogo\fontsize{13}{13}\selectfont\color{popink}OnBuch\color{poporange}+}}
\rhead{\poppill{\DOCMATIERE\ \textbullet\ \DOCNIVEAU\ \textbullet\ Leçon \DOCLECON}{white}{popink}}
\lfoot{\popxbold\fontsize{7.6}{7.6}\selectfont\color{popmuted}\MakeUppercase{Cours signé OnBuch+}\ \textbullet\ {\popsemi\DOCTITRE}}
\rfoot{\tikz[baseline=-0.6ex]\node[rounded corners=8.5pt,fill=popink,minimum height=17pt,minimum width=17pt,inner xsep=5pt,inner sep=0]{\popxbold\fontsize{8}{8}\selectfont\color{popcream}\thepage\,/\,\pageref*{LastPage}};}

% ============================================================
% Titres
% ============================================================
\newcommand{\chaptitre}[2]{%
  \begin{tcolorbox}[enhanced,colback=poporange,colframe=popink,boxrule=2.6pt,arc=11pt,
      drop shadow=popink,left=15pt,right=15pt,top=12pt,bottom=11pt,before skip=3pt,after skip=18pt]
    \poppill{#2}{popink}{popcream}\par\vspace{9pt}
    {\raggedright\hyphenpenalty=10000\exhyphenpenalty=10000\popdisplay\fontsize{22}{24}\selectfont\color{popcream}\MakeUppercase{#1}\par}
  \end{tcolorbox}
}
% Section numérotée : pastille orange avec le numéro + titre Archivo
\newcounter{popsec}
\newcommand{\coursec}[1]{%
  \stepcounter{popsec}\setcounter{popsub}{0}%
  \par\vspace{16pt}\noindent
  \begin{minipage}{\linewidth}
  \tikz[baseline=-0.7ex]\node[draw=popink,line width=1.6pt,fill=poporange,rounded corners=4pt,minimum size=22pt,inner sep=0]{\popdisplay\fontsize{12}{12}\selectfont\color{popcream}\thepopsec};%
  \hspace{8pt}{\popdisplay\fontsize{15}{17}\selectfont\color{popink}\MakeUppercase{#1}}\par
  \vspace{4pt}\noindent{\color{poporange}\rule{36pt}{3.6pt}}
  \end{minipage}\par\nopagebreak\vspace{8pt}
}
\newcounter{popsub}
\newcommand{\courssub}[1]{%
  \stepcounter{popsub}%
  \par\vspace{9pt}\noindent{\popxbold\fontsize{11.5}{13}\selectfont\color{popink}{\color{poporange}\thepopsec.\thepopsub}\hspace{6pt}#1}\par\nopagebreak\vspace{4pt}
}

% ============================================================
% Boîtes pédagogiques — même recette : fond blanc, bordure épaisse colorée,
% ombre dure encre, badge plein. Titre optionnel [#1].
% ============================================================
\tcbset{popbase/.style={breakable,enhanced,colback=white,boxrule=2.3pt,arc=9pt,
  shadow={3pt}{-3pt}{0pt}{popink},left=14pt,right=14pt,top=11pt,bottom=11pt,before skip=11pt,after skip=15pt,
  fontupper=\popsemi\fontsize{10.6}{14.5}\selectfont\color{popink},
  fontlower=\popsemi\fontsize{10.6}{14.5}\selectfont\color{popink}}}
% Coche dessinée (le glyphe ✓ n'existe pas dans les polices du document)
\newcommand{\popcheck}[1]{\tikz[baseline=-0.5ex]\draw[#1,line width=1.6pt,line cap=round,line join=round](-0.13,0.0)--(-0.04,-0.09)--(0.13,0.1);}
\providecommand{\coche}{\popcheck{popgreen}}
\providecommand{\resultat}[1]{\par\smallskip\noindent\fcolorbox{popgreen}{popgreenL}{\parbox{\dimexpr\linewidth-2\fboxsep-2\fboxrule\relax}{\textbf{Résultat.} #1}}\par\smallskip}
\newcommand{\popboxhead}[3]{\poppill{#1}{#2}{white}\ifx\relax#3\relax\else\hspace{6pt}{\popxbold\fontsize{10.6}{12}\selectfont\color{popink}#3}\fi\par\vspace{7pt}}

\newtcolorbox{aretenir}[1][]{popbase,colframe=popgreen,before upper={\popboxhead{À retenir}{popgreen}{#1}}}
% Texte en anglais (document de lecture, dialogue, extrait) : cadre
% d'étude. Un titre optionnel (source, auteur) s'affiche dans l'en-tête.
\newtcolorbox{textebox}[1][]{popbase,colframe=popblue,before upper={\popboxhead{Text}{popblue}{#1}\begin{otherlanguage}{english}},after upper={\end{otherlanguage}}}
% Marques ✓ / ✗ (forme correcte / fautive) souvent utilisées par les modèles
\providecommand{\cross}{\textcolor{poppink}{\ensuremath{\times}}}
\providecommand{\xmark}{\cross}\providecommand{\crossmark}{\cross}\providecommand{\cmark}{\ensuremath{\checkmark}}
% Ligne de réponse / justification à compléter (inventée par les modèles)
\providecommand{\justification}[1]{\par\noindent\textit{Justification~:} \dotfill\par}
% \textipa (package tipa non chargé) : la police ne couvre pas l'API -> simple passage
\providecommand{\textipa}[1]{#1}
% Soulignement (ulem non chargé) : \uline, \ul
\providecommand{\uline}[1]{\underline{#1}}\providecommand{\ul}[1]{\underline{#1}}
% \glossaire{\item ...} inventée par les modèles : liste de glossaire
\providecommand{\glossaire}[1]{\begin{itemize}#1\end{itemize}}
% \alert (beamer) inventée par les modèles : texte mis en évidence
\providecommand{\alert}[1]{\textbf{\textcolor{poppink}{#1}}}
% variantes ASCII d'ordinaux écrites en macro
\providecommand{\iere}{\textsuperscript{re}}\providecommand{\ieme}{\textsuperscript{e}}\providecommand{\ere}{\textsuperscript{re}}\providecommand{\eme}{\textsuperscript{e}}\providecommand{\er}{\textsuperscript{er}}\providecommand{\ie}{\textsuperscript{e}}
% Ordinaux français écrits en macro par les modèles : 1\ière, 3\ième
\newcommand{\ière}{\textsuperscript{re}}\newcommand{\ième}{\textsuperscript{e}}
% \exemple (inventée par les modèles) : étiquette « Exemple : »
\providecommand{\exemple}{\textbf{Exemple~:}\ }
% \square utilisable aussi hors mode math (cases à cocher)
\AtBeginDocument{\let\squareMath\square\renewcommand{\square}{\ensuremath{\squareMath}}}
% Mot ou phrase en anglais à l'intérieur d'un paragraphe français
\newcommand{\eng}[1]{\ifmmode\text{\textit{#1}}\else\begingroup\selectlanguage{english}\itshape #1\endgroup\fi}
\let\esp\eng % alias toléré
\newtcolorbox{exemplebox}[1][]{popbase,colframe=poporange,before upper={\popboxhead{Exemple}{poporange}{#1}}}
\newtcolorbox{definition}[1][]{popbase,colframe=popblue,before upper={\popboxhead{Définition}{popblue}{#1}}}
\newtcolorbox{propriete}[1][]{popbase,colframe=poppurple,before upper={\popboxhead{Propriété}{poppurple}{#1}}}
\newtcolorbox{methode}[1][]{popbase,colframe=popgold,before upper={\popboxhead{Méthode}{popgold}{#1}}}
\newtcolorbox{attention}[1][]{popbase,colframe=poppink,before upper={\popboxhead{Attention}{poppink}{#1}}}
\newtcolorbox{experience}[1][]{popbase,colframe=popblue,colback=popblueL,before upper={\popboxhead{Expérience}{popblue}{#1}}}
% « Le savais-tu ? » : carte violette pleine (contexte, culture, Cameroun)
\newtcolorbox{savaistu}[1][]{popbase,colback=poppurple,colframe=popink,
  fontupper=\popsemi\fontsize{10.4}{14.2}\selectfont\color{white},
  before upper={\poppill{Le savais-tu\,?}{white}{poppurple}\ifx\relax#1\relax\else\hspace{6pt}{\popxbold\color{white}#1}\fi\par\vspace{7pt}}}
% Exercice résolu : énoncé en haut, \tcblower puis la solution sur fond crème
\newcounter{popexo}\newcounter{popexr}
\newtcolorbox{exoresolu}[1][]{popbase,colframe=poporange,colbacklower=popcream,
  segmentation style={popink,line width=1.2pt,dashed},
  before upper={\stepcounter{popexr}\popboxhead{Exercice résolu \thepopexr}{poporange}{#1}},
  before lower={\poppill{Solution}{popink}{popcream}\par\vspace{6pt}}}
% Exercice (non résolu en place) — la correction est dans la partie Corrigés
\newtcolorbox{exercice}[1][]{popbase,colframe=popink,
  before upper={\stepcounter{popexo}\popboxhead{Exercice \thepopexo}{popink}{#1}}}
% Corrigé
\newtcolorbox{corrige}[1][]{popbase,colframe=popgreen,colback=popgreenL,
  before upper={\popboxhead{Corrigé}{popgreen}{#1}}}
% Objectifs / prérequis / bilan
\newtcolorbox{objectifs}{popbase,colframe=popink,colback=poporangeL,
  before upper={\popboxhead{Objectifs de la leçon}{popink}{}}}
\newtcolorbox{prerequis}{popbase,colframe=popink,colback=popblueL,
  before upper={\popboxhead{Prérequis}{popblue}{}}}
\newtcolorbox{bilan}{popbase,colframe=popink,colback=white,boxrule=2.8pt,
  before upper={\popboxhead{Fiche bilan}{poporange}{L'essentiel à connaître pour l'examen}}}
% Figure encadrée avec légende
\newtcolorbox{popfigure}[1][]{enhanced,breakable=false,colback=white,colframe=popline,boxrule=1.4pt,arc=8pt,
  left=10pt,right=10pt,top=10pt,bottom=8pt,before skip=10pt,after skip=12pt,halign=center,
  lower separated=false,halign lower=center,
  fontlower=\popsemi\fontsize{9}{11.5}\selectfont\color{popmuted},
  #1}
\newcommand{\legende}[1]{\par\vspace{6pt}{\popsemi\fontsize{9}{11.5}\selectfont\color{popmuted}#1}}

% ============================================================
% Signature OnBuch+
% ============================================================
\newcommand{\poplogoinline}[1]{{\poplogo\fontsize{#1}{#1}\selectfont\color{popink}OnBuch\color{poporange}+}}
\newcommand{\signatureonbuch}{%
  \begin{tcolorbox}[enhanced,colback=popink,colframe=popink,boxrule=2.2pt,arc=10pt,
      drop shadow=poporange,left=14pt,right=14pt,top=10pt,bottom=10pt,before skip=0pt,after skip=0pt]
    \tikz[baseline=-0.7ex]\node[circle,fill=popgreen,draw=popcream,line width=1.3pt,minimum size=22pt,inner sep=0]{\popcheck{white}};%
    \hspace{9pt}\begin{minipage}[c]{0.62\linewidth}
      {\popxbold\fontsize{12}{13}\selectfont\color{popcream}Signé\ \textbullet\ {\poplogo OnBuch\color{poporange}+}}\par\vspace{2pt}
      {\raggedright\popsemi\fontsize{8}{10.5}\selectfont\color{popline}\MakeUppercase{Équipe pédagogique OnBuch+}\\\MakeUppercase{Conforme au programme officiel MINESEC}\par}
    \end{minipage}\hfill
    {\poplogo\fontsize{22}{22}\selectfont\color{popcream}OnBuch\color{poporange}+}
  \end{tcolorbox}%
}

% ============================================================
% Page de garde
% #1 titre · #2 matière · #3 niveau · #4 module · #5 n° leçon · #6 durée ·
% #7 description / objectifs (texte ou itemize)
% ============================================================
\newcommand{\pagegarde}[7]{%
  \thispagestyle{empty}
  \newgeometry{a4paper,margin=1.7cm,top=1.6cm,bottom=1.4cm}%
  \begin{tikzpicture}[remember picture,overlay]
    \fill[poporange!18] ([xshift=-3.2cm,yshift=-3.6cm]current page.north east) circle (4.6cm);
    \fill[poppurple!14] ([xshift=2.4cm,yshift=7.6cm]current page.south west) circle (3.8cm);
  \end{tikzpicture}%
  {\poplogo\fontsize{30}{30}\selectfont\color{popink}OnBuch\color{poporange}+}\hfill
  \raisebox{0.6ex}{\poppill{Cours signé \textbullet\ Premium}{popink}{popcream}}\par
  \vspace{1pt}\popsquiggle{poppurple}\par
  \vspace{8pt}
  \begin{tcolorbox}[enhanced,colback=poporange,colframe=popink,boxrule=2.8pt,arc=13pt,
      drop shadow=popink,left=18pt,right=18pt,top=12pt,bottom=13pt,after skip=10pt]
    \poppill{#2 \textbullet\ #3}{popink}{popcream}\hspace{5pt}\poppill{#4}{white}{popink}\par\vspace{8pt}
    {\popdisplay\fontsize{14}{14}\selectfont\color{popink}LEÇON #5}\par\vspace{4pt}
    {\raggedright\hyphenpenalty=10000\exhyphenpenalty=10000\popdisplay\fontsize{25}{27}\selectfont\color{popcream}\MakeUppercase{#1}\par}
    \vspace{8pt}
    {\popxbold\fontsize{9.5}{9.5}\selectfont\color{popink}\MakeUppercase{Durée conseillée : #6}}
  \end{tcolorbox}
  \begin{tcolorbox}[enhanced,colback=white,colframe=popink,boxrule=2.2pt,arc=10pt,
      drop shadow=popink,left=16pt,right=16pt,top=10pt,bottom=10pt,after skip=8pt]
    \poppill{Dans cette leçon}{poppurple}{white}\par\vspace{5pt}
    \popsemi\fontsize{9.3}{12.6}\selectfont\color{popink}#7
  \end{tcolorbox}
  \noindent\begin{minipage}[t]{0.487\linewidth}
    \begin{tcolorbox}[enhanced,colback=popgreen,colframe=popink,boxrule=2.2pt,arc=10pt,
        drop shadow=popink,left=11pt,right=11pt,top=10pt,bottom=10pt,height=3.1cm,valign=center]
      \tcbox[colback=white,colframe=popink,boxrule=1.3pt,arc=5pt,boxsep=0pt,nobeforeafter,
        left=3pt,right=3pt,top=3pt,bottom=3pt,valign=center]{\qrcode[height=1.9cm]{https://play.google.com/store/apps/details?id=com.luvvix.onbuch&hl=fr}}%
      \hspace{8pt}\begin{minipage}[c]{0.5\linewidth}
        {\popdisplay\fontsize{11}{12.5}\selectfont\color{white}\MakeUppercase{Télécharge OnBuch}}\par\vspace{4pt}
        {\raggedright\popsemi\fontsize{8.4}{11}\selectfont\color{white}Gratuit \textbullet\ Google Play\\Tuteur IA, annales, concours, résultats\par}
      \end{minipage}
    \end{tcolorbox}
  \end{minipage}\hfill
  \begin{minipage}[t]{0.487\linewidth}
    \begin{tcolorbox}[enhanced,colback=poppurple,colframe=popink,boxrule=2.2pt,arc=10pt,
        drop shadow=popink,left=11pt,right=11pt,top=10pt,bottom=10pt,height=3.1cm,valign=center]
      \tikz[baseline=-0.7ex]\node[circle,fill=white,draw=popink,line width=1.3pt,minimum size=1.6cm,inner sep=0]{\popdisplay\fontsize{24}{24}\selectfont\color{poppurple}+};%
      \hspace{8pt}\begin{minipage}[c]{0.6\linewidth}
        {\popdisplay\fontsize{11}{12.5}\selectfont\color{white}\MakeUppercase{Passe à OnBuch+}}\par\vspace{4pt}
        {\raggedright\popsemi\fontsize{8.4}{11}\selectfont\color{white}Tous les cours signés, Tuteur IA illimité, fiches PDF — onglet OnBuch+ de l'app\par}
      \end{minipage}
    \end{tcolorbox}
  \end{minipage}
  \vspace{8pt}
  \popcontenu
  \vfill
  \signatureonbuch
  \restoregeometry
  \newpage
}

% Rangée « Ce cours contient » (page de garde)
\newcommand{\poptile}[3]{% #1 couleur #2 gros texte #3 légende
  \begin{tcolorbox}[enhanced,colback=#1,colframe=popink,boxrule=2pt,arc=9pt,shadow={2.5pt}{-2.5pt}{0pt}{popink},
      left=6pt,right=6pt,top=9pt,bottom=9pt,halign=center,valign=center,height=1.75cm,width=\linewidth,nobeforeafter]
    {\popdisplay\fontsize{12.5}{13}\selectfont\color{white}\MakeUppercase{#2}}\par\vspace{3pt}
    {\centering\popsemi\fontsize{7.8}{9.5}\selectfont\color{white}#3\par}
  \end{tcolorbox}}
\newcommand{\popcontenu}{%
  \par\noindent{\popxboldls\fontsize{7.5}{7.5}\selectfont\color{popmuted}CE COURS SIGNÉ CONTIENT}\par\vspace{7pt}
  \noindent\begin{minipage}[t]{0.235\linewidth}\poptile{popblue}{Cours}{complet, illustré, pas à pas}\end{minipage}\hfill
  \begin{minipage}[t]{0.235\linewidth}\poptile{poporange}{Méthodes}{et exercices résolus}\end{minipage}\hfill
  \begin{minipage}[t]{0.235\linewidth}\poptile{poppink}{Exercices}{type examen + corrigés}\end{minipage}\hfill
  \begin{minipage}[t]{0.235\linewidth}\poptile{popgreen}{Fiche bilan}{l'essentiel à retenir}\end{minipage}\par}

% Sommaire
\newtcolorbox{sommaire}{enhanced,colback=white,colframe=popink,boxrule=2.2pt,arc=10pt,
  drop shadow=popink,left=18pt,right=18pt,top=13pt,bottom=13pt,before skip=6pt,after skip=20pt,
  before upper={\poppill{Sommaire}{poppurple}{white}\par\vspace{9pt}\popsemi\fontsize{10.6}{14.5}\selectfont\color{popink}}}

% Signature de fin de cours — poussée tout en bas de la dernière page
\newcommand{\findecours}{%
  \par\vfill\nopagebreak
  \begin{center}\popsquiggle{poporange}\end{center}\vspace{4pt}
  \signatureonbuch
}
\AtBeginDocument{\def\llbracket{\mathopen{[\mkern-3.5mu[}}\def\rrbracket{\mathclose{]\mkern-3.5mu]}}} % intervalles d'entiers (glyphe ⟦ absent de la police)
% Glyphes absents de Fira Math : redéfinis à partir de glyphes disponibles
\AtBeginDocument{%
  \def\setminus{\mathbin{\backslash}}\let\smallsetminus\setminus
  \def\vdots{\mathord{\vbox{\baselineskip3.2pt\lineskiplimit0pt\kern4pt\hbox{.}\hbox{.}\hbox{.}}}}%
  \def\ddots{\mathinner{\mkern1mu\raise6.5pt\hbox{.}\mkern2mu\raise3.6pt\hbox{.}\mkern2mu\raise0.7pt\hbox{.}\mkern1mu}}%
  \def\triangle{\mathord{\scalebox{0.9}{$\symup{\Delta}$}}}%
  \def\nabla{\mathord{\raisebox{\depth}{\scalebox{1}[-1]{$\symup{\Delta}$}}}}%
  \def\checkmark{\popcheck{popdark}}%
  \def\star{\mathbin{\ast}}\def\bigstar{\mathord{\ast}}%
  \def\diamond{\mathbin{\scalebox{0.6}{$\Diamond$}}}%
  \def\bigcup{\mathop{\mathchoice{\raisebox{-0.3em}{\scalebox{1.7}{$\cup$}}}{\scalebox{1.25}{$\cup$}}{\cup}{\cup}}\displaylimits}%
  \def\bigcap{\mathop{\mathchoice{\raisebox{-0.3em}{\scalebox{1.7}{$\cap$}}}{\scalebox{1.25}{$\cap$}}{\cap}{\cap}}\displaylimits}%
  \def\lesseqqgtr{\mathrel{\lessgtr}}\def\bigoplus{\mathop{\oplus}\displaylimits}%
}
\providecommand{\ssi}{si et seulement si} % abréviation fréquente
\AtBeginDocument{\providecommand{\wideparen}{\overparen}} % arc de cercle

\begin{document}

\pagegarde{Lesson 32 — Grammar: 1st, 2nd, and 3rd conditionals}{Anglais}{Tle SH}{Chapter 3 : Environment, Well-being and Health}{EN32}{2 h}{Cette leçon de grammaire présente les trois conditionnels anglais (1st, 2nd, 3rd) : leur formation, leurs emplois et leurs différences de sens. À travers le thème de l'environnement et de la santé, les élèves apprennent à exprimer l'hypothèse réelle, irréelle et le regret, en évitant les erreurs typiques des francophones.
\vspace{4pt}
\begin{multicols}{2}\small
\begin{itemize}
\item À la fin de cette leçon, je sais identifier le 1st, le 2nd et le 3rd conditional dans un texte
\item À la fin de cette leçon, je sais construire correctement la structure If + présent / will + BV
\item À la fin de cette leçon, je sais construire la structure If + prétérit / would + BV
\item À la fin de cette leçon, je sais construire la structure If + past perfect / would have + participe passé
\item À la fin de cette leçon, je sais distinguer hypothèse réelle, hypothèse irréelle et regret sur le passé
\item À la fin de cette leçon, je sais utiliser 'If I were you' et les modaux alternatifs (could, might) dans les conditionnels
\end{itemize}\end{multicols}}

\begin{sommaire}
\begin{multicols}{2}
\begin{enumerate}
\item Observation et découverte
\item Le 1st conditional : l'hypothèse réelle
\item Le 2nd conditional : l'hypothèse irréelle
\item Le 3rd conditional : le regret sur le passé
\item Comparaison français/anglais et pièges
\item Entraînement guidé et production
\item Activité d'intégration
\item Méthodes et exercices résolus
\item Exercices
\item Corrigés des exercices
\item Fiche bilan
\end{enumerate}
\end{multicols}
\end{sommaire}

\begin{objectifs}
\begin{itemize}
\item À la fin de cette leçon, je sais identifier le 1st, le 2nd et le 3rd conditional dans un texte
\item À la fin de cette leçon, je sais construire correctement la structure If + présent / will + BV
\item À la fin de cette leçon, je sais construire la structure If + prétérit / would + BV
\item À la fin de cette leçon, je sais construire la structure If + past perfect / would have + participe passé
\item À la fin de cette leçon, je sais distinguer hypothèse réelle, hypothèse irréelle et regret sur le passé
\item À la fin de cette leçon, je sais utiliser 'If I were you' et les modaux alternatifs (could, might) dans les conditionnels
\item À la fin de cette leçon, je sais éviter les erreurs fréquentes des francophones (will après if, would dans les deux propositions)
\item À la fin de cette leçon, je sais employer les conditionnels dans des phrases et un court texte sur l'environnement
\end{itemize}
\end{objectifs}

\begin{prerequis}
\begin{itemize}
\item Le présent simple et le futur avec will (classe de Première)
\item Le prétérit simple et le past perfect
\item Les modaux would, could, might
\item Le participe passé des verbes réguliers et irréguliers
\item La notion de proposition subordonnée introduite par if
\end{itemize}
\end{prerequis}

\newpage

\coursec{Observation et découverte}

\courssub{Mise en situation : trois phrases, trois réalités}

Imagine : demain, tu te réveilles à Yaoundé et il pleut à torrents. Les gouttes frappent le toit en tôle, la rue se transforme en rivière, et les bendskins peinent dans la boue. Tu dis à ton ami : \eng{If it rains, I will stay at home.} Puis, un peu plus tard, tu rêves devant la télévision : \eng{If I were rich, I would plant trees all over the city.} Enfin, le soir, tu regrettes en pensant à tes résultats : \eng{If I had studied harder, I would have passed the Probatoire.}

Trois phrases. Trois structures différentes. Trois réalités différentes : une \cle{possibilité réelle} (la pluie est vraiment possible demain), un \cle{rêve irréel} (tu n'es pas riche), un \cle{regret sur le passé} (tu n'as pas assez travaillé, et c'est trop tard). Comment l'anglais exprime-t-il ces trois situations ? Pourquoi le temps du verbe change-t-il après \eng{if} ? C'est tout l'objet de cette leçon sur les \cle{conditionnels}, des structures absolument indispensables au Baccalauréat — en grammaire comme en composition — et pour parler de l'avenir de notre environnement.

\begin{aretenir}[Problématique de la leçon]
En anglais, le choix du temps après \eng{if} et dans la proposition principale n'est pas libre : il dépend du \textbf{degré de réalité} de l'hypothèse. Nous allons découvrir trois constructions : le \cle{1st conditional} (hypothèse réelle), le \cle{2nd conditional} (hypothèse irréelle au présent) et le \cle{3rd conditional} (hypothèse irréelle au passé).
\end{aretenir}

\courssub{Lecture d'un texte : le climat au Cameroun}

Lisons d'abord ce texte original sur l'environnement au Cameroun. Pendant la lecture, repère toutes les phrases qui contiennent le mot \eng{if} et observe le temps des verbes qui l'entourent.

\begin{textebox}[Saving our forests — a text about the Cameroonian climate]
Cameroon is one of the greenest countries in Africa, but its forests are in danger. Every year, farmers cut down trees to plant cocoa, and companies cut wood to sell it abroad. The consequences are already visible: the rainy season in Douala is shorter than before, and the dry season in the North is longer and hotter.

What can we do? The answer is simple. \underline{If we protect the forest, animals will survive}, and the climate will stay balanced. \underline{If students learn about the environment at school, they will become responsible adults.} Our government has a plan, but it needs everyone's help.

Some young people have big dreams. Amina, a student in Buea, says: \underline{If I were the Minister of the Environment, I would ban plastic bags} in all our markets. \underline{If I had more money, I would buy land and plant a thousand trees.} Her friend Roland adds: \underline{If people recycled their waste, our cities would be cleaner} and healthier.

Unfortunately, we have already made mistakes. \underline{If people had recycled more, the city would have been cleaner} ten years ago. \underline{If the government had protected the forests earlier, many species would not have disappeared.} We cannot change the past, but we can change the future. The choice is ours: will we act, or will we watch our beautiful country lose its green treasure?
\end{textebox}

\textbf{Glossaire :}

\begin{center}
\begin{tabularx}{\linewidth}{|l|l|X|}
\hline
\rowcolor{popblueL} \textbf{English} & \textbf{Nature} & \textbf{Français} \\
\hline
to cut down (trees) & verbe & abattre, couper (des arbres) \\
\hline
abroad & adverbe & à l'étranger \\
\hline
the rainy season & nom & la saison des pluies \\
\hline
balanced & adjectif & équilibré \\
\hline
to ban & verbe & interdire \\
\hline
waste & nom (indénombrable) & les déchets \\
\hline
a species & nom & une espèce \\
\hline
to disappear & verbe & disparaître \\
\hline
a treasure & nom & un trésor \\
\hline
\end{tabularx}
\end{center}

\textbf{Questions de compréhension (oral) :}

\begin{enumerate}
\item Why are Cameroon's forests in danger?
\item What would Amina do if she were the Minister of the Environment?
\item What mistakes does the text mention about the past?
\end{enumerate}

\courssub{Repérage guidé : observer les structures}

\begin{methode}[Comment observer une phrase conditionnelle]
Pour analyser une phrase avec \eng{if}, suis toujours ces quatre étapes :
\begin{enumerate}
\item \textbf{Repérer} le mot \eng{if} et souligner toute la proposition qui le contient (la \emph{subordonnée}).
\item \textbf{Noter le temps} du verbe après \eng{if} : present simple ? past simple ? past perfect ?
\item \textbf{Noter la forme} du verbe de la proposition principale : \eng{will} + verbe ? \eng{would} + verbe ? \eng{would have} + participe passé ?
\item \textbf{En déduire le sens} : l'hypothèse est-elle possible, imaginaire, ou impossible parce que passée ?
\end{enumerate}
\end{methode}

Appliquons cette méthode aux six phrases soulignées du texte. Complétons ensemble le tableau d'observation :

\begin{center}
\begin{tabularx}{\linewidth}{|X|l|l|}
\hline
\rowcolor{popblueL} \textbf{Phrase} & \textbf{Temps après if} & \textbf{Forme principale} \\
\hline
If we protect the forest, animals will survive. & present simple & will + verbe \\
\hline
If students learn, they will become responsible. & present simple & will + verbe \\
\hline
If I were the Minister, I would ban plastic bags. & past simple & would + verbe \\
\hline
If people recycled, cities would be cleaner. & past simple & would + verbe \\
\hline
If people had recycled more, the city would have been cleaner. & past perfect & would have + participe \\
\hline
If the government had protected the forests, species would not have disappeared. & past perfect & would have + participe \\
\hline
\end{tabularx}
\end{center}

\begin{exemplebox}[Deux exemples traduits pour bien sentir la différence]
\eng{If it rains, I will stay at home.} — « S'il pleut, je resterai à la maison. » (la pluie est une possibilité réelle : demain, à Yaoundé, cela peut vraiment arriver.)

\eng{If I were rich, I would plant trees.} — « Si j'étais riche, je planterais des arbres. » (ce n'est qu'un rêve : je ne suis pas riche, c'est imaginaire.)
\end{exemplebox}

\courssub{Question déclencheuse : quel changement de sens ?}

Comparons attentivement ces deux phrases :

\begin{itemize}
\item \eng{If it rains, I will stay.} — « S'il pleut, je resterai. »
\item \eng{If it rained, I would stay.} — « S'il pleuvait, je resterais. »
\end{itemize}

Dans la première, le locuteur regarde le ciel de Douala : les nuages sont noirs, la pluie est \textbf{probable}. Dans la seconde, il fait peut-être beau : la pluie est une pure \textbf{supposition}, une imagination. Remarque que le français fait exactement le même raisonnement : « s'il pleut » (fait possible) contre « s'il pleuvait » (conditionnel, irréel). L'anglais et le français fonctionnent ici de façon très proche : \textbf{plus le temps « recule » vers le passé, moins l'hypothèse est réelle}.

\courssub{Classement des exemples : trois groupes, trois degrés de réalité}

Classons maintenant les phrases du texte en trois groupes :

\begin{center}
\begin{tabularx}{\linewidth}{|l|X|X|}
\hline
\rowcolor{popblueL} \textbf{Groupe} & \textbf{Construction} & \textbf{Degré de réalité} \\
\hline
1st conditional & if + present simple \textrightarrow{} will + verbe & hypothèse \textbf{réelle, possible} (avenir) \\
\hline
2nd conditional & if + past simple \textrightarrow{} would + verbe & hypothèse \textbf{imaginaire, irréelle} (présent) \\
\hline
3rd conditional & if + past perfect \textrightarrow{} would have + participe & hypothèse \textbf{impossible} (passé, regret) \\
\hline
\end{tabularx}
\end{center}

\begin{exoresolu}[Classer des phrases dans le bon groupe]
Énoncé : classe chaque phrase dans le groupe qui convient (1st, 2nd ou 3rd conditional) et traduis-la.
\begin{enumerate}
\item \eng{If the rains come early, the farmers will be happy.}
\item \eng{If I lived in Bamenda, I would grow vegetables.}
\item \eng{If we had planted trees last year, the schoolyard would have been cooler.}
\end{enumerate}
\tcblower
Solution détaillée :
\begin{enumerate}
\item \eng{If the rains come} = if + present simple ; \eng{will be} = futur. C'est un \textbf{1st conditional} : la pluie peut vraiment arriver tôt. Traduction : « Si les pluies arrivent tôt, les agriculteurs seront contents. »
\item \eng{If I lived} = if + past simple ; \eng{would grow} = conditionnel. C'est un \textbf{2nd conditional} : je n'habite pas à Bamenda, c'est imaginaire. Traduction : « Si j'habitais à Bamenda, je cultiverais des légumes. »
\item \eng{If we had planted} = if + past perfect ; \eng{would have been} = conditionnel passé. C'est un \textbf{3rd conditional} : nous n'avons pas planté, c'est trop tard. Traduction : « Si nous avions planté des arbres l'année dernière, la cour de l'école aurait été plus fraîche. »
\end{enumerate}
\end{exoresolu}

\courssub{Dégagement collectif de la règle générale}

À partir de nos observations, nous pouvons formuler ensemble la règle fondamentale :

\begin{propriete}[La règle générale des conditionnels]
Dans une phrase conditionnelle anglaise, le temps de la subordonnée (après \eng{if}) et celui de la proposition principale sont \textbf{liés} :
\begin{itemize}
\item \eng{if} + \textbf{present simple} \textrightarrow{} \textbf{will} + base verbale : hypothèse réelle sur l'avenir.
\item \eng{if} + \textbf{past simple} \textrightarrow{} \textbf{would} + base verbale : hypothèse irréelle au présent.
\item \eng{if} + \textbf{past perfect} (\eng{had} + participe passé) \textrightarrow{} \textbf{would have} + participe passé : hypothèse irréelle sur le passé.
\end{itemize}
\textbf{Règle d'or :} après \eng{if}, on n'emploie \textbf{jamais} \eng{will} ni \eng{would} : le futur et le conditionnel se placent dans la proposition principale, pas dans la subordonnée.
\end{propriete}

\begin{attention}[Piège n°1 des francophones]
Ne traduis jamais mot à mot « si je \emph{serai} riche » par \eng{if I will be rich} ! En français on dit « si je suis riche, je serai content » : le futur n'apparaît que dans une seule proposition. En anglais, c'est pareil : \eng{If I am rich, I will be happy}, et jamais \eng{if I will be}. De même, évite \eng{if I would be rich} : la forme correcte est \eng{if I were rich}. Retiens aussi que dans le 2nd conditional, on préfère \eng{were} à \eng{was} pour toutes les personnes dans le style soigné : \eng{If I were you, I would apologise.} — « Si j'étais toi, je m'excuserais. »
\end{attention}

\begin{popfigure}
\begin{tikzpicture}[mindmap, grow cyclic, every node/.style={concept, align=center, font=\small},
root concept/.append style={concept color=popblue, font=\bfseries},
level 1/.append style={level distance=3.6cm, sibling angle=120},
level 2/.append style={level distance=2.6cm, sibling angle=45, font=\scriptsize}]
\node[root concept]{CONDITIONALS}
  child[concept color=popgreen]{ node[align=center]{1st\\ real future}
    child { node[align=center]{if + present\\ will + verb} }
    child { node[align=center]{If it rains,\\ I will stay.} }
  }
  child[concept color=poporange]{ node[align=center]{2nd\\ unreal present}
    child { node[align=center]{if + past\\ would + verb} }
    child { node[align=center]{If I were rich,\\ I would plant trees.} }
  }
  child[concept color=poppurple]{ node[align=center]{3rd\\ unreal past}
    child { node[align=center]{if + past perfect\\ would have + p.p.} }
    child { node[align=center]{If I had studied,\\ I would have passed.} }
  };
\end{tikzpicture}
\legende{Carte mentale : les trois conditionnels et leur degré de réalité.}
\end{popfigure}

\begin{exercice}[À toi de jouer]
Repère \eng{if} dans chaque phrase, note les deux temps utilisés, puis indique si l'hypothèse est réelle, imaginaire ou passée :
\begin{enumerate}
\item \eng{If we recycle plastic, Douala will be cleaner.}
\item \eng{If I were you, I would drink more water.}
\item \eng{If they had listened to the doctor, they would have felt better.}
\end{enumerate}
\end{exercice}

\begin{corrige}[Exercice : À toi de jouer]
\begin{enumerate}
\item \eng{recycle} = present simple ; \eng{will be} = futur. Hypothèse \textbf{réelle} (1st conditional) : « Si nous recyclons le plastique, Douala sera plus propre. »
\item \eng{were} = past simple ; \eng{would drink} = conditionnel. Hypothèse \textbf{imaginaire} (2nd conditional) : « Si j'étais toi, je boirais plus d'eau. »
\item \eng{had listened} = past perfect ; \eng{would have felt} = conditionnel passé. Hypothèse \textbf{passée, impossible} (3rd conditional) : « S'ils avaient écouté le médecin, ils se seraient sentis mieux. »
\end{enumerate}
\end{corrige}

\begin{aretenir}[Ce qu'il faut retenir de cette découverte]
L'anglais possède trois conditionnels qui correspondent à trois degrés de réalité. Le temps après \eng{if} « recule » d'un cran à chaque fois que l'hypothèse devient moins réelle : present \textrightarrow{} past \textrightarrow{} past perfect. Dans la proposition principale, on trouve respectivement \eng{will}, \eng{would}, \eng{would have}. Les sections suivantes étudieront chaque structure en détail : forme, emplois, exceptions et pièges.
\end{aretenir}

\coursec{Le 1st conditional : l'hypothèse réelle}

Dans la section précédente, nous avons observé un texte et un dialogue sur l'environnement dans lesquels apparaissaient des phrases comme \eng{If we protect the forest, animals will survive.} et \eng{If I had more time, I would plant trees.} Nous avons remarqué que toutes ces phrases contiennent le mot \cle{if} et expriment une \emph{condition}. Mais elles ne se construisent pas toutes de la même manière. Dans cette section, nous allons étudier en profondeur la première de ces structures : le \cle{1st conditional}, qui sert à exprimer une hypothèse \emph{réelle}, c'est-à-dire possible ou probable, dans le futur.

\courssub{Observation : un texte pour découvrir la structure}

Lisons d'abord ce court article de presse. Soulignez mentalement toutes les phrases qui contiennent \eng{if} et observez les temps utilisés dans chaque partie de la phrase.

\begin{textebox}[A warning from the Ministry of Environment — Yaoundé]
The rainy season is coming, and the Ministry of Environment has published a warning for the population of Yaoundé. “If people throw plastic bags in the gutters, the streets will flood,” explains Mrs Ngo Bell, an environmental engineer. “If the water does not flow, mosquitoes will multiply, and malaria will spread in the neighbourhoods.”

The message is clear: everyone can help. “If each family sorts its waste, the city will be cleaner. If you see litter in the street, pick it up. If we act together, we will avoid another disaster like last year's floods.”

Some citizens are optimistic. “If the government builds more drains, the situation will improve,” says a shopkeeper at Mokolo market. But others are worried: “Unless we change our habits now, the problem will get worse every year.”

The Ministry also reminds farmers in the surrounding villages. “If you burn the bushes, the soil will lose its fertility. If you plant trees instead, you will protect your land and your crops.”

Finally, a message for the young generation: “If you study environmental science, you can become part of the solution. Cameroon needs young people who care about nature.”
\end{textebox}

\textbf{Glossaire :}

\begin{center}
\begin{tabularx}{\linewidth}{|X|X|X|}
\hline
\rowcolor{popblueL} Mot anglais & Nature & Traduction \\
\hline
gutter & nom & caniveau, gouttière \\
\hline
to flood & verbe & inonder \\
\hline
to spread & verbe & se propager \\
\hline
litter & nom & déchets (dans la rue) \\
\hline
drain & nom & égout, canal d'évacuation \\
\hline
soil & nom & sol, terre \\
\hline
fertility & nom & fertilité \\
\hline
crops & nom pluriel & récoltes, cultures \\
\hline
to care about & verbe & se soucier de \\
\hline
\end{tabularx}
\end{center}

\textbf{Questions de compréhension :}
\begin{enumerate}
\item What will happen if people throw plastic bags in the gutters?
\item What does Mrs Ngo Bell advise people to do when they see litter?
\item According to the shopkeeper, what will improve the situation?
\item What advice does the Ministry give to farmers?
\end{enumerate}

\courssub{Analyse des exemples observés}

Reprenons quelques phrases du texte et décomposons-les :

\begin{itemize}
\item \eng{If people throw plastic bags, the streets will flood.} — \emph{Si les gens jettent des sachets plastiques, les rues seront inondées.}
\item \eng{If each family sorts its waste, the city will be cleaner.} — \emph{Si chaque famille trie ses déchets, la ville sera plus propre.}
\item \eng{If you burn the bushes, the soil will lose its fertility.} — \emph{Si vous brûlez les brousses, le sol perdra sa fertilité.}
\end{itemize}

Dans chaque phrase, on observe deux propositions :
\begin{enumerate}
\item la \cle{proposition subordonnée} introduite par \eng{if}, au \textbf{présent simple} : \eng{If people throw}, \eng{If each family sorts}, \eng{If you burn} ;
\item la \cle{proposition principale}, au \textbf{futur avec will + base verbale} : \eng{the streets will flood}, \eng{the city will be}, \eng{the soil will lose}.
\end{enumerate}

\begin{propriete}[Règle du 1st conditional : forme]
Le \cle{1st conditional} exprime une hypothèse \textbf{réelle, possible ou probable} dans le futur. Sa structure est :

\begin{center}
\textbf{If + sujet + présent simple, \quad sujet + will + base verbale (BV)}
\end{center}

\begin{itemize}
\item \eng{If it rains tomorrow, we will stay at home.} — \emph{S'il pleut demain, nous resterons à la maison.}
\item \eng{If she works hard, she will pass the Probatoire.} — \emph{Si elle travaille dur, elle réussira le Probatoire.}
\end{itemize}

\textbf{Règle d'or :} on ne met \textbf{jamais} \eng{will} dans la proposition introduite par \eng{if}. Le futur n'apparaît que dans la proposition principale.
\end{propriete}

\begin{popfigure}
\begin{center}
\begin{tikzpicture}[
box/.style={rectangle, rounded corners, draw=popblue, fill=popblueL, align=center, minimum height=1.2cm, font=\small},
arrow/.style={-{Stealth[length=3mm]}, very thick, poporange}]
\node[box, align=center] (cond) at (0,0){\textbf{Condition}\\ If + présent simple\\ \eng{If it rains,}};
\node[box, fill=popgreenL, draw=popgreen, align=center] (res) at (7.5,0){\textbf{Résultat futur}\\ will + base verbale\\ \eng{we will stay at home.}};
\draw[arrow] (cond) -- (res) node[midway, above, font=\small\itshape, text=popdark] {futur réel};
\end{tikzpicture}
\end{center}
\legende{Structure du 1st conditional : une condition au présent entraîne un résultat au futur.}
\end{popfigure}

\courssub{Emploi : quand utiliser le 1st conditional ?}

\begin{aretenir}[Emplois du 1st conditional]
On utilise le 1st conditional pour parler d'une situation \textbf{future} dont la condition est jugée \textbf{possible ou probable} par le locuteur :
\begin{itemize}
\item une \textbf{prédiction} : \eng{If we pollute the rivers, fish will die.} — \emph{Si nous polluons les rivières, les poissons mourront.}
\item un \textbf{avertissement} (warning) : \eng{If you don't sleep, you will be tired tomorrow.} — \emph{Si tu ne dors pas, tu seras fatigué demain.}
\item une \textbf{promesse} : \eng{If you help me, I will help you too.} — \emph{Si tu m'aides, je t'aiderai aussi.}
\item une \textbf{menace} : \eng{If you touch that wire, you will get an electric shock.} — \emph{Si tu touches ce fil, tu prendras une décharge électrique.}
\item une \textbf{offre ou un conseil} : \eng{If you feel ill, we will call the doctor.} — \emph{Si tu te sens malade, nous appellerons le médecin.}
\end{itemize}
Le locuteur estime que la condition a de \textbf{bonnes chances de se réaliser} : c'est ce qui distingue le 1st conditional du 2nd conditional (hypothèse irréelle), que nous étudierons dans la section suivante.
\end{aretenir}

\courssub{L'ordre des propositions : avec ou sans virgule}

Les deux propositions peuvent s'inverser librement, mais la ponctuation change :

\begin{propriete}[Inversion des propositions]
\begin{itemize}
\item Quand la proposition en \eng{if} est \textbf{au début}, on met une \textbf{virgule} entre les deux propositions : \eng{If it rains, I will stay at home.}
\item Quand la proposition principale est \textbf{au début}, il n'y a \textbf{pas de virgule} : \eng{I will stay at home if it rains.}
\end{itemize}
Le sens ne change pas ; seule la mise en relief diffère.
\end{propriete}

\begin{exemplebox}[Exemples traduits — inversion]
\begin{itemize}
\item \eng{If she recycles, she will help the planet.} = \eng{She will help the planet if she recycles.} — \emph{Si elle recycle, elle aidera la planète.}
\item \eng{If we plant trees, the air will be cleaner.} = \eng{The air will be cleaner if we plant trees.} — \emph{Si nous plantons des arbres, l'air sera plus pur.}
\item \eng{If you come to Douala, you will visit the port.} = \eng{You will visit the port if you come to Douala.} — \emph{Si tu viens à Douala, tu visiteras le port.}
\end{itemize}
\end{exemplebox}

\courssub{Variantes : les modaux à la place de will}

Dans la proposition principale, \eng{will} peut être remplacé par d'autres \textbf{auxiliaires modaux}, selon le sens voulu :

\begin{center}
\begin{tabularx}{\linewidth}{|l|X|X|}
\hline
\rowcolor{popblueL} Modal & Exemple anglais & Traduction et nuance \\
\hline
\eng{will} & \eng{If you study, you will pass.} & Si tu étudies, tu réussiras. (certitude) \\
\hline
\eng{can} & \eng{If you study, you can pass.} & Si tu étudies, tu peux réussir. (capacité, possibilité) \\
\hline
\eng{may} & \eng{If it rains, the match may be cancelled.} & S'il pleut, le match pourrait être annulé. (éventualité) \\
\hline
\eng{might} & \eng{If we hurry, we might catch the bus.} & Si on se dépêche, on pourrait attraper le bus. (possibilité plus faible) \\
\hline
\eng{should} & \eng{If you feel sick, you should see a doctor.} & Si tu te sens malade, tu devrais voir un médecin. (conseil) \\
\hline
\eng{must} & \eng{If you drive, you must wear a seatbelt.} & Si tu conduis, tu dois porter la ceinture. (obligation) \\
\hline
\end{tabularx}
\end{center}

\begin{exemplebox}[Exemples traduits — modaux]
\begin{itemize}
\item \eng{If you study environmental science, you can become part of the solution.} — \emph{Si tu étudies les sciences de l'environnement, tu peux faire partie de la solution.}
\item \eng{If the rain continues, the Wouri river may overflow.} — \emph{Si la pluie continue, le fleuve Wouri pourrait déborder.}
\item \eng{If we leave now, we might arrive before the storm.} — \emph{Si nous partons maintenant, nous arriverons peut-être avant l'orage.}
\end{itemize}
\end{exemplebox}

\courssub{L'impératif dans la proposition principale}

Pour donner un \textbf{ordre}, un \textbf{conseil} ou une \textbf{instruction} sous condition, on remplace la proposition principale par un \textbf{impératif} :

\begin{propriete}[If + présent, impératif]
\begin{center}
\textbf{If + sujet + présent simple, \quad impératif (base verbale)}
\end{center}
\begin{itemize}
\item \eng{If you see litter, pick it up.} — \emph{Si tu vois des déchets, ramasse-les.}
\item \eng{If you feel tired, take a rest.} — \emph{Si tu te sens fatigué, repose-toi.}
\item \eng{If the tap leaks, call the plumber.} — \emph{Si le robinet fuit, appelle le plombier.}
\item \eng{If you go to Buea, don't forget your umbrella.} — \emph{Si tu vas à Buea, n'oublie pas ton parapluie.}
\end{itemize}
\end{propriete}

\courssub{Unless = if... not}

Le mot \cle{unless} signifie \eng{if... not} (« à moins que... ne », « sauf si »). Il introduit une condition \textbf{négative} sans utiliser de négation dans le verbe qui le suit.

\begin{propriete}[Unless]
\eng{Unless + présent simple} = \eng{if + sujet + not + présent simple}.
\begin{itemize}
\item \eng{Unless we act now, the climate will get worse.} = \eng{If we do not act now, the climate will get worse.} — \emph{À moins que nous n'agissions maintenant, le climat empirera.}
\item \eng{Unless it rains, the crops will fail.} = \eng{If it does not rain, the crops will fail.} — \emph{À moins qu'il ne pleuve, les récoltes seront mauvaises.}
\item \eng{You will not pass unless you work harder.} — \emph{Tu ne réussiras pas à moins de travailler davantage.}
\end{itemize}
\textbf{Attention :} après \eng{unless}, le verbe est à la forme \textbf{affirmative} (la négation est déjà contenue dans \eng{unless}).
\end{propriete}

\courssub{Comparaison avec le français et pièges}

\begin{center}
\begin{tabularx}{\linewidth}{|X|X|X|}
\hline
\rowcolor{popblueL} Français & English & Remarque \\
\hline
S'il pleut, je resterai à la maison. & \eng{If it rains, I will stay at home.} & Présent après \emph{si} dans les deux langues. \\
\hline
S'il pleuvait... (futur après si : impossible) & \eng{If it will rain...} (fautif) & Jamais \eng{will} après \eng{if}. \\
\hline
Si tu étudies, tu peux réussir. & \eng{If you study, you can pass.} & Modal possible dans la principale. \\
\hline
À moins qu'il ne pleuve, les récoltes seront mauvaises. & \eng{Unless it rains, the crops will fail.} & \eng{Unless} + forme affirmative. \\
\hline
\end{tabularx}
\end{center}

\begin{attention}[Erreurs fréquentes des francophones]
\begin{enumerate}
\item \textbf{Le piège majeur :} en français, on peut être tenté de mettre le futur après « si » par analogie avec la phrase principale. En anglais, c'est une faute grave : ne dites jamais \eng{If it will rain, I will stay at home.} Dites : \eng{If it rains, I will stay at home.}
\item Ne confondez pas \eng{if} (condition) et \eng{when} (certitude temporelle) : \eng{When I arrive, I will call you} (je suis sûr d'arriver) ≠ \eng{If I arrive early, I will call you} (ce n'est pas certain).
\item N'oubliez pas la \textbf{virgule} quand la phrase commence par \eng{if} : \eng{If it rains, ...} et non \eng{If it rains I will...}.
\item Après \eng{unless}, ne rajoutez pas de négation : \eng{Unless you don't work} est fautif ; dites \eng{Unless you work}.
\item La base verbale après \eng{will} ne prend jamais de \eng{-s} ni de \eng{to} : \eng{she will helps} et \eng{she will to help} sont fautifs ; dites \eng{she will help}.
\item Attention au faux ami : \eng{actually} ne signifie pas « actuellement » mais « en fait » ; « actuellement » se dit \eng{at present} ou \eng{currently}. Dans les phrases conditionnelles, ce piège revient souvent : \eng{If you actually come, I will be happy.} = \emph{Si tu viens vraiment, je serai content.}
\end{enumerate}
\end{attention}

\begin{savaistu}[Le 1st conditional au Bac]
Au baccalauréat, le 1st conditional est très fréquent dans les sujets de composition sur l'environnement, la santé ou l'éducation. Les correcteurs apprécient les phrases du type : \eng{If the government does nothing, pollution will destroy our cities.} ou \eng{If young people plant trees, Cameroon will be greener.} Maîtriser cette structure — sans jamais mettre \eng{will} après \eng{if} — est un moyen simple et efficace de gagner des points en expression écrite. C'est aussi la structure favorite des slogans écologiques anglophones : \eng{If you can't reuse it, refuse it!}
\end{savaistu}

\begin{exoresolu}[Exercice résolu — complétez au 1st conditional]
Énoncé : complétez les phrases avec la forme correcte des verbes entre parenthèses.

1. If we \_\_\_\_\_\_ (not / protect) the forest, many animals \_\_\_\_\_\_ (disappear).

2. If she \_\_\_\_\_\_ (recycle) her plastic bottles, she \_\_\_\_\_\_ (help) the planet.

3. I \_\_\_\_\_\_ (call) you if the doctor \_\_\_\_\_\_ (arrive).

4. Unless they \_\_\_\_\_\_ (stop) smoking, they \_\_\_\_\_\_ (have) health problems.

5. If you \_\_\_\_\_\_ (see) litter in the school yard, \_\_\_\_\_\_ (pick) it up.

\tcblower
Solution détaillée :

1. \eng{If we do not protect the forest, many animals will disappear.} — Après \eng{if} : présent simple négatif (\eng{do not protect}) ; proposition principale : \eng{will} + BV.

2. \eng{If she recycles her plastic bottles, she will help the planet.} — Attention à la 3\textsuperscript{e} personne du singulier au présent : \eng{recycles} avec un \eng{-s}.

3. \eng{I will call you if the doctor arrives.} — Phrase inversée : pas de virgule ; \eng{arrives} avec un \eng{-s} (3\textsuperscript{e} personne).

4. \eng{Unless they stop smoking, they will have health problems.} — Après \eng{unless}, forme affirmative : \eng{stop} (et non \eng{don't stop}).

5. \eng{If you see litter in the school yard, pick it up.} — Proposition principale à l'impératif : base verbale seule, sans sujet ni \eng{will}.
\end{exoresolu}

\begin{exercice}[À vous de jouer]
Mettez les verbes entre parenthèses à la forme correcte, puis traduisez les phrases 1 et 3 en français.

1. If the rains \_\_\_\_\_\_ (come) late, farmers in the North \_\_\_\_\_\_ (suffer).

2. You \_\_\_\_\_\_ (not / pass) the Bac unless you \_\_\_\_\_\_ (revise) your lessons.

3. If you \_\_\_\_\_\_ (feel) sick, \_\_\_\_\_\_ (go) to the health centre immediately.

4. If we \_\_\_\_\_\_ (burn) the bushes, the soil \_\_\_\_\_\_ (lose) its fertility.

5. If the government \_\_\_\_\_\_ (build) more hospitals, people in the villages \_\_\_\_\_\_ (can) receive better care.
\end{exercice}

\begin{corrige}[Exercice « À vous de jouer »]
1. \eng{If the rains come late, farmers in the North will suffer.}

2. \eng{You will not pass the Bac unless you revise your lessons.}

3. \eng{If you feel sick, go to the health centre immediately.}

4. \eng{If we burn the bushes, the soil will lose its fertility.}

5. \eng{If the government builds more hospitals, people in the villages can receive better care.}

Traductions : 1. \emph{Si les pluies arrivent tard, les agriculteurs du Nord souffriront.} 3. \emph{Si tu te sens malade, va immédiatement au centre de santé.}
\end{corrige}

\begin{aretenir}[L'essentiel du 1st conditional]
\begin{itemize}
\item Forme : \textbf{If + présent simple, will + BV} (ou \eng{can}, \eng{may}, \eng{might}, impératif).
\item Emploi : hypothèse \textbf{réelle ou probable} dans le futur (prédiction, avertissement, promesse, conseil).
\item Jamais de \eng{will} après \eng{if}.
\item Inversion possible : virgule si \eng{if} commence la phrase, pas de virgule sinon.
\item \eng{Unless} = \eng{if... not}, suivi d'une forme affirmative.
\end{itemize}
\end{aretenir}

\coursec{Le 2nd conditional : l'hypothèse irréelle}

Dans la section précédente, nous avons étudié le \cle{1st conditional} : \eng{If it rains, the match will be cancelled.} L'hypothèse était \textbf{réelle} — il peut vraiment pleuvoir demain à Douala. Mais que se passe-t-il quand l'hypothèse est \textbf{imaginaire}, impossible ou contraire à la réalité ? Que diriez-vous si vous étiez soudain ministre de la Santé ? C'est exactement le rôle du \cle{2nd conditional}, que nous allons maintenant découvrir.

\courssub{Observation : un texte pour découvrir}

\begin{textebox}[Rêver d'une ville plus saine — rédaction d'un élève]
If I were the mayor of Yaoundé, I would change many things. First, I would build more public toilets in every neighbourhood, because people would not throw waste in the streets if they had clean facilities. If our city had more rubbish bins, we could keep the markets and the roads clean. If the water supply were reliable, families would not buy expensive bottled water, and children would be healthier.

My friend Amina has different dreams. If she were rich, she would open a free clinic in her village. “If I were you,” she told me, “I would study medicine. You could help so many people!” I smiled and answered: “If I had more time, I might study both medicine and engineering!”

Of course, we are not the mayor, and we are not rich. But if every student planted one tree, our school would be greener. And if we all stopped burning plastic, the air would be cleaner. Small dreams can become real plans — if we work hard enough.
\end{textebox}

\textbf{Glossaire :}
\begin{itemize}
\item \eng{mayor} (n.) : maire ; \eng{waste} (n.) : déchets ; \eng{rubbish bin} (n.) : poubelle
\item \eng{water supply} (n.) : approvisionnement en eau ; \eng{reliable} (adj.) : fiable
\item \eng{to plant a tree} (v.) : planter un arbre ; \eng{to burn} (v.) : brûler
\item \eng{neighbourhood} (n.) : quartier ; \eng{clinic} (n.) : dispensaire, clinique
\end{itemize}

\textbf{Questions de compréhension :}
\begin{enumerate}
\item What would the writer do if he were the mayor of Yaoundé?
\item What would Amina do if she were rich?
\item What advice does Amina give to her friend?
\item Is the writer really the mayor? How do you know?
\end{enumerate}

\textbf{Observation grammaticale.} Relevons quelques phrases du texte :

\begin{exemplebox}[Phrases à observer]
\begin{itemize}
\item \eng{If I were the mayor of Yaoundé, I would change many things.} « Si j'étais le maire de Yaoundé, je changerais beaucoup de choses. »
\item \eng{If our city had more rubbish bins, we could keep the markets clean.} « Si notre ville avait plus de poubelles, nous pourrions garder les marchés propres. »
\item \eng{If I were you, I would study medicine.} « Si j'étais toi, j'étudierais la médecine. »
\end{itemize}
\end{exemplebox}

Remarquez deux choses essentielles : (1) le verbe après \eng{if} est au \textbf{prétérit} (\eng{were}, \eng{had}), pourtant la phrase parle du \textbf{présent}, pas du passé ; (2) l'autre proposition contient \eng{would} + base verbale. C'est la signature du 2nd conditional.

\courssub{La règle : forme et emploi}

\begin{propriete}[Règle du 2nd conditional]
\textbf{Forme :} \cle{If + sujet + prétérit simple, sujet + would + base verbale (BV)}

\eng{If I had money, I would travel.} « Si j'avais de l'argent, je voyagerais. »

\textbf{Point capital :} le prétérit après \eng{if} n'exprime \textbf{pas le passé} : il exprime l'\textbf{irréel du présent} — une situation imaginaire, impossible ou contraire à la réalité actuelle.

\textbf{Emplois :}
\begin{enumerate}
\item \textbf{Situation imaginaire ou contraire à la réalité présente} : \eng{If we had clean water, we would be healthier.} (mais nous n'avons pas d'eau potable)
\item \textbf{Rêve, souhait} : \eng{If I were a pilot, I would fly over Mount Cameroon.}
\item \textbf{Conseil poli} : \eng{If I were you, I would see a doctor.}
\end{enumerate}

\textbf{Inversion possible :} comme au 1st conditional, on peut inverser les deux propositions sans virgule : \eng{I would travel if I had money.}
\end{propriete}

\begin{center}
\begin{tabular}{|l|l|l|}
\hline
\rowcolor{popblueL} \textbf{Proposition en if (condition)} & \textbf{Proposition principale (résultat)} & \textbf{Traduction} \\
\hline
If I were rich, & I would build a hospital. & Si j'étais riche, je construirais un hôpital. \\
\hline
If she lived in Buea, & she would see the mountain every day. & Si elle vivait à Buea, elle verrait la montagne chaque jour. \\
\hline
If they had more bins, & they would keep the street clean. & S'ils avaient plus de poubelles, ils garderaient la rue propre. \\
\hline
If we didn't smoke, & we would breathe better. & Si nous ne fumions pas, nous respirerions mieux. \\
\hline
\end{tabular}
\end{center}

\textbf{Forme négative et interrogative :}
\begin{itemize}
\item Négation dans la condition : \eng{If I didn't have a headache, I would come.} « Si je n'avais pas mal à la tête, je viendrais. »
\item Négation dans le résultat : \eng{If it rained, the crops would not die.} « S'il pleuvait, les récoltes ne mourraient pas. »
\item Question : \eng{What would you do if you won the lottery?} « Que ferais-tu si tu gagnais à la loterie ? » ; \eng{If you were sick, would you go to the hospital?}
\end{itemize}

\courssub{Le cas spécial : If I were you}

\begin{definition}[La formule « If I were you »]
\eng{If I were you...} est une \textbf{formule figée} qui sert à donner un \textbf{conseil poli}. Elle signifie « Si j'étais toi / vous... », c'est-à-dire « À ta place... ». Elle est toujours suivie de \eng{would} + BV.

\eng{If I were you, I would see a doctor.} « Si j'étais toi, je verrais un médecin. » (= Tu devrais voir un médecin.)
\end{definition}

\begin{propriete}[Were pour toutes les personnes]
En style soigné et à l'écrit, on utilise \cle{were} (et non \eng{was}) pour \textbf{toutes les personnes} dans le 2nd conditional, même à la 1\textsuperscript{re} et à la 3\textsuperscript{e} personne du singulier :

\eng{If I were taller, I would play basketball.} « Si j'étais plus grand, je jouerais au basket. »

\eng{If she were here, she would help us.} « Si elle était là, elle nous aiderait. »

À l'oral familier, \eng{was} (\eng{If I was you...}) se rencontre, mais dans un examen et dans une rédaction, utilisez toujours \eng{were}.
\end{propriete}

\begin{exemplebox}[Conseils polis avec If I were you]
\begin{itemize}
\item \eng{If I were you, I would drink more water.} « Si j'étais toi, je boirais plus d'eau. »
\item \eng{If I were you, I would not eat that street food.} « À ta place, je ne mangerais pas cette nourriture de rue. »
\item \eng{If I were you, I would sleep earlier before the exam.} « Si j'étais toi, je me coucherais plus tôt avant l'examen. »
\end{itemize}
\end{exemplebox}

\courssub{Variantes avec could et might}

À la place de \eng{would}, on peut utiliser \cle{could} (capacité) ou \cle{might} (possibilité incertaine) dans la proposition principale :

\begin{center}
\begin{tabularx}{\linewidth}{|X|X|X|}
\hline
\rowcolor{popblueL} \textbf{Auxiliaire} & \textbf{Sens} & \textbf{Exemple et traduction} \\
\hline
would & certitude dans l'hypothèse & \eng{If I had money, I would buy a bicycle.} « Si j'avais de l'argent, j'achèterais un vélo. » \\
\hline
could & capacité (= would be able to) & \eng{If we had more bins, we could keep the city clean.} « Si nous avions plus de poubelles, nous pourrions garder la ville propre. » \\
\hline
might & possibilité douteuse (= perhaps) & \eng{If you exercised, you might feel better.} « Si tu faisais du sport, tu te sentirais peut-être mieux. » \\
\hline
\end{tabularx}
\end{center}

\courssub{Probabilité ou imagination ? 1st vs 2nd conditional}

C'est le choix du locuteur qui décide : le 1st conditional présente l'hypothèse comme \textbf{possible, probable} ; le 2nd conditional la présente comme \textbf{imaginaire, improbable ou impossible}.

\begin{exemplebox}[Comparaison de sens]
\begin{itemize}
\item \eng{If it rains tomorrow, I will stay at home.} « S'il pleut demain, je resterai à la maison. » (il peut vraiment pleuvoir — saison des pluies à Bamenda !)
\item \eng{If it rained in the Sahara, the desert would become green.} « S'il pleuvait au Sahara, le désert deviendrait vert. » (imagination pure)
\item \eng{If I pass my exam, I will celebrate.} « Si je réussis mon examen, je fêterai cela. » (réaliste)
\item \eng{If I were the President, I would build more hospitals.} « Si j'étais le Président, je construirais plus d'hôpitaux. » (je ne suis pas le Président !)
\end{itemize}
\end{exemplebox}

\begin{popfigure}
\begin{tikzpicture}[
  box/.style={draw=popblue, thick, rounded corners=4pt, fill=popblueL, align=center, text width=5.6cm, inner sep=6pt},
  lab/.style={font=\bfseries, text=popdark, align=center}
]
\node[lab] (t1) at (0,2.6) {1st conditional — réel, probable};
\node[box, align=center] (b1) at (0,1.2){\eng{If it rains, I will stay.}\\ If + présent, will + BV\\ « S'il pleut, je resterai. »};
\node[lab] (t2) at (0,-0.4) {2nd conditional — irréel, imaginaire};
\node[box, draw=poporange, fill=poporangeL, align=center] (b2) at (0,-2.0){\eng{If it rained, I would stay.}\\ If + prétérit, would + BV\\ « S'il pleuvait, je resterais. »};
\draw[-{Stealth}, very thick, poppurple] (b1) -- (t2) node[midway, right, text=poppurple, align=left, font=\small]{présent $\rightarrow$ prétérit\\ will $\rightarrow$ would\\ (recul d'un temps = irréel)};
\end{tikzpicture}
\legende{Du 1st au 2nd conditional : recul d'un temps pour marquer l'irréel.}
\end{popfigure}

\courssub{Comparaison avec le français}

Bonne nouvelle : le français fonctionne presque pareil ! Le français utilise l'\textbf{imparfait} après « si » et le \textbf{conditionnel} dans l'autre proposition ; l'anglais utilise le \textbf{prétérit} et \eng{would}.

\begin{center}
\begin{tabularx}{\linewidth}{|X|X|X|}
\hline
\rowcolor{popblueL} \textbf{Français} & \textbf{English} & \textbf{Remarque} \\
\hline
Si j'étais riche, je voyagerais. & \eng{If I were rich, I would travel.} & Imparfait = prétérit ; conditionnel = would + BV. \\
\hline
Si nous avions de l'eau potable, nous serions en meilleure santé. & \eng{If we had clean water, we would be healthier.} & Correspondance directe. \\
\hline
Si j'étais toi, je verrais un médecin. & \eng{If I were you, I would see a doctor.} & Formule de conseil identique dans les deux langues. \\
\hline
\end{tabularx}
\end{center}

\begin{attention}[Piège n° 1 : jamais de would dans la proposition en if]
L'erreur la plus fréquente des francophones est de calquer le conditionnel français sur les deux moitiés de la phrase. En anglais, \eng{would} est \textbf{interdit} après \eng{if} :

\begin{itemize}
\item Faux : \eng{If I would be rich, I would travel.}
\item Correct : \eng{If I were rich, I would travel.}
\item Faux : \eng{If we would have more bins, we could keep the city clean.}
\item Correct : \eng{If we had more bins, we could keep the city clean.}
\end{itemize}

Retenez : \textbf{un seul would par phrase} — dans la proposition principale, jamais dans la condition.
\end{attention}

\begin{attention}[Piège n° 2 : le prétérit n'est pas un passé]
Ne traduisez pas \eng{If I had money} par « Si j'avais eu de l'argent ». Le prétérit du 2nd conditional parle du \textbf{présent imaginaire} : « Si j'avais de l'argent (maintenant) ». Pour parler du passé, il faudra le 3rd conditional (section suivante).
\end{attention}

\courssub{Entraînement guidé}

\begin{exoresolu}[Complétez avec le 2nd conditional]
Mettez les verbes entre parenthèses à la forme correcte.

1. If I (be) the minister of Health, I (build) more hospitals in rural areas.

2. If we (have) clean water in our quarter, we (be) healthier.

3. If she (not smoke), she (breathe) better.

4. What you (do) if you (win) one million francs?

5. If they (plant) more trees, the air in Douala (be) cleaner.
\tcblower
\textbf{Solution détaillée :}

1. \eng{If I \textbf{were} the minister of Health, I \textbf{would build} more hospitals in rural areas.} — Après \eng{if}, prétérit de \eng{be} : \eng{were} pour toutes les personnes en style soigné. Résultat : would + BV.

2. \eng{If we \textbf{had} clean water in our quarter, we \textbf{would be} healthier.} — Prétérit de \eng{have} : \eng{had}. Attention : \eng{would be}, car \eng{be} reste à la base verbale après \eng{would}.

3. \eng{If she \textbf{didn't smoke}, she \textbf{would breathe} better.} — Négation au prétérit : \eng{didn't} + BV (\eng{smoke}, sans -s ni -ed).

4. \eng{What \textbf{would} you \textbf{do} if you \textbf{won} one million francs?} — Question : \eng{would} se place avant le sujet dans la proposition principale ; la condition reste au prétérit (\eng{won}, prétérit irrégulier de \eng{win}).

5. \eng{If they \textbf{planted} more trees, the air in Douala \textbf{would be} cleaner.} — Prétérit régulier : \eng{plant} $\rightarrow$ \eng{planted}.
\end{exoresolu}

\begin{exercice}[À vous de jouer]
\begin{enumerate}
\item Traduisez en anglais : « Si j'étais toi, je consulterais un médecin. »
\item Traduisez : « Si nous avions plus de poubelles, nous pourrions garder la ville propre. »
\item Corrigez la faute : \eng{If I would have a car, I would drive to Kribi.}
\item Complétez : If Cameroon (win) the World Cup, everybody (celebrate) in the streets.
\item Écrivez trois phrases au 2nd conditional sur votre école : que changeriez-vous si vous étiez le proviseur ?
\end{enumerate}
\end{exercice}

\begin{corrige}[Exercice « À vous de jouer »]
\begin{enumerate}
\item \eng{If I were you, I would see a doctor.} — Formule figée de conseil : \eng{were}, jamais \eng{would be} après \eng{if}.
\item \eng{If we had more bins, we could keep the city clean.} — « Pourrions » = capacité imaginaire : \eng{could}.
\item \eng{If I \textbf{had} a car, I would drive to Kribi.} — Jamais de \eng{would} dans la proposition en \eng{if} ; prétérit de \eng{have} : \eng{had}.
\item \eng{If Cameroon \textbf{won} the World Cup, everybody \textbf{would celebrate} in the streets.} — Hypothèse imaginaire (hélas !) : prétérit + would.
\item Réponses libres. Modèle : \eng{If I were the principal, I would build a bigger library. If we had more computers, we could learn faster. If the school were cleaner, we would be healthier.}
\end{enumerate}
\end{corrige}

\begin{aretenir}[L'essentiel du 2nd conditional]
\begin{itemize}
\item \textbf{Forme :} If + prétérit simple, sujet + would + BV.
\item \textbf{Sens :} hypothèse imaginaire, impossible ou contraire à la réalité présente ; le prétérit marque l'irréel, pas le passé.
\item \textbf{Conseil poli :} \eng{If I were you, I would...} avec \eng{were} pour toutes les personnes.
\item \textbf{Variantes :} \eng{could} (capacité), \eng{might} (possibilité incertaine).
\item \textbf{Interdit :} \eng{would} dans la proposition en \eng{if}.
\item \textbf{Choix 1st/2nd :} probable $\rightarrow$ 1st conditional ; imaginaire $\rightarrow$ 2nd conditional.
\end{itemize}
\end{aretenir}

Et si l'hypothèse porte non pas sur le présent, mais sur le \textbf{passé} — un regret, un reproche : « Si seulement j'avais étudié davantage... » ? C'est le domaine du \cle{3rd conditional}, que nous étudierons dans la section suivante.

\coursec{Le 3rd conditional : le regret sur le passé}

Dans la section précédente, nous avons étudié le \cle{2nd conditional}, qui exprime une hypothèse \emph{irréelle au présent} : \eng{If I had money, I would buy a solar panel.} (« Si j'avais de l'argent — mais je n'en ai pas —, j'achèterais un panneau solaire. »). Mais que se passe-t-il quand l'irréel porte sur le \emph{passé}, sur quelque chose qui est terminé et qu'on ne peut plus changer ? C'est exactement le rôle du \cle{3rd conditional} : il exprime un passé qui aurait pu être différent — un regret, un reproche, un « et si... ».

\courssub{Observation : un texte pour découvrir la structure}

Lisons d'abord ce court article original sur les inondations à Douala. Soulignez mentalement toutes les phrases qui commencent par \eng{If}.

\begin{textebox}[The floods that could have been avoided — Douala, last rainy season]
Last October, heavy rains fell on Douala for three days without stopping. Several neighbourhoods near the Wouri River were flooded, and hundreds of families lost their belongings. After the disaster, many people asked the same question: could this have been avoided?

“If the city council had cleaned the gutters before the rainy season, the water would have flowed away,” explained a local engineer. “If residents had not thrown plastic bags into the drains, the pipes would not have been blocked.”

A shopkeeper in Akwa added: “If I had moved my goods to the first floor, I would not have lost my whole stock. If someone had warned us earlier, we could have saved many things.”

The mayor admitted that mistakes had been made. “If we had listened to the weather forecasts, we might have organised an evacuation,” he said. “If the government had invested in better drainage years ago, Douala would be a safer city today.”

Nobody can change the past now. But if the authorities learn from this disaster, the next rainy season will not be so dramatic.
\end{textebox}

\textbf{Glossaire}

\begin{center}
\begin{tabularx}{\linewidth}{|l|l|X|}
\hline
\rowcolor{popblueL} \textbf{Mot anglais} & \textbf{Nature} & \textbf{Traduction} \\
\hline
to flood & verb & inonder \\
\hline
belongings & noun & affaires (personnelles) \\
\hline
gutter & noun & caniveau, gouttière \\
\hline
to flow away & phrasal verb & s'écouler \\
\hline
drain & noun & égout, conduit d'évacuation \\
\hline
goods & noun (plur.) & marchandises \\
\hline
stock & noun & stock, réserve \\
\hline
forecast & noun & prévision (météo) \\
\hline
drainage & noun & système d'évacuation des eaux \\
\hline
\end{tabularx}
\end{center}

\textbf{Questions de compréhension}

\begin{enumerate}
\item What happened in Douala last October?
\item According to the engineer, what should the city council have done?
\item What does the shopkeeper regret?
\item Find in the text one sentence with \eng{could have} and one with \eng{might have}.
\item Which sentence in the last paragraph uses a \emph{1st conditional}? Why is it different from the others?
\end{enumerate}

\courssub{Observation guidée : que remarquez-vous ?}

Reprenons deux phrases du texte :

\eng{If the city council had cleaned the gutters, the water would have flowed away.} « Si la municipalité avait nettoyé les caniveaux, l'eau se serait écoulée. »

\eng{If I had moved my goods, I would not have lost my whole stock.} « Si j'avais déplacé mes marchandises, je n'aurais pas perdu tout mon stock. »

Observons la structure :

\begin{center}
\begin{tabularx}{\linewidth}{|X|X|}
\hline
\rowcolor{popblueL} \textbf{Proposition avec IF (condition)} & \textbf{Proposition principale (résultat)} \\
\hline
If the city council \textbf{had cleaned}... & the water \textbf{would have flowed} away. \\
\hline
If I \textbf{had moved} my goods... & I \textbf{would not have lost} my stock. \\
\hline
If we \textbf{had listened}... & we \textbf{might have organised} an evacuation. \\
\hline
\end{tabularx}
\end{center}

On constate deux choses essentielles :
\begin{itemize}
\item après \eng{if}, on trouve \eng{had} + participe passé : c'est le \cle{past perfect} (plus-que-parfait anglais) ;
\item dans l'autre proposition, on trouve \eng{would have} + participe passé (ou \eng{might have}, \eng{could have}).
\end{itemize}

Et surtout : est-ce que la municipalité a nettoyé les caniveaux ? \textbf{Non.} Est-ce que le commerçant a déplacé ses marchandises ? \textbf{Non.} La condition est \emph{contraire à la réalité passée} : on imagine un passé différent de celui qui s'est réellement produit.

\begin{attention}[Un cas particulier dans le texte : le conditionnel mixte]
Remarquez la phrase du maire : \eng{“If the government had invested in better drainage years ago, Douala would be a safer city today.”} La condition est passée (\eng{had invested}), mais le résultat (\eng{would be}) est au \emph{présent}. C'est un \textbf{conditionnel mixte} (passé $\rightarrow$ présent) que nous étudierons en fin de leçon.
\end{attention}

\courssub{La règle : forme et emploi}

\begin{propriete}[Le 3rd conditional : forme]
Le \cle{3rd conditional} se construit ainsi :

\textbf{If + sujet + past perfect (had + participe passé), sujet + would have + participe passé.}

\begin{itemize}
\item \textbf{Proposition de condition} : \eng{If} + \eng{had} + participe passé (pour toutes les personnes, sans exception : \eng{If I had}, \eng{If she had}, \eng{If they had}).
\item \textbf{Proposition de résultat} : sujet + \eng{would have} + participe passé.
\item \textbf{Négation} : \eng{If I had not (hadn't) known...} / \eng{...I would not have come.}
\item \textbf{Inversion possible} : les deux propositions peuvent changer de place, sans virgule quand \eng{if} est en seconde position : \eng{I would have come if I had known.}
\end{itemize}
\end{propriete}

\begin{definition}[Emploi du 3rd conditional]
Le 3rd conditional exprime une \cle{hypothèse contraire à la réalité passée} : la condition ne s'est \emph{pas} réalisée, et le résultat ne s'est \emph{pas} produit non plus. Il sert à exprimer :
\begin{itemize}
\item le \textbf{regret} : \eng{If I had worked harder, I would have passed the Bac.} (« Si j'avais travaillé davantage, j'aurais réussi le Bac. » — mais je n'ai pas travaillé, et j'ai échoué.)
\item le \textbf{reproche} : \eng{If you had told me, I would have helped you.} (« Si tu me l'avais dit, je t'aurais aidé. » — tu ne me l'as pas dit.)
\item le \textbf{soulagement} ou l'\textbf{imagination} d'un passé différent : \eng{If the driver had been drunk, the accident would have been worse.} (« Si le chauffeur avait été ivre, l'accident aurait été pire. » — heureusement, il ne l'était pas.)
\end{itemize}
\end{definition}

\begin{exemplebox}[Exemples traduits]
\eng{If they had planted trees, the soil would not have eroded.} « S'ils avaient planté des arbres, le sol ne se serait pas érodé. »

\eng{If I had known, I would have come.} « Si j'avais su, je serais venu. »

\eng{If she had taken the bendskin earlier, she would not have missed the exam.} « Si elle avait pris le bendskin plus tôt, elle n'aurait pas raté l'examen. »

\eng{If we had protected the forest, the climate would have stayed cooler.} « Si nous avions protégé la forêt, le climat serait resté plus frais. »

\eng{If he had not wasted water, the tank would not have been empty.} « S'il n'avait pas gaspillé l'eau, le réservoir n'aurait pas été vide. »
\end{exemplebox}

\textbf{Tableau récapitulatif des formes}

\begin{center}
\begin{tabular}{|l|l|l|}
\hline
\rowcolor{popblueL} \textbf{Forme} & \textbf{Proposition IF} & \textbf{Proposition de résultat} \\
\hline
Affirmative & If I had studied... & I would have passed. \\
\hline
Négative (if) & If I had not studied... & I would have failed. \\
\hline
Négative (résultat) & If I had studied... & I would not have failed. \\
\hline
Interrogative & If you had called... & would you have come? \\
\hline
\end{tabular}
\end{center}

\courssub{Variantes avec could have et might have}

Dans la proposition de \emph{résultat}, on peut remplacer \eng{would have} par \eng{could have} ou \eng{might have} pour nuancer le sens :

\begin{center}
\begin{tabularx}{\linewidth}{|l|X|X|}
\hline
\rowcolor{popblueL} \textbf{Forme} & \textbf{Sens} & \textbf{Exemple} \\
\hline
would have + p.p. & certitude (dans l'hypothèse) & If you had left earlier, you \textbf{would have caught} the bus. (certain) \\
\hline
could have + p.p. & possibilité, capacité & If you had left earlier, you \textbf{could have caught} the bus. (tu aurais \emph{pu} l'attraper) \\
\hline
might have + p.p. & possibilité incertaine & If you had left earlier, you \textbf{might have caught} the bus. (peut-être) \\
\hline
\end{tabularx}
\end{center}

\eng{If you had left earlier, you could have caught the bus.} « Si tu étais parti plus tôt, tu aurais pu attraper le bus. »

\eng{If we had listened to the forecasts, we might have organised an evacuation.} « Si nous avions écouté les prévisions, nous aurions peut-être organisé une évacuation. »

\begin{attention}[could have dans la proposition IF ?]
\eng{Could have} et \eng{might have} se placent dans la proposition de \textbf{résultat}, jamais après \eng{if}. Après \eng{if}, on garde toujours le past perfect : \eng{If you \textbf{had left} earlier...} et non \emph{If you could have left earlier...} pour exprimer la condition.
\end{attention}

\courssub{Les conditionnels mixtes (complément pour le Bac)}

Parfois, la condition porte sur le \emph{passé}, mais le résultat se situe dans le \emph{présent}. On combine alors un 3rd conditional (pour la condition) et un 2nd conditional (pour le résultat) :

\begin{propriete}[Conditionnel mixte passé $\rightarrow$ présent]
\textbf{If + past perfect, sujet + would + base verbale.}

\eng{If I had studied medicine, I would be a doctor now.} « Si j'avais fait médecine, je serais médecin aujourd'hui. »

La condition est passée et irréelle (je n'ai pas fait médecine), mais la conséquence imaginaire concerne ma situation \emph{actuelle}.
\end{propriete}

Autres exemples :

\eng{If we had planted more trees ten years ago, our town would be greener today.} « Si nous avions planté plus d'arbres il y a dix ans, notre ville serait plus verte aujourd'hui. »

\eng{If she had learned English at school, she would have a better job now.} « Si elle avait appris l'anglais à l'école, elle aurait un meilleur emploi maintenant. »

\courssub{Vue d'ensemble : la frise des trois conditionnels}

\begin{popfigure}
\begin{center}
\begin{tikzpicture}[scale=1.0]
% axe du temps
\draw[-{Stealth[length=3mm]}, thick, popdark] (-5.5,0) -- (6,0);
\node[below, popdark] at (6.1,0) {\small temps};
% graduations
\draw[thick, popdark] (-3.5,-0.15) -- (-3.5,0.15);
\node[below, popdark, align=center] at (-3.5,-0.2) {\small PASSÉ};
\draw[thick, popdark] (0,-0.15) -- (0,0.15);
\node[below, popdark, align=center] at (0,-0.2) {\small PRÉSENT};
\draw[thick, popdark] (3.5,-0.15) -- (3.5,0.15);
\node[below, popdark, align=center] at (3.5,-0.2) {\small FUTUR};
% 3rd conditional
\node[draw, rounded corners, fill=poppinkL, align=center, text width=3.6cm] at (-3.5,1.6) {\small \textbf{3rd conditional}\\ If + past perfect,\\ would have + p.p.};
\node[align=center, poppink!60!black] at (-3.5,-1.2) {\small regret : \emph{If I had known,}\\ \emph{I would have come.}};
\draw[-{Stealth}, poppink!60!black] (-3.5,0.2) -- (-3.5,1.0);
% 2nd conditional
\node[draw, rounded corners, fill=poppurpleL, align=center, text width=3.6cm] at (0,1.6) {\small \textbf{2nd conditional}\\ If + past simple,\\ would + base verbale};
\node[align=center, poppurple] at (0,-1.2) {\small irréel présent : \emph{If I had}\\ \emph{money, I would buy it.}};
\draw[-{Stealth}, poppurple] (0,0.2) -- (0,1.0);
% 1st conditional
\node[draw, rounded corners, fill=popgreenL, align=center, text width=3.6cm] at (3.5,1.6) {\small \textbf{1st conditional}\\ If + present simple,\\ will + base verbale};
\node[align=center, popgreen!60!black] at (3.5,-1.2) {\small possible : \emph{If it rains,}\\ \emph{I will stay home.}};
\draw[-{Stealth}, popgreen!60!black] (3.5,0.2) -- (3.5,1.0);
% conditionnel mixte
\draw[-{Stealth}, dashed, thick, poporange] (-2.2,2.3) .. controls (-1,3.1) .. (0.6,2.3);
\node[align=center, poporange] at (-0.8,3.3) {\small mixte : passé $\rightarrow$ présent};
\end{tikzpicture}
\end{center}
\legende{Frise des temps : le 3rd conditional regarde vers un passé irréel, le 2nd vers un présent imaginaire, le 1st vers un futur possible. La flèche en pointillés montre le conditionnel mixte.}
\end{popfigure}

\courssub{Comparaison avec le français}

Bonne nouvelle : le français possède une structure presque identique, avec le plus-que-parfait après « si » et le conditionnel passé dans l'autre proposition. C'est l'un des rares cas où la traduction est presque mot à mot !

\begin{center}
\begin{tabularx}{\linewidth}{|X|X|X|}
\hline
\rowcolor{popblueL} \textbf{Français} & \textbf{English} & \textbf{Remarque} \\
\hline
Si j'avais su, je serais venu. & If I had known, I would have come. & plus-que-parfait = past perfect ; conditionnel passé = would have + p.p. \\
\hline
Si tu étais parti plus tôt, tu aurais pu attraper le bus. & If you had left earlier, you could have caught the bus. & « aurais pu » = could have + p.p. \\
\hline
S'ils avaient planté des arbres, le sol ne se serait pas érodé. & If they had planted trees, the soil would not have eroded. & attention à la négation : would \emph{not} have eroded \\
\hline
Si j'avais fait médecine, je serais médecin aujourd'hui. & If I had studied medicine, I would be a doctor now. & conditionnel mixte : résultat au présent \\
\hline
\end{tabularx}
\end{center}

\begin{attention}[L'erreur n° 1 des francophones : would have après IF]
Ne dites \textbf{jamais} : \emph{If I would have known, I would have come.} En anglais, \eng{would have} ne se place \textbf{jamais} dans la proposition introduite par \eng{if}. Après \eng{if}, on utilise uniquement le past perfect :

\eng{If I \textbf{had known}, I would have come.} (correct)

\emph{If I would have known...} (incorrect — faute très grave à l'examen)

De même : \eng{If it \textbf{had rained}...} et non \emph{If it would have rained...}. Retenez le slogan : \textbf{« No would after if ! »}
\end{attention}

\begin{attention}[Autres pièges fréquents]
\begin{itemize}
\item \textbf{Oublier le participe passé} : \emph{I would have go} est faux ; dites \eng{I would have \textbf{gone}}. Révisez vos verbes irréguliers : go–went–gone, know–knew–known, come–came–come, take–took–taken, see–saw–seen.
\item \textbf{Confondre les conditionnels} : pour un événement passé, n'utilisez pas le 2nd conditional. \emph{If I knew, I would come} parle du présent ; pour le passé : \eng{If I had known, I would have come.}
\item \textbf{Traduire « si » par deux temps identiques} : en français on dit « si j'aurais su » dans la langue familière — c'est une faute en français aussi, et elle devient catastrophique en anglais.
\end{itemize}
\end{attention}

\courssub{Prononciation des formes contractées}

À l'oral, les anglophones contractent presque toujours ces formes. Il faut les reconnaître à la compréhension orale et oser les utiliser à l'expression orale :

\begin{center}
\begin{tabular}{|l|l|l|}
\hline
\rowcolor{popblueL} \textbf{Forme complète} & \textbf{Forme contractée} & \textbf{Prononciation (approximation française)} \\
\hline
I would have & I'd have & « aïd av » \\
\hline
you would have & you'd have & « youd av » \\
\hline
she would have & she'd have & « chid av » \\
\hline
we would have & we'd have & « ouid av » \\
\hline
had not & hadn't & « a-dent » \\
\hline
would not have & wouldn't have & « wou-dent av » \\
\hline
\end{tabular}
\end{center}

Remarquez que \eng{have} se prononce souvent de façon très réduite, presque « av » ou « euhv » : \eng{I would have come} se prononce « aï wou-deu-v kom ». Attention : à l'\textbf{écrit formel} (composition du Bac), évitez les contractions ; écrivez les formes complètes.

\begin{savaistu}[Le 3rd conditional, l'arme du reproche]
Dans les conversations quotidiennes en anglais, le 3rd conditional est l'outil préféré pour exprimer le \cle{reproche} : \eng{If you had listened to me, none of this would have happened!} (« Si tu m'avais écouté, rien de tout cela ne serait arrivé ! »). Les parents, les entraîneurs de football et les professeurs anglophones l'emploient constamment : \eng{If you had trained harder, you would have won the match.} C'est aussi le temps des chansons et des films qui imaginent un passé différent — le fameux « what if...? » (« et si... ? ») des histoires. Au Cameroun, après une défaite des Lions Indomptables, les débats à la radio anglophone sont pleins de 3rd conditionals : \eng{If the striker had scored that penalty, we would have qualified!}
\end{savaistu}

\courssub{Exercice résolu et entraînement}

\begin{exoresolu}[Compléter avec le 3rd conditional]
Énoncé — Complétez les phrases en mettant les verbes entre parenthèses à la forme correcte.

1. If the farmers \ldots\ldots\ldots\ldots (plant) trees, the soil \ldots\ldots\ldots\ldots (not / erode).
2. If I \ldots\ldots\ldots\ldots (know) about the meeting, I \ldots\ldots\ldots\ldots (come).
3. If you \ldots\ldots\ldots\ldots (leave) earlier, you \ldots\ldots\ldots\ldots (can / catch) the bus.
4. If she \ldots\ldots\ldots\ldots (study) medicine, she \ldots\ldots\ldots\ldots (be) a doctor now.

\tcblower
Solution détaillée :

1. \eng{If the farmers \textbf{had planted} trees, the soil \textbf{would not have eroded}.} Condition passée irréelle : past perfect après \eng{if}, \eng{would have} + participe passé dans le résultat, avec négation (\eng{would not have} + p.p.).

2. \eng{If I \textbf{had known} about the meeting, I \textbf{would have come}.} Verbe irrégulier : know–knew–\textbf{known} ; come–came–\textbf{come} (participe passé identique à la base).

3. \eng{If you \textbf{had left} earlier, you \textbf{could have caught} the bus.} Ici le sens est « tu aurais pu attraper » : on utilise \eng{could have} + participe passé (catch–caught–\textbf{caught}).

4. \eng{If she \textbf{had studied} medicine, she \textbf{would be} a doctor now.} Le mot \eng{now} indique un résultat au présent : c'est un \textbf{conditionnel mixte} — past perfect après \eng{if}, mais \eng{would} + base verbale dans le résultat.
\end{exoresolu}

\begin{exercice}[À vous de jouer]
1. Mettez les verbes entre parenthèses à la forme correcte :

a) If we \ldots\ldots\ldots\ldots (protect) the forest, the rain \ldots\ldots\ldots\ldots (not / disappear).

b) If they \ldots\ldots\ldots\ldots (listen) to the doctor, they \ldots\ldots\ldots\ldots (not / fall) ill.

c) If he \ldots\ldots\ldots\ldots (save) his money, he \ldots\ldots\ldots\ldots (can / buy) a bicycle.

2. Traduisez en anglais :

a) Si j'avais pris mon parapluie, je n'aurais pas été mouillé.

b) Si tu m'avais écouté, tu aurais réussi l'examen.

c) Si elle avait appris l'anglais, elle aurait un meilleur emploi aujourd'hui.

3. Transformez en phrase au 3rd conditional : \eng{I did not water the plants, so they died.} (If...)
\end{exercice}

\begin{corrige}[Exercice « À vous de jouer »]
1. a) \eng{If we \textbf{had protected} the forest, the rain \textbf{would not have disappeared}.} b) \eng{If they \textbf{had listened} to the doctor, they \textbf{would not have fallen} ill.} (fall–fell–\textbf{fallen}) c) \eng{If he \textbf{had saved} his money, he \textbf{could have bought} a bicycle.} (buy–bought–\textbf{bought})

2. a) \eng{If I had taken my umbrella, I would not have got wet.} b) \eng{If you had listened to me, you would have passed the exam.} c) \eng{If she had learned English, she would have a better job today.} (conditionnel mixte : résultat au présent)

3. \eng{If I had watered the plants, they would not have died.}
\end{corrige}

\begin{aretenir}[L'essentiel du 3rd conditional]
\begin{itemize}
\item \textbf{Forme} : \eng{If} + sujet + \eng{had} + participe passé, sujet + \eng{would have} + participe passé.
\item \textbf{Emploi} : hypothèse contraire à la réalité passée — regret, reproche, « et si... ».
\item \textbf{Variantes} : \eng{could have} (possibilité) et \eng{might have} (incertitude) remplacent \eng{would have} dans le résultat.
\item \textbf{Mixte} : passé $\rightarrow$ présent : \eng{If I had studied medicine, I would be a doctor now.}
\item \textbf{Piège mortel} : jamais de \eng{would have} après \eng{if} — « No would after if ! »
\item \textbf{Oral} : formes contractées \eng{I'd have}, \eng{hadn't}, \eng{wouldn't have} ; à l'écrit formel, formes complètes.
\end{itemize}
\end{aretenir}

\coursec{Comparaison français/anglais et pièges}

Nous venons d'étudier le \cle{3rd conditional}, cette structure du regret qui permet de réécrire le passé : \eng{If I had planted trees, the soil would not have eroded.} Vous maîtrisez maintenant les trois conditionnels. Il reste une étape décisive : comprendre comment ces structures correspondent au français — et surtout où elles en diffèrent. C'est précisément là que les élèves francophones perdent des points à l'examen. Cette section vous donne les clés pour ne plus jamais tomber dans les pièges.

\courssub{La correspondance globale des temps}

Observons d'abord trois phrases jumelles, l'une en français, l'autre en anglais :

\begin{exemplebox}[Les trois correspondances fondamentales]
\begin{enumerate}
\item « Si tu \textbf{viens}, je \textbf{serai} content. » = \eng{If you \textbf{come}, I \textbf{will be} happy.}
\item « Si tu \textbf{venais}, je \textbf{serais} content. » = \eng{If you \textbf{came}, I \textbf{would be} happy.}
\item « Si tu \textbf{étais venu}, j'\textbf{aurais été} content. » = \eng{If you \textbf{had come}, I \textbf{would have been} happy.}
\end{enumerate}
\end{exemplebox}

La bonne nouvelle, c'est que la logique est \textbf{identique} dans les deux langues. Le français et l'anglais utilisent les mêmes temps dans la subordonnée introduite par \eng{if} / « si » :

\begin{propriete}[Correspondance des temps français/anglais]
\begin{itemize}
\item Français « si + \cle{présent} » = anglais \eng{if + present simple} → 1st conditional (hypothèse réelle).
\item Français « si + \cle{imparfait} » = anglais \eng{if + past simple} → 2nd conditional (hypothèse irréelle au présent).
\item Français « si + \cle{plus-que-parfait} » = anglais \eng{if + past perfect} → 3rd conditional (hypothèse contraire au passé).
\end{itemize}
Dans la proposition principale, le futur français correspond à \eng{will}, le conditionnel français à \eng{would}, et le conditionnel passé français à \eng{would have + participe passé}.
\end{propriete}

\begin{center}
\begin{tabularx}{\linewidth}{|X|X|X|}
\hline
\rowcolor{popblueL} Français & English & Remarque \\
\hline
Si + présent & If + present simple & Correspondance exacte \\
\hline
Si + imparfait & If + past simple & L'imparfait français = le prétérit anglais \\
\hline
Si + plus-que-parfait & If + past perfect & Correspondance exacte \\
\hline
Futur (je serai) & will + base verbale & Jamais de futur après \eng{if} \\
\hline
Conditionnel (je serais) & would + base verbale & Jamais de \eng{would} après \eng{if} \\
\hline
Conditionnel passé (j'aurais été) & would have + p.p. & Ne pas oublier \eng{have} \\
\hline
\end{tabularx}
\end{center}

\begin{popfigure}
\begin{center}
\begin{tikzpicture}[
  box/.style={draw=popblue, fill=popblueL, rounded corners=3pt, align=center, minimum height=1.1cm, text width=3.6cm, font=\small},
  head/.style={draw=popdark, fill=popdark, text=white, rounded corners=3pt, align=center, minimum height=0.9cm, text width=3.6cm, font=\small\bfseries}
]
\node[head] (h1) at (0,0) {Type};
\node[head] (h2) at (4.4,0) {Structure};
\node[head] (h3) at (8.8,0) {Traduction française};
\node[box] (a1) at (0,-1.3) {1st conditional};
\node[box, align=center] (a2) at (4.4,-1.3){If + present,\\ will + V};
\node[box, align=center] (a3) at (8.8,-1.3){Si + présent,\\ futur};
\node[box] (b1) at (0,-2.7) {2nd conditional};
\node[box, align=center] (b2) at (4.4,-2.7){If + past simple,\\ would + V};
\node[box, align=center] (b3) at (8.8,-2.7){Si + imparfait,\\ conditionnel};
\node[box] (c1) at (0,-4.1) {3rd conditional};
\node[box, align=center] (c2) at (4.4,-4.1){If + past perfect,\\ would have + p.p.};
\node[box, align=center] (c3) at (8.8,-4.1){Si + plus-que-parfait,\\ conditionnel passé};
\end{tikzpicture}
\end{center}
\legende{Correspondance des trois conditionnels entre l'anglais et le français}
\end{popfigure}

\courssub{La différence majeure : jamais de futur ni de would après if}

Voici maintenant la différence cruciale. En français, le mot « si » peut introduire une \textbf{question indirecte}, et dans ce cas on utilise le futur ou le conditionnel après « si » : « Je me demande s'il \textbf{viendra}. » « Je ne sais pas si elle \textbf{accepterait}. »

En anglais, c'est impossible. Après \eng{if}, on n'utilise \textbf{jamais} \eng{will} ni \eng{would}, que ce soit dans une condition ou dans une question indirecte :

\begin{exemplebox}[Questions indirectes : attention au temps]
\eng{I wonder if he will come.} « Je me demande s'il viendra. » — ici \eng{will} est correct car ce \eng{if} signifie « si » au sens de « whether » (question indirecte), pas une condition.\par
\eng{If it rains tomorrow, we will stay at home.} « S'il pleut demain, nous resterons à la maison. » — condition : présent après \eng{if}, même si l'action est future.
\end{exemplebox}

\begin{attention}[Piège n°1 : « If it will rain »]
L'élève francophone pense : « demain = futur, donc \eng{will} ». Faux ! Dans une \textbf{condition}, le futur est exprimé par le présent après \eng{if}.\par
Incorrect : \eng{If it \textbf{will rain}, we will cancel the match.}\par
Correct : \eng{If it \textbf{rains}, we will cancel the match.} « S'il pleut, nous annulerons le match. »\par
Même piège avec \eng{when} : \eng{When I \textbf{arrive} (pas \emph{will arrive}), I will call you.} « Quand j'arriverai, je t'appellerai. »
\end{attention}

\begin{attention}[Piège n°2 : « If I would be rich »]
En français familier, on entend parfois « si j'aurais su » ou « si je serais riche » — c'est déjà une faute en français soigné, mais elle devient catastrophique en anglais. On ne met \textbf{jamais} \eng{would} dans la proposition introduite par \eng{if}.\par
Incorrect : \eng{If I \textbf{would be} rich, I would build a hospital.}\par
Correct : \eng{If I \textbf{were} rich, I would build a hospital.} « Si j'étais riche, je construirais un hôpital. »\par
Incorrect : \eng{If I \textbf{would have known}, I would have helped.}\par
Correct : \eng{If I \textbf{had known}, I would have helped.} « Si j'avais su, j'aurais aidé. »
\end{attention}

\courssub{Piège n°3 : oublier « have » dans le 3rd conditional}

Le 3rd conditional exige \eng{would have + participe passé} dans la proposition principale. Sous la pression de l'oral, beaucoup d'élèves écrivent \eng{would passed}, \eng{would gone}, \eng{would been} — sans le \eng{have}. C'est une erreur grave qui détruit la structure.

\begin{attention}[« would passed » n'existe pas]
Incorrect : \eng{If she had studied, she \textbf{would passed} the exam.}\par
Correct : \eng{If she had studied, she \textbf{would have passed} the exam.} « Si elle avait étudié, elle aurait réussi l'examen. »\par
Astuce : à l'oral, \eng{would have} se contracte en \eng{would've} (prononcé « wou-dèv »), ce qui explique que certains croient entendre « would of ». On n'écrit \textbf{jamais} \eng{would of} !
\end{attention}

\begin{exemplebox}[3rd conditional correctement formé]
\eng{If the farmers had used fertilisers, the harvest would have been better.} « Si les agriculteurs avaient utilisé des engrais, la récolte aurait été meilleure. »\par
\eng{We would have visited the hospital if we had had time.} « Nous aurions visité l'hôpital si nous avions eu le temps. »\par
\eng{If they had left Douala earlier, they would not have missed the bus.} « S'ils étaient partis de Douala plus tôt, ils n'auraient pas manqué le bus. »
\end{exemplebox}

\courssub{Piège n°4 : confondre when (certitude) et if (hypothèse)}

Le français utilise souvent « quand » et « si » de manière interchangeable dans la tête des élèves, mais l'anglais les distingue strictement :

\begin{definition}[When ou if ?]
\eng{When} exprime une \cle{certitude} : l'événement aura lieu, on ne sait juste pas exactement quand.\par
\eng{If} exprime une \cle{hypothèse} : l'événement peut arriver ou ne pas arriver.
\end{definition}

\begin{exemplebox}[When ≠ If]
\eng{When I grow up, I will become a doctor.} « Quand je serai grand, je deviendrai médecin. » — l'enfant grandira forcément : certitude.\par
\eng{If it rains this evening, the match will be postponed.} « S'il pleut ce soir, le match sera reporté. » — il peut pleuvoir ou non : hypothèse.\par
\eng{When the rains come, the farmers will plant maize.} « Quand les pluies arriveront, les agriculteurs planteront le maïs. » — les pluies viennent chaque année : certitude.\par
\eng{If the rains come early, the farmers will be happy.} « Si les pluies arrivent tôt, les agriculteurs seront contents. » — hypothèse.
\end{exemplebox}

Remarquez que dans les deux cas, on utilise le \textbf{présent} après \eng{when} et après \eng{if}, jamais le futur : \eng{When I \textbf{am} older} (pas \emph{when I will be older}).

\courssub{Méthode : choisir le bon conditional}

\begin{methode}[Trois questions pour choisir la bonne structure]
\begin{enumerate}
\item La situation est-elle \textbf{possible, réelle}, envisagée pour l'avenir ? → utilisez le \cle{1st conditional} : \eng{If + present, will + V}. Exemple : \eng{If we recycle plastic, we will protect the environment.}
\item La situation est-elle \textbf{imaginaire, irréelle} au présent (un rêve, une hypothèse contraire à la réalité actuelle) ? → utilisez le \cle{2nd conditional} : \eng{If + past simple, would + V}. Exemple : \eng{If I were the Minister of Health, I would build more hospitals.}
\item La situation est-elle \textbf{contraire au passé} (un regret, un reproche sur ce qui ne s'est pas passé) ? → utilisez le \cle{3rd conditional} : \eng{If + past perfect, would have + p.p.}. Exemple : \eng{If we had vaccinated earlier, fewer people would have fallen ill.}
\end{enumerate}
Réflexe de traduction : « si + présent » → 1st ; « si + imparfait » → 2nd ; « si + plus-que-parfait » → 3rd.
\end{methode}

\begin{exoresolu}[Choisissez et justifiez le bon conditional]
Énoncé — Complétez avec la forme correcte du verbe entre parenthèses, puis indiquez le type de conditional :\par
1. If the government (invest) in clean water, cholera cases (decrease).\par
2. If I (be) a nurse, I (work) in a rural health centre.\par
3. If they (plant) trees last year, the erosion (not / be) so serious.\par
4. « Si tu manges trop de sucre, tu tomberas malade. » — Traduisez.
\tcblower
Solution détaillée :\par
1. \eng{If the government \textbf{invests} in clean water, cholera cases \textbf{will decrease}.} — 1st conditional : mesure tout à fait possible, hypothèse réelle sur l'avenir. « Si le gouvernement investit dans l'eau potable, les cas de choléra diminueront. »\par
2. \eng{If I \textbf{were} a nurse, I \textbf{would work} in a rural health centre.} — 2nd conditional : je ne suis pas infirmier, situation imaginaire au présent. « Si j'étais infirmier, je travaillerais dans un centre de santé rural. »\par
3. \eng{If they \textbf{had planted} trees last year, the erosion \textbf{would not have been} so serious.} — 3rd conditional : ils n'ont pas planté d'arbres, regret sur le passé. « S'ils avaient planté des arbres l'année dernière, l'érosion n'aurait pas été si grave. »\par
4. \eng{If you \textbf{eat} too much sugar, you \textbf{will fall} ill.} — 1st conditional : « si + présent » en français, donc \eng{if + present} en anglais, et le futur français « tomberas » se traduit par \eng{will fall}.
\end{exoresolu}

\begin{exercice}[À vous de jouer]
Traduisez en anglais en choisissant le bon conditional :\par
1. Si nous protégeons la forêt, les animaux survivront.\par
2. Si j'avais plus de temps, je ferais du sport tous les jours.\par
3. Si elle avait consulté le médecin plus tôt, elle aurait guéri plus vite.\par
4. Quand les pluies arriveront, nous planterons du cacao.\par
5. Je me demande s'il pleuvra demain.
\end{exercice}

\begin{corrige}[Exercice : À vous de jouer]
1. \eng{If we protect the forest, the animals will survive.} — 1st conditional (« si + présent »).\par
2. \eng{If I had more time, I would do sport every day.} — 2nd conditional (« si + imparfait »).\par
3. \eng{If she had seen the doctor earlier, she would have recovered faster.} — 3rd conditional (« si + plus-que-parfait ») ; ne pas oublier \eng{have}.\par
4. \eng{When the rains come, we will plant cocoa.} — \eng{when} car les pluies sont une certitude saisonnière ; présent après \eng{when}.\par
5. \eng{I wonder if it will rain tomorrow.} — question indirecte : \eng{will} est permis car ce \eng{if} n'exprime pas une condition.
\end{corrige}

\begin{aretenir}[L'essentiel de la comparaison]
\begin{itemize}
\item Les temps après « si » correspondent parfaitement : présent = \eng{present}, imparfait = \eng{past simple}, plus-que-parfait = \eng{past perfect}.
\item Jamais de \eng{will} ni de \eng{would} dans la proposition en \eng{if} de condition : \eng{If it rains}, \eng{If I were}, \eng{If I had known}.
\item Le 3rd conditional exige \eng{would have + participe passé} : jamais \eng{would passed}.
\item \eng{When} = certitude, \eng{if} = hypothèse ; dans les deux cas, présent après la conjonction.
\item Exception : dans les questions indirectes (\eng{I wonder if...}), \eng{will} et \eng{would} sont possibles après \eng{if}.
\end{itemize}
\end{aretenir}

\coursec{Entraînement guidé et production}

Nous savons maintenant distinguer les trois structures conditionnelles et repérer les pièges de traduction entre le français et l'anglais. Il est temps de passer à la pratique : c'est en manipulant les formes, en transformant les phrases et en produisant nos propres énoncés que la règle devient un réflexe. Cette dernière partie alterne exercices de reconnaissance, transformations, traductions et production libre sur le thème du chapitre : l'environnement au Cameroun.

\courssub{Exercice 1 : mettre le verbe au bon temps (les trois types mélangés)}

Avant de conjuguer, pose-toi toujours la même question : \emph{la situation est-elle réelle, imaginaire au présent, ou imaginaire au passé ?} La figure suivante résume cette démarche de décision.

\begin{popfigure}
\begin{center}
\begin{tikzpicture}[
  node distance=1.2cm and 1.5cm,
  decision/.style={diamond, draw=popblue, fill=popblueL, text width=3.4cm, align=center, inner sep=2pt, font=\small, aspect=2},
  box/.style={rectangle, rounded corners, draw=popdark, fill=popcream, text width=3.8cm, align=center, font=\small},
  arr/.style={-{Stealth}, thick, popdark},
  label/.style={font=\small, text=popdark, midway}
]
\node[decision] (q1) {La situation est-elle réelle et probable ?};
\node[box, below left=of q1, xshift=-1cm, align=center] (first){\textbf{1st conditional}\\ If + present,\\ will + V};
\node[decision, below right=of q1, xshift=1cm] (q2) {Imaginaire : présent ou passé ?};
\node[box, below left=of q2, xshift=-1cm, align=center] (second){\textbf{2nd conditional}\\ If + past,\\ would + V};
\node[box, below right=of q2, xshift=1cm, align=center] (third){\textbf{3rd conditional}\\ If + past perfect,\\ would have + p.p.};
\draw[arr] (q1) -- node[label, above left]{oui} (first);
\draw[arr] (q1) -- node[label, above right]{non} (q2);
\draw[arr] (q2) -- node[label, above left]{présent} (second);
\draw[arr] (q2) -- node[label, above right]{passé} (third);
\end{tikzpicture}
\end{center}
\legende{Organigramme de décision : quel conditional choisir ?}
\end{popfigure}

\begin{exoresolu}[Exercice 1 — Complète avec le verbe au bon temps]
Énoncé. Mets le verbe entre parenthèses à la forme correcte.
\begin{enumerate}
\item If it rains tomorrow, the football match \ldots\ (be) cancelled.
\item If I were rich, I \ldots\ (plant) a thousand trees in my village.
\item If the fishermen had listened to the warning, they \ldots\ (not / lose) their nets.
\item If we recycle plastic bottles, we \ldots\ (protect) the rivers of Douala.
\item If she \ldots\ (study) medicine, she would become a doctor.
\end{enumerate}
\tcblower
Solution détaillée.
\begin{enumerate}
\item \eng{will be} — 1st conditional : la pluie de demain est une possibilité réelle (\eng{if} + présent, \eng{will} + V). La forme passive \eng{will be cancelled} est attendue.
\item \eng{would plant} — 2nd conditional : « si j'étais riche », hypothèse imaginaire au présent (signal : \eng{were}).
\item \eng{would not have lost} — 3rd conditional : signal \eng{had listened} (past perfect), donc \eng{would have} + participe passé (\eng{lost}).
\item \eng{will protect} — 1st conditional : après \eng{if} + présent, la principale prend \eng{will} + V pour exprimer une conséquence future probable. \emph{Note : le présent simple (\eng{protect}) n'est utilisé que pour des vérités générales (Zero conditional), pas pour une action future spécifique.}
\item \eng{studied} — 2nd conditional : la principale contient \eng{would become}, donc la subordonnée prend le prétérit (\eng{studied}).
\end{enumerate}
\end{exoresolu}

\begin{exercice}[À toi de jouer — Exercice 1 (suite)]
Mets le verbe entre parenthèses à la forme correcte.
\begin{enumerate}
\item If the government builds more dams, the region \ldots\ (have) electricity.
\item If people stopped burning bushes, the soil \ldots\ (stay) fertile.
\item If we had planted trees last year, the schoolyard \ldots\ (be) cooler today.
\item If you throw rubbish in the gutter, it \ldots\ (block) the drains.
\item If I \ldots\ (be) the Minister of Health, I would build more hospitals.
\item If they had left earlier, they \ldots\ (not / miss) the bus to Bamenda.
\end{enumerate}
\end{exercice}

\begin{corrige}[Exercice 1 (suite)]
1. \eng{will have} (1st conditional, situation réelle future). 2. \eng{would stay} (2nd conditional, signal \eng{stopped} = prétérit). 3. \eng{would be} (3rd conditional avec conséquence au présent : \eng{had planted} $\rightarrow$ \eng{would be} today). 4. \eng{will block} (1st conditional, conséquence future d'une action présente ; \eng{blocks} serait un Zero conditional, vérité générale). 5. \eng{were} (2nd conditional, \eng{were} pour toutes les personnes dans le style soigné). 6. \eng{would not have missed} (3rd conditional, \eng{had left} = past perfect).
\end{corrige}

\courssub{Exercice 2 : transformation — du 1st au 2nd conditional}

Transformer une phrase du 1st au 2nd conditional n'est pas un simple exercice mécanique : le sens change. Le 1st présente l'événement comme probable ; le 2nd le recule dans l'imaginaire.

\begin{methode}[Transformer : du 1st au 2nd conditional]
\begin{enumerate}
\item Repère les deux verbes : celui de la subordonnée (\eng{if} + présent) et celui de la principale (\eng{will} + V).
\item Dans la subordonnée, remplace le présent par le prétérit simple : \eng{don't act} devient \eng{didn't act}.
\item Dans la principale, remplace \eng{will} par \eng{would} : \eng{will change} devient \eng{would change}.
\item Observe le changement de sens : la phrase passe de « probable » à « imaginaire, peu probable ».
\end{enumerate}
\end{methode}

\begin{exemplebox}[Exemple traité pas à pas]
\eng{If we don't act, the climate will change.} « Si nous n'agissons pas, le climat va changer. » (1st : danger réel, probable.)

Transformation : \eng{If we didn't act, the climate would change.} « Si nous n'agissions pas, le climat changerait. » (2nd : hypothèse plus théorique, comme si l'inaction était peu envisageable ou purement spéculative.)

Autre exemple : \eng{If it rains, we will stay at home.} $\rightarrow$ \eng{If it rained, we would stay at home.} « S'il pleuvait, on resterait à la maison » (mais le ciel est dégagé).
\end{exemplebox}

\begin{exercice}[À toi de jouer — Exercice 2]
Transforme ces phrases du 1st au 2nd conditional, puis traduis les deux versions.
\begin{enumerate}
\item If the factory pollutes the river, the villagers will complain.
\item If she works hard, she will pass the Probatoire.
\item If we plant trees, the air will be cleaner.
\end{enumerate}
\end{exercice}


\begin{corrige}[Exercice 2]
1. \eng{If the factory polluted the river, the villagers would complain.} « Si l'usine polluait la rivière, les villageois se plaindraient. » 2. \eng{If she worked hard, she would pass the Probatoire.} « Si elle travaillait dur, elle réussirait le Probatoire. » 3. \eng{If we planted trees, the air would be cleaner.} « Si nous plantions des arbres, l'air serait plus pur. » Dans chaque cas, l'événement devient moins probable, plus hypothétique.
\end{corrige}

\courssub{Exercice 3 : traduction français $\rightarrow$ anglais}

\begin{attention}[Avant de traduire : regarde le temps après « si »]
En français, le temps qui suit \emph{si} te donne la clé : présent $\rightarrow$ 1st conditional anglais ; imparfait $\rightarrow$ 2nd ; plus-que-parfait $\rightarrow$ 3rd. Ne traduis jamais mot à mot : le futur et le conditionnel français après \emph{si} n'existent pas en anglais dans la subordonnée. « Si j'aurais su » (fautif même en français) ne peut pas devenir \eng{if I would have known}.
\end{attention}

\begin{center}
\begin{tabularx}{\linewidth}{|X|X|X|}
\hline
\rowcolor{popblueL} Français (temps après « si ») & English & Type \\
\hline
S'il pleut (présent) & If it rains (present) & 1st \\
\hline
S'il pleuvait (imparfait) & If it rained (past) & 2nd \\
\hline
S'il avait plu (plus-que-parfait) & If it had rained (past perfect) & 3rd \\
\hline
\end{tabularx}
\end{center}

\begin{exoresolu}[Exercice 3 — Traduis en anglais]
Énoncé. Traduis : « Si nous avions protégé la forêt, les pluies seraient revenues. »
\tcblower
Solution détaillée. Après \emph{si}, le français utilise le plus-que-parfait (« avions protégé ») : il s'agit d'un regret sur le passé, donc 3rd conditional. La subordonnée : \eng{If we had protected the forest}. La principale française est au conditionnel présent (« seraient revenues »), mais la logique anglaise exige \eng{would have} + participe passé : \eng{the rains would have come back}. Phrase complète : \eng{If we had protected the forest, the rains would have come back.}
\end{exoresolu}

\begin{exercice}[À toi de jouer — Exercice 3]
Traduis en anglais.
\begin{enumerate}
\item Si tu éteins la lumière, tu économiseras de l'énergie.
\item Si j'étais toi, je planterais des arbres autour de l'école. (corrige d'abord le français si nécessaire)
\item Si les élèves avaient nettoyé la cour, le directeur aurait été content.
\item Si le gouvernement construit des pistes cyclables, moins de gens prendront des bendskins.
\end{enumerate}
\end{exercice}


\begin{corrige}[Exercice 3]
1. \eng{If you switch off the light, you will save energy.} 2. Français correct : « Si j'étais toi, je planterais des arbres… » $\rightarrow$ \eng{If I were you, I would plant trees around the school.} 3. \eng{If the students had cleaned the yard, the headmaster would have been pleased.} 4. \eng{If the government builds cycle lanes, fewer people will take bendskins.}
\end{corrige}

\courssub{Exercice 4 : corrige les phrases fautives}

Voici des erreurs authentiques d'élèves francophones. Pour chacune, identifie la faute, corrige-la et justifie.

\begin{exoresolu}[Exercice 4 — Trouve et corrige l'erreur]
Énoncé. \eng{If I would have money, I will buy a bicycle.}
\tcblower
Solution détaillée. Deux fautes. Premièrement, \eng{would} ne se met jamais dans la subordonnée introduite par \eng{if} : il faut le prétérit \eng{had}. Deuxièmement, si la subordonnée est au prétérit (2nd conditional), la principale prend \eng{would}, pas \eng{will}. Correction : \eng{If I had money, I would buy a bicycle.} « Si j'avais de l'argent, j'achèterais un vélo. »
\end{exoresolu}

\begin{exercice}[À toi de jouer — Exercice 4]
Corrige chaque phrase et justifie.
\begin{enumerate}
\item If it will rain, we will cancel the match.
\item If I was you, I will stop smoking.
\item If they had came earlier, they would have seen the doctor.
\item If we don't hurry, we would miss the bus.
\end{enumerate}
\end{exercice}


\begin{corrige}[Exercice 4]
1. \eng{If it rains, we will cancel the match.} — jamais de \eng{will} après \eng{if} (subordonnée au présent). 2. \eng{If I were you, I would stop smoking.} — \eng{were} (style soigné, subjonctif) et \eng{would} dans la principale (2nd conditional). 3. \eng{If they had come earlier…} — après \eng{had}, le participe passé : \eng{come}, pas \eng{came} (forme irrégulière : come / came / come). 4. \eng{If we don't hurry, we will miss the bus.} — 1st conditional : présent + \eng{will} (cohérence des temps).
\end{corrige}

\courssub{Production guidée : six phrases sur l'environnement au Cameroun}

\begin{experience}[Rédaction guidée — My environment, my conditionals]
Rédige six phrases complètes sur l'environnement au Cameroun : deux de chaque type.
\begin{itemize}
\item \textbf{1st conditional} (actions réelles et proches) : \eng{If we clean the gutters before the rainy season, the floods will be less serious.}
\item \textbf{2nd conditional} (rêves et projets imaginaires) : \eng{If Douala had more parks, the air would be fresher.}
\item \textbf{3rd conditional} (regrets sur le passé) : \eng{If people had not cut the mangroves, the coast would have been better protected.}
\end{itemize}
Contraintes : utilise au moins quatre mots du lexique du chapitre (\eng{pollution, waste, to recycle, drought, flood, to protect, climate}) ; souligne les deux verbes de chaque phrase ; fais relire ton travail par un voisin qui vérifiera la structure avec l'organigramme.
\end{experience}

\courssub{Jeu de rôle : If I were the Mayor of Douala...}

\begin{experience}[Speaking — If I were the Mayor of Douala]
La classe simule un conseil municipal. Chaque élève propose une mesure pour la ville en commençant obligatoirement par : \eng{If I were the Mayor of Douala, I would...}

Exemples de mesures : \eng{If I were the Mayor of Douala, I would build more public toilets in the markets.} / \eng{...I would ban plastic bags.} / \eng{...I would create a bicycle lane on every main road.}

Variante pour aller plus loin : un deuxième élève réagit avec un 1st conditional (\eng{If you do that, people will be healthier!}) et un troisième avec un 3rd conditional de regret (\eng{If the previous mayor had done that, the city would have been cleaner.}). La classe vote ensuite pour les trois meilleures mesures.
\end{experience}

\begin{aretenir}[L'essentiel de la leçon]
\begin{itemize}
\item 1st conditional : \eng{If + present, will + V} — situation réelle, probable.
\item 2nd conditional : \eng{If + past, would + V} — situation imaginaire au présent ; \eng{were} pour toutes les personnes.
\item 3rd conditional : \eng{If + past perfect, would have + participe passé} — regret sur le passé.
\item Jamais de \eng{will} ni de \eng{would} dans la subordonnée en \eng{if}.
\item Pour traduire, regarde d'abord le temps français après \emph{si} : présent $\rightarrow$ 1st, imparfait $\rightarrow$ 2nd, plus-que-parfait $\rightarrow$ 3rd.
\end{itemize}
\end{aretenir}

\newpage

\coursec{Activité d'intégration}

\coursec{Lesson 32 — Grammar: 1st, 2nd, and 3rd Conditionals}

\courssub{Objectifs de l'activité}

Cette activité d'intégration vise à mobiliser l'ensemble des connaissances acquises sur les trois types de conditionnels (\eng{1st}, \eng{2nd}, \eng{3rd conditional}) dans une tâche de production écrite collaborative, ancrée dans une problématique concrète et citoyenne : l'environnement au sein de l'établissement scolaire.

À l'issue de cette activité, l'élève sera capable de :
\begin{itemize}
    \item \textbf{Identifier} le type d'hypothèse véhiculé (réelle, irréelle présente, regret passé) pour choisir la structure conditionnelle appropriée.
    \item \textbf{Construire} des phrases grammaticalement correctes en respectant l'accord des temps (\eng{If + present simple $\rightarrow$ will + BV} ; \eng{If + preterit $\rightarrow$ would + BV} ; \eng{If + past perfect $\rightarrow$ would have + V-en}).
    \item \textbf{Manipuler} le vocabulaire thématique de l'environnement (\eng{waste management, recycling, biodiversity, sustainable development, carbon footprint}) et des engagements citoyens (\eng{pledge, commit to, regret, wish}).
    \item \textbf{Coopérer} en groupe pour négocier le sens, corriger les erreurs entre pairs et produire un texte cohérent de 9 phrases.
    \item \textbf{Présenter} oralement le fruit du travail collectif et argumenter ses choix linguistiques.
\end{itemize}

\courssub{Situation}

\begin{textebox}{Contexte : Le projet « Green School » au Lycée Bilingue de Buea}
Le Lycée Bilingue de Buea, situé au pied du Mont Cameroun, lance une initiative intitulée \eng{“Green School Challenge”} pour sensibiliser sa communauté éducative aux enjeux écologiques. L'administration, en partenariat avec le club \eng{“Eco-Guards”} et l'ONG locale \eng{“Green Cameroon”}, invite chaque classe de Terminale à rédiger une \textbf{Green Charter} (Charte Verte) qui sera affichée dans le hall principal et soumise au vote de l'ensemble des élèves et du personnel.

Votre classe est divisée en groupes de 4 élèves. Chaque groupe représente une « commission environnementale » chargée de proposer une charte en \textbf{9 engagements formulés en anglais}. La charte doit refléter trois dimensions temporelles et modales distinctes pour montrer la maturité de la réflexion des élèves :
\begin{enumerate}
    \item \textbf{3 engagements concrets et réalisables immédiatement} (hypothèses réelles $\rightarrow$ \eng{1st Conditional}).
    \item \textbf{3 projets ambitieux ou rêves pour l'école idéale} (hypothèses irréelles présentes $\rightarrow$ \eng{2nd Conditional}).
    \item \textbf{3 constats de regrets sur des occasions manquées par le passé} (hypothèses passées non réalisées $\rightarrow$ \eng{3rd Conditional}).
\end{enumerate}

Le document final sera présenté lors d'une assemblée simulée (\eng{“Town Hall Meeting”}) où chaque groupe lit sa charte. La meilleure charte — celle qui allie justesse grammaticale, richesse lexicale, pertinence écologique et force de persuasion — sera laminée et affichée durablement dans l'établissement.
\end{textebox}

\courssub{Consignes}

\textbf{Étape 1 : Analyse de la situation et remobilisation lexicale (15 min)}\par
\begin{enumerate}
    \item En groupe, relevez dans le texte de situation les mots-clés liés à l'environnement et à l'action citoyenne. Complétez le tableau ci-dessous en ajoutant la traduction et la nature grammaticale.
    \begin{center}
    \begin{tabularx}{\linewidth}{|X|X|X|X|}
    \hline
    \rowcolor{popblueL} \textbf{Mot anglais} & \textbf{Nature} & \textbf{Traduction} & \textbf{Collocation utile} \\
    \hline
    Waste management & Noun (U) & Gestion des déchets & \eng{improve waste management} \\
    \hline
    Recycling bins & Noun (pl.) & Bacs de recyclage & \eng{install recycling bins} \\
    \hline
    Biodiversity & Noun (U) & Biodiversité & \eng{protect biodiversity} \\
    \hline
    Carbon footprint & Noun (C) & Empreinte carbone & \eng{reduce our carbon footprint} \\
    \hline
    Sustainable & Adjective & Durable & \eng{sustainable development} \\
    \hline
    Pledge (to) & Verb & S'engager à & \eng{pledge to plant trees} \\
    \hline
    Missed opportunity & Noun (C) & Occasion manquée & \eng{a missed opportunity} \\
    \hline
    \end{tabularx}
    \end{center}
    \item Discutez en français : Quelles actions concrètes peut-on mener \textit{demain matin} au lycée ? Quels rêves fous aimeriez-vous voir se réaliser ? Quelles actions n'avez-vous \textbf{pas} menées l'an dernier et que vous regrettez ?
\end{enumerate}

\textbf{Étape 2 : Rédaction collaborative de la Charte (30 min)}\par
\begin{enumerate}
    \setcounter{enumi}{2}
    \item Rédigez vos 9 phrases sur une grande feuille A3 (format affiche). Respectez scrupuleusement la structure imposée :
    \begin{itemize}
        \item \textbf{Phrases 1--3 (1st Conditional) :} \eng{If + Present Simple, ... will + Base Verb}. 
        \textit{Exemple :} \eng{If we separate plastic from organic waste, the school grounds will stay clean.}
        \item \textbf{Phrases 4--6 (2nd Conditional) :} \eng{If + Preterit, ... would + Base Verb}. 
        \textit{Exemple :} \eng{If the school had a compost system, we would produce fertilizer for a garden.}
        \item \textbf{Phrases 7--9 (3rd Conditional) :} \eng{If + Past Perfect, ... would have + Past Participle}. 
        \textit{Exemple :} \eng{If we had started the recycling program last year, we would have saved hundreds of plastic bottles.}
    \end{itemize}
    \item \textbf{Contrainte de cohérence :} Les 9 phrases doivent former un tout logique (pas 9 idées disjointes). Utilisez des connecteurs logiques (\eng{Moreover, Furthermore, However, Consequently}) pour lier les phrases entre elles.
    \item \textbf{Auto-correction guidée :} Avant de finaliser, chaque groupe vérifie sa production à l'aide de la grille ci-dessous (un élève = un vérificateur par item) :
    \begin{itemize}
        \item [ ] Vérificateur « Structure » : Le \eng{If} est-il suivi du bon temps ? Le principal clause utilise-t-il le bon modal (\eng{will/would/would have}) ?
        \item [ ] Vérificateur « Verbes irréguliers » : Les participes passés sont-ils corrects ? (\eng{begin $\rightarrow$ begun, do $\rightarrow$ done, see $\rightarrow$ seen, write $\rightarrow$ written}).
        \item [ ] Vérificateur « Vocabulaire » : Au moins 5 mots du tableau de l'étape 1 sont-ils utilisés ?
        \item [ ] Vérificateur « Ponctuation/Majuscules » : Virgule après la proposition \eng{If} ? Majuscules aux noms propres (\eng{Cameroon, Buea, Mount Cameroon}) ?
    \end{itemize}
\end{enumerate}

\textbf{Étape 3 : Présentation orale et vote (15 min)}\par
\begin{enumerate}
    \setcounter{enumi}{5}
    \item Chaque groupe désigne un \eng{spokesperson} (porte-parole) qui lit la charte devant la classe (simulation d'assemblée). Les autres membres peuvent intervenir pour expliquer un choix de vocabulaire ou de structure.
    \item La classe vote à bulletin secret sur trois critères (notation sur 5) :
    \begin{enumerate}
        \item Exactitude grammaticale (0 faute = 5/5, 1--2 fautes = 4/5, etc.).
        \item Richesse lexicale et pertinence écologique.
        \item Originalité et faisabilité des propositions.
    \end{enumerate}
    \item L'enseignant annonce le groupe gagnant, commente les points forts linguistiques de la charte victorieuse et corrige collectivement les erreurs récurrentes relevées lors des présentations.
\end{enumerate}

\courssub{Ressources à mobiliser}

\begin{propriete}[Rappel des trois structures conditionnelles]
\centering
\begin{tabularx}{\linewidth}{|X|X|X|X|}
\hline
\rowcolor{popblueL} \textbf{Type} & \textbf{Structure (If-clause $\rightarrow$ Main clause)} & \textbf{Sens / Emploi} & \textbf{Marqueurs temporels fréquents} \\
\hline
\textbf{1st Conditional} & \eng{If + Present Simple $\rightarrow$ will + BV} & Hypothèse \textbf{réelle}, possible dans le futur proche. Engagement, prédiction, avertissement. & \eng{tomorrow, next week, if we act now, as soon as} \\
\hline
\textbf{2nd Conditional} & \eng{If + Preterit $\rightarrow$ would + BV} & Hypothèse \textbf{irréelle} au présent/futur. Rêve, conseil, situation imaginaire. & \eng{now, today, if I were you, in an ideal world} \\
\hline
\textbf{3rd Conditional} & \eng{If + Past Perfect $\rightarrow$ would have + V-en} & Hypothèse \textbf{passée non réalisée}. Regret, critique, spéculation sur le passé. & \eng{yesterday, last year, previously, if we had known} \\
\hline
\end{tabularx}
\end{propriete}

\begin{propriete}[Comparaison Français / Anglais — Structures conditionnelles]
\centering
\begin{tabularx}{\linewidth}{|X|X|X|}
\hline
\rowcolor{popblueL} \textbf{Français (Structure standard)} & \textbf{English (Structure standard)} & \textbf{Remarque cruciale pour francophones} \\
\hline
Si + \textbf{présent de l'indicatif} $\rightarrow$ \textbf{futur simple} & \eng{If + Present Simple $\rightarrow$ will + BV} & En anglais, \textbf{jamais de futur} (will) après \eng{if}. Le présent simple a valeur de futur. \\
\hline
Si + \textbf{imparfait} $\rightarrow$ \textbf{conditionnel présent} & \eng{If + Preterit $\rightarrow$ would + BV} & Le \textbf{preterit} anglais n'est \textbf{pas} un passé : il marque l'irréel du présent. Pas de \eng{would} après \eng{if}. \\
\hline
Si + \textbf{plus-que-parfait} $\rightarrow$ \textbf{conditionnel passé} & \eng{If + Past Perfect $\rightarrow$ would have + V-en} & \eng{Had + V-en} (Past Perfect) pour l'antériorité passée. \eng{Would have + V-en} pour le regret/résultat passé. \\
\hline
\end{tabularx}
\end{propriete}

\begin{attention}[Pièges majeurs pour francophones — Check-list de révision]
\begin{itemize}
    \item \textbf{Piège 1 : \eng{Will} / \eng{Would} dans la proposition \eng{If}.} \eng{*If we will sort...} $\rightarrow$ \textbf{FAUX}. Jamais de \eng{will/would} après \eng{if}. \eng{If we sort...} (Present simple).
    \item \textbf{Piège 2 : Le \eng{Preterit} au 2nd conditional n'est pas un passé.} \eng{If we had a garden...} ne signifie pas « Si nous \textbf{avions eu} un jardin » (passé) mais « Si nous \textbf{avions} un jardin » (présent imaginaire). Pour le passé, on utilise le \eng{Past Perfect} (3rd conditional).
    \item \textbf{Piège 3 : \eng{Would have} + \textbf{Participe Passé} (pas infinitif).} \eng{*We would have save} $\rightarrow$ \textbf{FAUX}. \eng{We would have saved}.
    \item \textbf{Piège 4 : Verbes irréguliers au 3rd conditional.} \eng{Begin $\rightarrow$ Begun}, \eng{Do $\rightarrow$ Done}, \eng{Go $\rightarrow$ Gone}, \eng{See $\rightarrow$ Seen}, \eng{Write $\rightarrow$ Written}, \eng{Take $\rightarrow$ Taken}, \eng{Make $\rightarrow$ Made}. Apprenez-les par cœur.
    \item \textbf{Piège 5 : \eng{If I was} vs \eng{If I were}.} À l'écrit formel et aux examens, préférez \eng{\textbf{were}} pour toutes les personnes (\eng{If I were you, If he were here}). \eng{Was} est toléré à l'oral informel seulement.
    \item \textbf{Piège 6 : Virgule.} Virgule obligatoire si la phrase commence par \eng{If...} (\eng{If we act now, we will succeed.}) ; pas de virgule si \eng{If} est au milieu (\eng{We will succeed if we act now.}).
    \item \textbf{Piège 7 : \eng{Would} contracté en \eng{'d}.} À l'oral et à l'écrit informel : \eng{I'd, you'd, he'd} = \eng{I would, you would, he would}. Ne pas confondre avec \eng{I had} (aussi \eng{I'd}) : le contexte (BV vs V-en) tranche.
\end{itemize}
\end{attention}

\begin{definition}[Vocabulaire thématique utile — Environnement \& Action]
\begin{itemize}
    \item \eng{To sort / separate waste} : trier les déchets.
    \item \eng{Recycling bins / containers} : bacs de recyclage.
    \item \eng{Compost (n/v)} : compost / composter.
    \item \eng{Organic waste} : déchets organiques.
    \item \eng{Carbon footprint} : empreinte carbone.
    \item \eng{To reduce / cut down on} : réduire.
    \item \eng{To pledge to + BV} : s'engager à.
    \item \eng{A missed opportunity} : une occasion manquée.
    \item \eng{Sustainable / eco-friendly} : durable / respectueux de l'environnement.
    \item \eng{To raise awareness} : sensibiliser.
    \item \eng{Single-use plastic} : plastique à usage unique.
    \item \eng{Reusable} : réutilisable.
\end{itemize}
\end{definition}

\courssub{Proposition de production et corrigé commenté}

\begin{exoresolu}[Modèle de « Green Charter » — Groupe « Eco-Leaders »]
\textbf{Production écrite (Modèle élève) :}

\eng{Green Charter of the Terminale A Class — Lycée Bilingue de Buea}

\eng{1. If we sort our waste into the new recycling bins every day, the school compound will become much cleaner.}
\eng{2. If the administration installs water fountains, students will stop buying single-use plastic bottles.}
\eng{3. If we organize a monthly “Clean-Up Day”, we will protect the biodiversity around Mount Cameroon.}

\eng{4. If our school had a large vegetable garden, we would grow organic food for the canteen.}
\eng{5. If solar panels covered the roof of the main building, we would reduce our carbon footprint significantly.}
\eng{6. If every student had a reusable water bottle, we would eliminate plastic waste completely.}

\eng{7. If we had started this program last year, we would have saved hundreds of trees used for paper.}
\eng{8. If the previous students had planted more trees, the courtyard would have provided more shade over the years.}
\eng{9. If we had raised awareness earlier, we would have changed bad habits before they became routine.}

\tcblower

\textbf{Commentaire linguistique détaillé :}

\begin{itemize}
    \item \textbf{Phrases 1--3 (1st Conditional — Engagements réels) :}
    \begin{itemize}
        \item \eng{If we sort... will become} : Structure canonique. \eng{Sort} (present simple) $\rightarrow$ \eng{will become}. « Si nous trions..., le terrain de l'école deviendra bien plus propre. » Vocabulaire précis : \eng{recycling bins, school compound} (cour de récréation, terme camerounais courant).
        \item \eng{If the administration installs... will stop} : Sujet 3e personne du singulier (\eng{installs} avec \eng{-s}) correctement géré. \eng{Single-use plastic bottles} : collocation standard environnement. « Si l'administration installe des fontaines, les élèves cesseront d'acheter des bouteilles en plastique à usage unique. »
        \item \eng{If we organize... will protect} : \eng{Organize} (present simple). \eng{Clean-Up Day} : nom composé culturel (journée de nettoyage). \eng{Biodiversity around Mount Cameroon} : ancrage local fort. « Si nous organisons une "Journée Propreté" mensuelle, nous protégerons la biodiversité autour du Mont Cameroun. »
    \end{itemize}
    \item \textbf{Phrases 4--6 (2nd Conditional — Rêves/Projets irréels) :}
    \begin{itemize}
        \item \eng{If our school had... would grow} : \eng{Had} (preterit de \eng{have}) $\rightarrow$ \eng{would grow}. \eng{Vegetable garden, organic food, canteen} : champ lexical cohérent (agriculture scolaire). « Si notre école avait un grand potager, nous cultiverions de la nourriture bio pour la cantine. »
        \item \eng{If solar panels covered... would reduce} : \eng{Covered} (preterit régulier). Passif possible (\eng{were covered}) mais actif accepté pour la concision. \eng{Carbon footprint} : mot-clé du programme. « Si des panneaux solaires couvraient le toit du bâtiment principal, nous réduirions significativement notre empreinte carbone. »
        \item \eng{If every student had... would eliminate} : \eng{Every student} = singulier $\rightarrow$ \eng{had}. \eng{Reusable water bottle} : objet concret. \eng{Eliminate completely} : adverbe de degré bien placé. « Si chaque élève avait une gourde réutilisable, nous éliminerions complètement les déchets plastiques. »
    \end{itemize}
    \item \textbf{Phrases 7--9 (3rd Conditional — Regrets passés) :}
    \begin{itemize}
        \item \eng{If we had started... would have saved} : \eng{Had started} (Past Perfect) $\rightarrow$ \eng{would have saved} (Participe passé régulier \eng{saved}). \eng{Hundreds of trees} : quantification. « Si nous avions lancé ce programme l'an dernier, nous aurions sauvé des centaines d'arbres utilisés pour le papier. »
        \item \eng{If the previous students had planted... would have provided} : \eng{Had planted} $\rightarrow$ \eng{would have provided} (irrégulier \eng{provide $\rightarrow$ provided} régulier). \eng{Courtyard} (cour intérieure) $\rightarrow$ \eng{shade} (ombre). Phrase strictement au passé (pas de "today" qui créerait un conditionnel mixte). « Si les élèves d'avant avaient planté plus d'arbres, la cour aurait offert plus d'ombre au fil des ans. »
        \item \eng{If we had raised... would have changed} : \eng{Had raised} $\rightarrow$ \eng{would have changed}. \eng{Raise awareness} : collocation forte apprise en lexique. \eng{Before they became routine} : subordonnée temporelle au preterit (passé réel) bien gérée. « Si nous avions sensibilisé plus tôt, nous aurions changé les mauvaises habitudes avant qu'elles ne deviennent routinières. »
    \end{itemize}
    \item \textbf{Cohésion globale :} Parallélisme des structures (\eng{If we..., If the administration..., If we...}), ancrage camerounais (\eng{Mount Cameroon, school compound, canteen}), progression logique : actions immédiates $\rightarrow$ rêves d'avenir $\rightarrow$ leçons du passé.
\end{itemize}
\end{exoresolu}

\courssub{Bilan de l'activité}

Cette activité d'intégration a permis de vérifier l'acquisition opérationnelle des trois conditionnels dans un contexte de \textbf{production écrite authentique} et \textbf{collaborative}.

\begin{itemize}
    \item \textbf{Sur le plan grammatical :} La contrainte de produire exactement trois phrases par type a forcé les élèves à discriminer consciemment la valeur modale de leur énoncé (promesse vs rêve vs regret) avant de choisir la structure syntaxique. L'auto-correction par rôles (vérificateurs) a favorisé la métalinguistique : les élèves ont dû expliciter \textit{why} \eng{had started} et non \eng{started} ou \eng{would start}.
    \item \textbf{Sur le plan lexical :} L'ancrage dans le contexte camerounais (\eng{school compound, Mount Cameroon, canteen, bendskins} — bien que non utilisé ici, il était disponible) a donné du sens à l'apprentissage. L'usage de collocations (\eng{raise awareness, carbon footprint, single-use plastic}) montre une progression vers le niveau B2 visé en Terminale.
    \item \textbf{Sur le plan pragmatique :} La simulation d'assemblée (\eng{Town Hall Meeting}) a travaillé l'expression orale en continu (lecture expressive) et l'argumentation (défense de la charte). Le vote par les pairs instaure une émulation positive et valide la dimension communicationnelle de la langue.
    \item \textbf{Pistes de remédiation (si difficultés observées) :}
    \begin{itemize}
        \item Si confusion persistante 2nd/3rd conditional : réviser la frise temporelle (ligne du temps) avec marqueurs \eng{Now / Past}.
        \item Si erreurs sur participes passés irréguliers : fiche de révision ciblée \eng{Top 20 Irregular Verbs for Conditionals} (\eng{begin/begun, do/done, go/gone, see/seen, write/written, take/taken, make/made}).
        \item Si phrases trop simples : exiger l'ajout d'un complément de circonstance (lieu, temps, manière) pour la prochaine production.
    \end{itemize}
\end{itemize}

La charte gagnante, affichée dans le hall, devient un \textbf{texte ressource} vivant pour l'établissement : les élèves de classes inférieures pourront l'analyser plus tard comme modèle de conditionnels, bouclant ainsi la boucle didactique entre production senior et observation junior.

\newpage

\coursec{Méthodes et exercices résolus}

\courssub{Observation et découverte : Un texte authentique}
\begin{textebox}[Blog post : “Green Cameroon, Healthy Citizens”]
\eng{Protecting our environment is not just a slogan; it is a survival strategy for Cameroon. \textbf{If the government enforces} the new waste management law strictly, plastic pollution in the Wouri River \textbf{will drop} significantly within two years. This is a realistic goal. 

However, many citizens behave as if the problem were someone else's concern. \textbf{If people were more aware} of the link between blocked drains and malaria outbreaks, they \textbf{would stop} dumping refuse in gutters. Unfortunately, this awareness is still low. 

Looking back, the cost of inaction is heavy. \textbf{If the authorities had invested} in recycling plants a decade ago, the mountains of waste in Nkolfoulou and Makepe \textbf{would not have become} such environmental hazards today.}

\textbf{Glossaire :} \eng{enforce} (v., appliquer une loi), \eng{waste management} (n., gestion des déchets), \eng{drop} (v., baisser/diminuer), \eng{dump} (v., déverser/jeter sauvagement), \eng{refuse} (n., ordures/déchets solides), \eng{gutter} (n., caniveau), \eng{outbreak} (n., épidémie/soudaine apparition), \eng{recycling plant} (n., usine de recyclage), \eng{hazard} (n., danger/risque).
\end{textebox}

\courssub{Synthèse des trois conditionnels (Tableau récapitulatif)}
\begin{propriete}[Structure, emploi et traduction des conditionnels]
\begin{center}
\begin{tabularx}{\linewidth}{|>{\centering\arraybackslash}X|>{\centering\arraybackslash}X|>{\centering\arraybackslash}X|>{\centering\arraybackslash}X|X|}
\hline
\rowcolor{popblueL} 
\textbf{Type} & \textbf{If-clause (Condition)} & \textbf{Main Clause (Conséquence)} & \textbf{Emploi / Nuance} & \textbf{Exemple + Traduction} \\ \hline
\textbf{1st Conditional} 
(Real Future) & \eng{If + Present Simple} 
(\eng{If it rains}, \eng{If she comes}) & \eng{will / can / may / must / shall + BV} 
(\eng{will stay}, \eng{can go}) & \textbf{Hypothèse réelle, possible} dans le futur. 
Prédiction, promesse, avertissement, offre. & \eng{If you boil the water, it will be safe.} 
« Si tu fais bouillir l'eau, elle sera potable. » \\ \hline
\textbf{2nd Conditional} 
(Unreal Present) & \eng{If + Past Simple} 
(\eng{If I knew}, \eng{If he were}) 
\cle{Note : \eng{were} $\forall$ personnes} & \eng{would / could / might + BV} 
(\eng{would tell}, \eng{could buy}) & \textbf{Hypothèse imaginaire, improbable} au présent/futur. 
Conseil (\eng{If I were you...}), rêve. & \eng{If I were rich, I would travel.} 
« Si j'étais riche, je voyagerais. » \\ \hline
\textbf{3rd Conditional} 
(Unreal Past) & \eng{If + Past Perfect} 
(\eng{If I had known}, \eng{If they had come}) & \eng{would / could / might + have + V3} 
(\eng{would have told}, \eng{could have come}) & \textbf{Regret, reproche, critique} sur le passé. 
Action non réalisée (contraire aux faits). & \eng{If you had studied, you would have passed.} 
« Si tu avais étudié, tu aurais réussi. » \\ \hline
\end{tabularx}
\end{center}
\cle{Règle d'or :} \textbf{Jamais de \eng{will} ou \eng{would} dans la \eng{if-clause}} (sauf \eng{will} = accepter de/politesse : \eng{If you will wait...}).
\end{propriete}

\courssub{Méthodes et exercices résolus}

\courssub{Identifier le type de conditionnel (Observation et analyse)}
\begin{methode}[Identifier le type de conditionnel]
\textbf{Objectif :} Déterminer s'il s'agit d'un 1st, 2nd ou 3rd conditional pour appliquer la bonne structure.
\begin{enumerate}
    \item \textbf{Repérer les deux propositions :} La \cle{if-clause} (proposition conditionnelle, introduite par \eng{if} ou \eng{unless}) et la \cle{main clause} (proposition principale).
    \item \textbf{Analyser le temps du verbe dans la \eng{if-clause} :}
    \begin{itemize}
        \item \textbf{Present Simple} ($\rightarrow$ \textbf{1st Conditional}) : Hypothèse \emph{réelle}, possible dans le futur.
        \item \textbf{Past Simple} (ou \eng{were} pour tous les sujets) ($\rightarrow$ \textbf{2nd Conditional}) : Hypothèse \emph{irréelle}, peu probable ou imaginaire au présent/futur.
        \item \textbf{Past Perfect} (\eng{had + V3}) ($\rightarrow$ \textbf{3rd Conditional}) : Hypothèse \emph{passée}, irréelle (regret, reproche), contraire aux faits réels.
    \end{itemize}
    \item \textbf{Vérifier la cohérence dans la \eng{main clause} :}
    \begin{itemize}
        \item 1st : \eng{will / can / may / must + BV} (ou impératif).
        \item 2nd : \eng{would / could / might + BV}.
        \item 3rd : \eng{would / could / might + have + V3}.
    \end{itemize}
    \item \textbf{Attention aux marqueurs temporels :} \eng{tomorrow, next week, soon} (1st) ; \eng{now, today, if I were you} (2nd) ; \eng{yesterday, last year, in 2010, a decade ago} (3rd).
\end{enumerate}
\end{methode}

\begin{exoresolu}[Analyse de phrases conditionnelles]
\textbf{Énoncé :} Pour chacune des phrases suivantes, a) soulignez la \eng{if-clause}, b) entourez le verbe de la \eng{main clause}, c) identifiez le type de conditionnel (1st, 2nd ou 3rd) et d) justifiez votre choix en français.
\begin{enumerate}
    \item If the Ministry of Environment \textbf{enforces} the new law, pollution \textbf{will decrease} in Douala.
    \item If I \textbf{were} the Minister of Health, I \textbf{would build} more hospitals in rural areas.
    \item If the farmers \textbf{had used} organic fertilizer, the soil \textbf{would not have been} degraded.
    \item Unless it \textbf{rains} heavily tonight, the match \textbf{will take place} in Buea.
    \item If she \textbf{had known} about the cholera outbreak, she \textbf{would have boiled} the water.
\end{enumerate}
\tcblower
\textbf{Solution détaillée :}
\begin{enumerate}
    \item \textbf{If-clause :} \eng{If the Ministry of Environment enforces the new law}. \textbf{Main clause verb :} \eng{will decrease}. \textbf{Type : 1st Conditional}. \textbf{Justification :} Verbe \eng{enforces} au \emph{Present Simple} $\rightarrow$ hypothèse réelle et possible pour le futur. Conséquence : \eng{will + BV}.
    \item \textbf{If-clause :} \eng{If I were the Minister of Health}. \textbf{Main clause verb :} \eng{would build}. \textbf{Type : 2nd Conditional}. \textbf{Justification :} Verbe \eng{were} (subjonctif/prétérit modal) $\rightarrow$ hypothèse imaginaire au présent (je ne suis pas ministre). Conséquence : \eng{would + BV}. \cle{Note} : \eng{were} pour toutes les personnes (I/he/she/it).
    \item \textbf{If-clause :} \eng{If the farmers had used organic fertilizer}. \textbf{Main clause verb :} \eng{would not have been}. \textbf{Type : 3rd Conditional}. \textbf{Justification :} Verbe \eng{had used} au \emph{Past Perfect} $\rightarrow$ action passée non réalisée (les fermiers n'ont pas utilisé d'engrais bio). Conséquence passée : \eng{would + have + V3}.
    \item \textbf{If-clause :} \eng{Unless it rains heavily tonight} (\eng{Unless} = \eng{If not}). \textbf{Main clause verb :} \eng{will take place}. \textbf{Type : 1st Conditional}. \textbf{Justification :} \eng{rains} au \emph{Present Simple} après \eng{unless} $\rightarrow$ condition réelle future. \eng{will + BV} dans la principale.
    \item \textbf{If-clause :} \eng{If she had known about the cholera outbreak}. \textbf{Main clause verb :} \eng{would have boiled}. \textbf{Type : 3rd Conditional}. \textbf{Justification :} \eng{had known} (Past Perfect) $\rightarrow$ regret sur le passé (elle n'a pas su). \eng{would have boiled} (Perfect Conditional) $\rightarrow$ action passée non réalisée.
\end{enumerate}
\textbf{Conclusion :} L'identification repose \emph{exclusivement} sur le temps de la \eng{if-clause}. Ne vous laissez pas distraire par le vocabulaire.
\end{exoresolu}

\courssub{Construire et utiliser le 1st Conditional (Hypothèse réelle)}
\begin{methode}[Construire le 1st Conditional]
\textbf{Structure :} \eng{If + Present Simple, ... will/can/may/must + Base Verbale} (ou Impératif).
\textbf{Emploi :} Prédiction, avertissement, promesse, offre, menace basés sur une condition \textbf{réalisable}.
\begin{enumerate}
    \item \textbf{Poser la condition au présent simple :} Sujet + V (+\eng{s} à la 3e pers. sg). \emph{Ex: If the government bans plastic bags...}
    \item \textbf{Choisir le modal adapté à la nuance :}
    \begin{itemize}
        \item \eng{will} : certitude, prédiction neutre.
        \item \eng{can} : possibilité, capacité.
        \item \eng{may/might} : probabilité plus faible.
        \item \eng{must/should} : obligation, conseil fort.
        \item \textbf{Impératif} : ordre, conseil direct (\eng{If you feel sick, see a doctor.}).
    \end{itemize}
    \item \textbf{Ordre des propositions :} \eng{If-clause} en premier (virgule) ou \eng{Main clause} en premier (pas de virgule). \eng{We will plant trees if the rain starts.}
    \item \textbf{Négation :} \eng{If + don't/doesn't + BV} / \eng{won't/can't + BV}. \emph{Ex: If we don't protect the forest, wildlife will disappear.}
\end{enumerate}
\end{methode}

\begin{exoresolu}[Mise en situation : Prévisions environnementales]
\textbf{Énoncé :} Complétez les phrases en conjuguant les verbes entre parenthèses au 1st Conditional. Traduisez ensuite chaque phrase.
\begin{enumerate}
    \item If the temperature \_\_\_\_\_\_ (rise) by 2°C, sea levels \_\_\_\_\_\_ (threaten) coastal cities like Kribi.
    \item The NGO \_\_\_\_\_\_ (launch) a cleaning campaign if volunteers \_\_\_\_\_\_ (sign up) this week.
    \item \_\_\_\_\_\_ (not / throw) your waste in the river if you \_\_\_\_\_\_ (want) to preserve the fish.
    \item If the hospital \_\_\_\_\_\_ (receive) the new vaccines, they \_\_\_\_\_\_ (can / start) the immunization next Monday.
\end{enumerate}
\tcblower
\textbf{Solution détaillée :}
\begin{enumerate}
    \item \textbf{rises / will threaten}. \eng{rise} $\rightarrow$ 3e pers. sg \eng{rises}. \eng{will threaten} (prédiction). \textit{Traduction : « Si la température augmente de 2°C, le niveau de la mer menacera les villes côtières comme Kribi. »}
    \item \textbf{will launch / sign up}. Principale d'abord : pas de virgule. \eng{sign up} (présent simple). \textit{Traduction : « L'ONG lancera une campagne de nettoyage si des bénévoles s'inscrivent cette semaine. »}
    \item \textbf{Do not throw / want}. \textbf{Impératif négatif} dans la principale (conseil/ordre) + \eng{if-clause} au présent. \textit{Traduction : « Ne jetez pas vos déchets dans la rivière si vous voulez préserver les poissons. »}
    \item \textbf{receives / can start}. \eng{receives} (3e pers. sg). \eng{can start} (capacité/possibilité future). \textit{Traduction : « Si l'hôpital reçoit les nouveaux vaccins, ils pourront commencer la vaccination lundi prochain. »}
\end{enumerate}
\textbf{Point de vigilance :} À la phrase 3, l'impératif remplace \eng{will}. C'est une structure courante pour les conseils de santé/sécurité.
\end{exoresolu}

\courssub{Construire et utiliser le 2nd Conditional (Hypothèse irréelle présente)}
\begin{methode}[Construire le 2nd Conditional]
\textbf{Structure :} \eng{If + Past Simple (V-ed / V2), ... would/could/might + Base Verbale}.
\textbf{Emploi :} Situation \textbf{imaginaire, hypothétique, improbable} au présent ou futur ; conseils (\eng{If I were you...}).
\begin{enumerate}
    \item \textbf{If-clause :} Mettez le verbe au \textbf{Prétérit (Past Simple)}. Pour \eng{be} $\rightarrow$ \textbf{\eng{were}} \emph{pour toutes les personnes} (formel/standard) ; \eng{was} possible à l'oral pour I/he/she/it (moins formel). \emph{Ex: If I \textbf{were} rich...}
    \item \textbf{Main clause :} \textbf{\eng{would} (+ \eng{not}) + BV} (conséquence imaginaire). \eng{could} = capacité/possibilité ; \eng{might} = possibilité faible.
    \item \textbf{Contracter :} \eng{I'd, you'd, he'd, we'd, they'd} (à l'oral et écrit informel). \eng{Wouldn't} pour la négation.
    \item \textbf{Traduction française :} Utilisez l'\textbf{imparfait} dans le \eng{si}-clause et le \textbf{conditionnel présent} dans la principale. \eng{If I knew, I would tell you.} = « Si je \emph{savais}, je te le \emph{dirais}. »
    \item \textbf{Conseil type :} \eng{If I were you, I would + BV} $\rightarrow$ « À ta place, je... »
\end{enumerate}
\end{methode}

\begin{exoresolu}[Conseils imaginaires : Bien-être et santé]
\textbf{Énoncé :} Réécrivez les phrases en utilisant le 2nd Conditional pour donner un conseil ou exprimer une situation imaginaire. Commencez par \eng{If I were you...} ou le \eng{if} indiqué.
\begin{enumerate}
    \item You eat too much \eng{suya} (grilled meat) late at night. You have stomach aches. (\eng{If you...})
    \item She doesn't sleep enough. She is always tired. (\eng{If she...})
    \item We don't have a mosquito net. We get malaria often. (\eng{If we...})
    \item I am not a doctor. I cannot prescribe medicine. (\eng{If I...})
\end{enumerate}
\tcblower
\textbf{Solution détaillée :}
\begin{enumerate}
    \item \eng{If you \textbf{ate} less \emph{suya} late at night, you \textbf{would not have} stomach aches.} (Ou \eng{If I were you, I would not eat...}). \textit{Passage du présent réel (eat/have) à l'irréel (ate/would not have).}
    \item \eng{If she \textbf{slept} more, she \textbf{would not be} always tired.} (\eng{sleep} $\rightarrow$ \eng{slept} ; \eng{is} $\rightarrow$ \eng{would not be}). \textit{Traduction : « Si elle dormait plus, elle ne serait pas toujours fatiguée. »}
    \item \eng{If we \textbf{had} a mosquito net, we \textbf{would not get} malaria often.} (\eng{have} $\rightarrow$ \eng{had} ; \eng{get} $\rightarrow$ \eng{would not get}). \cle{Attention} : \eng{have} (possession) $\rightarrow$ prétérit \eng{had}, pas \eng{had got}.
    \item \eng{If I \textbf{were} a doctor, I \textbf{could prescribe} medicine.} (\eng{am} $\rightarrow$ \eng{were} pour toutes les personnes ; \eng{can} $\rightarrow$ \eng{could}). \textit{Traduction : « Si j'étais médecin, je pourrais prescrire des médicaments. »}
\end{enumerate}
\textbf{Conclusion :} Le 2nd conditional exprime un \emph{écart avec la réalité actuelle}. Les verbes d'état (\eng{be, have, know}) prennent aussi le prétérit.
\end{exoresolu}

\courssub{Construire et utiliser le 3rd Conditional (Regret / Hypothèse passée)}
\begin{methode}[Construire le 3rd Conditional]
\textbf{Structure :} \eng{If + Past Perfect (had + V3), ... would/could/might + have + V3 (Participe Passé)}.
\textbf{Emploi :} \textbf{Regret, reproche, critique} sur une action passée qui \textbf{n'a pas eu lieu} (ou a eu lieu contrairement à la condition). « Ce qui se serait passé si... »
\begin{enumerate}
    \item \textbf{If-clause :} \textbf{\eng{Had + Participe Passé (V3)}}. Négation : \eng{Had not (Hadn't) + V3}. \emph{Ex: If they \textbf{had planted} trees...}
    \item \textbf{Main clause :} \textbf{\eng{Would / Could / Might + Have + V3}}. Négation : \eng{Would not (Wouldn't) + Have + V3}.
    \item \textbf{Contracter (oral/écrit courant) :} \eng{If I'd known... I'd have come.} (\eng{I'd} = \eng{I had} dans le \eng{if}, \eng{I'd} = \eng{I would} dans la principale).
    \item \textbf{Inversion formelle (style soutenu/littéraire) :} \eng{Had I known, I would have acted.} (Suppression de \eng{if}).
    \item \textbf{Traduction française :} \textbf{Plus-que-parfait} dans le \eng{si}-clause + \textbf{Conditionnel Passé} (\eng{aurais + participe}) dans la principale. \eng{If you had listened, you would have understood.} = « Si tu \emph{avais écouté}, tu \emph{aurais compris}. »
\end{enumerate}
\end{methode}

\begin{exoresolu}[Regrets sur la gestion des déchets à Douala]
\textbf{Énoncé :} Lisez les faits réels (réalité passée). Exprimez le regret ou la conséquence imaginaire en utilisant le 3rd Conditional.
\begin{enumerate}
    \item \textbf{Fait :} The city council \textbf{did not collect} the garbage on time. The drains \textbf{got blocked}.
    \item \textbf{Fait :} Residents \textbf{burned} plastic waste. The air \textbf{became} toxic.
    \item \textbf{Fait :} The government \textbf{did not invest} in recycling plants. The environment \textbf{suffered}.
    \item \textbf{Fait :} I \textbf{forgot} my reusable bag. I \textbf{had to buy} a plastic one at the market.
\end{enumerate}
\tcblower
\textbf{Solution détaillée :}
\begin{enumerate}
    \item \eng{If the city council \textbf{had collected} the garbage on time, the drains \textbf{would not have got(ten)} blocked.}
    \begin{itemize}
        \item \eng{did not collect} (négation réel) $\rightarrow$ \eng{had collected} (affirmatif conditionnel).
        \item \eng{got blocked} (réel) $\rightarrow$ \eng{would not have got blocked} (négation conditionnel). \textit{Participe de \eng{get} : \eng{got} (UK) ou \eng{gotten} (US).}
    \end{itemize}
    \item \eng{If residents \textbf{had not burned} plastic waste, the air \textbf{would not have become} toxic.}
    \begin{itemize}
        \item \eng{burned} (affirmatif réel) $\rightarrow$ \eng{had not burned}.
        \item \eng{became} $\rightarrow$ \eng{would not have become}.
    \end{itemize}
    \item \eng{If the government \textbf{had invested} in recycling plants, the environment \textbf{would not have suffered}.}
    \begin{itemize}
        \item \eng{did not invest} $\rightarrow$ \eng{had invested}.
        \item \eng{suffered} $\rightarrow$ \eng{would not have suffered}.
    \end{itemize}
    \item \eng{If I \textbf{had not forgotten} my reusable bag, I \textbf{would not have had to buy} a plastic one at the market.}
    \begin{itemize}
        \item \eng{forgot} $\rightarrow$ \eng{had not forgotten}.
        \item \eng{had to buy} (semi-modal) $\rightarrow$ \eng{would not have had to buy}.
    \end{itemize}
\end{enumerate}
\textbf{Astuce Bac :} Pour passer du réel au 3rd conditional : \textbf{Négation $\leftrightarrow$ Affirmation}. Si le réel est négatif (\eng{did not}), le conditionnel est affirmatif (\eng{had + V3}), et inversement.
\end{exoresolu}

\courssub{Transformer des phrases : du réel à l'irréel (Entraînement mixte)}
\begin{methode}[Transformer une phrase d'un conditionnel à l'autre]
\textbf{Objectif :} Passer du 1st au 2nd (réel $\rightarrow$ imaginaire présent), du 2nd au 3rd (présent $\rightarrow$ regret passé), ou inversement.
\begin{enumerate}
    \item \textbf{Identifier le temps de départ} (Présent / Prétérit / Past Perfect) et le modal (\eng{will/would/would have}).
    \item \textbf{Appliquer le décalage temporel (Backshifting) :}
    \begin{center}
    \begin{tabular}{|c|c|c|}\hline
    \rowcolor{popblueL} \textbf{Depuis / Vers} & \textbf{If-clause} & \textbf{Main clause} \\ \hline
    1st $\rightarrow$ 2nd & Present Simple $\rightarrow$ \textbf{Past Simple} & \eng{will + BV} $\rightarrow$ \textbf{\eng{would + BV}} \\ \hline
    2nd $\rightarrow$ 3rd & Past Simple $\rightarrow$ \textbf{Past Perfect (\eng{had + V3})} & \eng{would + BV} $\rightarrow$ \textbf{\eng{would + have + V3}} \\ \hline
    1st $\rightarrow$ 3rd & Present Simple $\rightarrow$ \textbf{Past Perfect} & \eng{will + BV} $\rightarrow$ \textbf{\eng{would + have + V3}} \\ \hline
    \end{tabular}
    \end{center}
    \item \textbf{Adapter le sens :} 1st = \eng{It will rain} ; 2nd = \eng{If it rained (but it doesn't)} ; 3rd = \eng{If it had rained (but it didn't)}.
    \item \textbf{Vérifier l'accord du sujet} (3e pers. sg au présent, \eng{were} au prétérit pour \eng{be}).
    \item \textbf{Garder le même sujet et le même verbe lexical} (sauf changement de sens imposé).
\end{enumerate}
\end{methode}

\begin{exoresolu}[Exercices de transformation type Bac]
\textbf{Énoncé :} Transformez les phrases suivantes selon l'indication entre parenthèses.
\begin{enumerate}
    \item If the NGO gets funding, they will drill a well in the village. (\textbf{2nd Conditional})
    \item If I were the Principal, I would ban single-use plastics in the canteen. (\textbf{3rd Conditional} — contexte : \eng{Last year, I was not the Principal.})
    \item If you boil the water, it will be safe to drink. (\textbf{3rd Conditional} — contexte : \eng{Yesterday, you didn't boil it.})
    \item Unless the factory stops polluting the river, the fish will die. (\textbf{2nd Conditional} avec \eng{if})
\end{enumerate}
\tcblower
\textbf{Solution détaillée :}
\begin{enumerate}
    \item \eng{If the NGO \textbf{got} funding, they \textbf{would drill} a well in the village.}
    \begin{itemize}
        \item \eng{gets} (Present) $\rightarrow$ \eng{got} (Past).
        \item \eng{will drill} $\rightarrow$ \eng{would drill}.
        \item \textit{Sens :} L'ONG n'a pas le financement actuellement (hypothèse imaginaire).
    \end{itemize}
    \item \eng{If I \textbf{had been} the Principal \textbf{last year}, I \textbf{would have banned} single-use plastics in the canteen.}
    \begin{itemize}
        \item \eng{were} (Past) $\rightarrow$ \eng{had been} (Past Perfect). Ajout du marqueur temporel \eng{last year} pour ancrer dans le passé.
        \item \eng{would ban} $\rightarrow$ \eng{would have banned}.
        \item \textit{Sens :} Regret sur le passé (je n'étais pas principal).
    \end{itemize}
    \item \eng{If you \textbf{had boiled} the water, it \textbf{would have been} safe to drink.}
    \begin{itemize}
        \item \eng{boil} (Present) $\rightarrow$ \eng{had boiled} (Past Perfect). Contexte \eng{Yesterday}.
        \item \eng{will be} $\rightarrow$ \eng{would have been}.
    \end{itemize}
    \item \eng{If the factory \textbf{stopped} polluting the river, the fish \textbf{would not die}.}
    \begin{itemize}
        \item \eng{Unless} $\rightarrow$ \eng{If ... not} / \eng{If ... stopped}. \eng{stops} (Present) $\rightarrow$ \eng{stopped} (Past).
        \item \eng{will die} $\rightarrow$ \eng{would not die} (pas \eng{would die not}).
        \item \textit{Note :} \eng{Unless} disparaît au profit de \eng{If + négation} ou verbe d'action négative (\eng{stopped}).
    \end{itemize}
\end{enumerate}
\textbf{Conclusion :} La transformation mécanique (backshifting) doit s'accompagner d'un ajustement du \textbf{sens} (marqueurs de temps, contexte).
\end{exoresolu}

\courssub{Rédiger un paragraphe argumentatif avec les conditionnels (Production écrite)}
\begin{methode}[Rédiger un court texte intégrant les trois conditionnels]
\textbf{Objectif :} Produire un paragraphe cohérent (80-100 mots) sur l'environnement/la santé en utilisant \textbf{au moins une fois} chaque conditionnel.
\begin{enumerate}
    \item \textbf{Planifier la structure logique :}
    \begin{itemize}
        \item \textbf{1st Conditional :} \textbf{Avenir réel / Appel à l'action.} (\eng{If we act now, we will save...})
        \item \textbf{2nd Conditional :} \textbf{Scénario idéal / Conseil / Critique générale.} (\eng{If people were more careful... / If I were the Minister...})
        \item \textbf{3rd Conditional :} \textbf{Regret sur le passé / Leçon de l'histoire.} (\eng{If we had protected the forest earlier...})
    \end{itemize}
    \item \textbf{Rédiger les phrases-clés (Draft) :} Écrivez d'abord les 3 phrases "squelettes" avec les structures cibles. Vérifiez la grammaire (\eng{had + V3}, \eng{were}, \eng{will/would/would have}).
    \item \textbf{Relier par des connecteurs logiques :} \eng{However, Moreover, Consequently, For instance,} pour fluidifier le paragraphe.
    \item \textbf{Enrichir le vocabulaire thématique :} \eng{deforestation, waste management, awareness, sustainable, outbreak, resilience, stakeholder}.
    \item \textbf{Relire (Check-list) :} 1. Les 3 types présents ? 2. Pas de \eng{will} dans le \eng{if} ? 3. \eng{Were} pour \eng{be} au 2nd ? 4. \eng{Had + V3} au 3rd ? 5. Orthographe vocabulaire ?
\end{enumerate}
\end{methode}

\begin{exoresolu}[Sujet de composition : « Environment and Health »]
\textbf{Énoncé :} \eng{Write a paragraph of about 90 words on the following topic: "Protecting our environment is protecting our health." Use the 1st, 2nd and 3rd conditionals at least once each.}
\tcblower
\textbf{Solution détaillée (Modèle de rédaction) :}
\eng{Protecting our environment is undoubtedly protecting our health. \textbf{If the authorities enforce} stricter regulations on industrial waste, the quality of water in the Wouri River \textbf{will improve} significantly (1st Conditional). \textbf{If I were} the Minister of Environment, I \textbf{would launch} a nationwide campaign to educate citizens on plastic recycling (2nd Conditional). \textbf{If we had invested} in sustainable agriculture ten years ago, the soil erosion in the West Region \textbf{would not have reached} such a critical level (3rd Conditional). Ultimately, every citizen must act responsibly today to secure a healthier tomorrow.}

\textbf{Analyse grammaticale du modèle :}
\begin{itemize}
    \item \textbf{Phrase 1 (1st) :} \eng{enforce} (Present), \eng{will improve} (Future). Contexte : mesure concrète possible.
    \item \textbf{Phrase 2 (2nd) :} \eng{were} (Subjonctif/Prétérit \eng{be}), \eng{would launch}. Contexte : hypothèse imaginaire (je ne suis pas ministre).
    \item \textbf{Phrase 3 (3rd) :} \eng{had invested} (Past Perfect), \eng{would not have reached} (Conditionnel Passé). Contexte : regret passé (action non faite il y a 10 ans).
    \item \textbf{Vocabulaire :} \eng{enforce, regulations, industrial waste, nationwide campaign, sustainable agriculture, soil erosion, critical level}.
\end{itemize}
\textbf{Note :} Ce paragraphe obtient la note maximale car il respecte la contrainte de forme (3 conditionnels), la cohérence thématique et la correction linguistique.
\end{exoresolu}

\begin{attention}[Les 5 erreurs qui coûtent des points au Bac]
\begin{enumerate}
    \item \textbf{\eng{Will / Would} dans la \eng{if-clause} (Calque du français « Si je \emph{ferais}... »)} \\
    \eng{\textcolor{red}{If it will rain, I stay home.}} $\rightarrow$ \eng{\textcolor{popgreen}{If it rains, I will stay home.}} (1st) \\
    \eng{\textcolor{red}{If I would know, I would tell you.}} $\rightarrow$ \eng{\textcolor{popgreen}{If I knew, I would tell you.}} (2nd) \\
    \eng{\textcolor{red}{If you would have come, we would have eaten.}} $\rightarrow$ \eng{\textcolor{popgreen}{If you had come, we would have eaten.}} (3rd) \\
    \cle{Règle d'or :} \textbf{Jamais de \eng{will/would} après \eng{if} (sauf politesse : \eng{If you will wait a moment...}).}
    \item \textbf{Confusion 2nd / 3rd Conditional (Oubli du \eng{have} au 3rd)} \\
    \eng{\textcolor{red}{If I had studied, I would pass the exam.}} (Mélange : \eng{if} passé, principale présent) \\
    $\rightarrow$ \eng{\textcolor{popgreen}{If I had studied, I would \textbf{have passed} the exam.}} (3rd) \\
    Ou : \eng{\textcolor{popgreen}{If I studied, I would pass.}} (2nd). \textbf{Vérifiez la cohérence temporelle.}
    \item \textbf{\eng{Was} au lieu de \eng{Were} au 2nd Conditional (Style formel exigé à l'écrit)} \\
    \eng{\textcolor{red}{If I was you, I would go.}} $\rightarrow$ \eng{\textcolor{popgreen}{If I \textbf{were} you, I would go.}} \\
    \textit{Au Bac, l'écrit formel impose \eng{were} pour toutes les personnes (I, he, she, it).}
    \item \textbf{Mauvaise formation du Past Perfect (\eng{Had + V2} au lieu de \eng{Had + V3})} \\
    \eng{\textcolor{red}{If he had went / had did / had saw...}} $\rightarrow$ \eng{\textcolor{popgreen}{If he had \textbf{gone / done / seen}...}} \\
    \cle{Apprenez par cœur la liste des verbes irréguliers (V2 $\neq$ V3).}
    \item \textbf{Traduction littérale de « Si j'aurais su » $\rightarrow$ \eng{If I would have known}} \\
    Le français utilise le conditionnel dans les deux propositions (« Si j'\emph{aurais} su, j'\emph{aurais} venu »). L'anglais \textbf{ne le fait jamais}. \\
    \textbf{Français :} Si + Imparfait / Plus-que-parfait $\rightarrow$ Conditionnel Présent / Passé. \\
    \textbf{Anglais :} If + Past / Past Perfect $\rightarrow$ Would / Would Have. \\
    \textbf{Mémo :} \eng{If I \textbf{knew} (Imparfait), I \textbf{would know} (Cond. Présent).} \quad \eng{If I \textbf{had known} (PQP), I \textbf{would have known} (Cond. Passé).}
\end{enumerate}
\end{attention}

\courssub{Entraînement personnel (Exercices non résolus)}
\begin{exercice}[Entraînement mixte : Conditionnels 1, 2, 3]
\textbf{A. Mettez les verbes entre parenthèses au temps qui convient (1st, 2nd ou 3rd Conditional).}
\begin{enumerate}
    \item If the rain \_\_\_\_\_\_ (stop) soon, we \_\_\_\_\_\_ (play) the final match in Bamenda.
    \item If I \_\_\_\_\_\_ (be) you, I \_\_\_\_\_\_ (not / trust) that herbalist without a licence.
    \item The cholera epidemic \_\_\_\_\_\_ (not / spread) so fast if the population \_\_\_\_\_\_ (have) access to clean water last year.
    \item Unless the council \_\_\_\_\_\_ (collect) the bins tomorrow, the market \_\_\_\_\_\_ (smell) terrible.
    \item If the researcher \_\_\_\_\_\_ (publish) her findings earlier, the government \_\_\_\_\_\_ (can / react) faster to the outbreak.
\end{enumerate}

\textbf{B. Réécrivez les phrases en commençant par les mots donnés (Transformation).}
\begin{enumerate}
    \item The farmers didn't use compost, so the yield was low. \\
    \eng{If the farmers \_\_\_\_\_\_ \_\_\_\_\_\_ \_\_\_\_\_\_ \_\_\_\_\_\_, the yield \_\_\_\_\_\_ \_\_\_\_\_\_ \_\_\_\_\_\_ \_\_\_\_\_\_.}
    \item We don't have enough incubators, so premature babies are at risk. \\
    \eng{If we \_\_\_\_\_\_ \_\_\_\_\_\_ \_\_\_\_\_\_ \_\_\_\_\_\_, premature babies \_\_\_\_\_\_ \_\_\_\_\_\_ \_\_\_\_\_\_ \_\_\_\_\_\_.}
    \item "You should see a doctor immediately," the nurse told me. \\
    \eng{If I \_\_\_\_\_\_ \_\_\_\_\_\_ \_\_\_\_\_\_, I \_\_\_\_\_\_ \_\_\_\_\_\_ \_\_\_\_\_\_ \_\_\_\_\_\_.}
\end{enumerate}

\textbf{C. Expression écrite (30-40 mots).}
\eng{Complete the paragraph: "Yesterday, a bendskin rider had an accident because he wasn't wearing a helmet. If he \_\_\_\_\_\_ \_\_\_\_\_\_ \_\_\_\_\_\_ (wear) one, he \_\_\_\_\_\_ \_\_\_\_\_\_ \_\_\_\_\_\_ \_\_\_\_\_\_ (not / injure) his head. From now on, \_\_\_\_\_\_ \_\_\_\_\_\_ \_\_\_\_\_\_ (1st conditional advice)."}
\end{exercice}

\begin{corrige}[Exercice Entraînement personnel]
\textbf{A.}
\begin{enumerate}
    \item \textbf{stops / will play} (1st : condition réelle future).
    \item \textbf{were / would not trust} (2nd : conseil imaginaire, \eng{were} pour \eng{I}).
    \item \textbf{would not have spread / had had} (3rd : regret passé, \eng{had had} = Past Perfect de \eng{have}).
    \item \textbf{does not collect / will smell} (1st : \eng{Unless} + Present Simple $\rightarrow$ \eng{will}).
    \item \textbf{had published / could have reacted} (3rd : regret passé, \eng{could have + V3}).
\end{enumerate}

\textbf{B.}
\begin{enumerate}
    \item \eng{If the farmers \textbf{had used} compost, the yield \textbf{would have been} higher.} (3rd Conditional)
    \item \eng{If we \textbf{had} enough incubators, premature babies \textbf{would not be} at risk.} (2nd Conditional : situation présente imaginaire).
    \item \eng{If I \textbf{were} you, I \textbf{would see} a doctor immediately.} (2nd Conditional : conseil \eng{If I were you}).
\end{enumerate}

\textbf{C.}
\eng{Yesterday, a bendskin rider had an accident because he wasn't wearing a helmet. If he \textbf{had worn} one, he \textbf{would not have injured} his head. From now on, \textbf{if you ride a motorbike, you must wear a helmet.} (ou \eng{...you will protect your life.})}
\end{corrige}

\newpage

\coursec{Exercices}

\courssub{Je vérifie mes connaissances}

\begin{exercice}[QCM : le bon conditionnel]
Entoure la bonne réponse.
\begin{enumerate}
\item If it rains tomorrow, we \textbf{...} at home. \\
a) will stay \quad b) would stay \quad c) would have stayed
\item If I were rich, I \textbf{...} a hospital in my village. \\
a) will build \quad b) would build \quad c) would have built
\item If they had planted trees, the soil \textbf{...} so eroded. \\
a) will not be \quad b) would not be \quad c) would not have been
\item If you heat water to 100°C, it \textbf{...}. \\
a) boils \quad b) would boil \quad c) would have boiled
\end{enumerate}
\end{exercice}

\begin{exercice}[Vrai ou faux — justifie]
Dis si chaque phrase est vraie (V) ou fausse (F) et justifie ta réponse en une phrase en français.
\begin{enumerate}
\item In the first conditional, we use \eng{will} in the if-clause.
\item \eng{If I were you} is correct, even with the subject \eng{I}.
\item The third conditional expresses a real possibility in the future.
\item \eng{If she had studied, she would have passed} expresses regret about the past.
\item In the second conditional, the if-clause is in the past simple.
\end{enumerate}
\end{exercice}

\begin{exercice}[Vocabulaire : l'environnement et la santé]
Complète chaque phrase avec le mot anglais qui convient, choisi dans la liste : \eng{pollution — recycling — disease — drought — waste}.
\begin{enumerate}
\item If factories reduce their \textbf{...}, the air in Douala will be cleaner.
\item If people had sorted their \textbf{...}, the dump would not have grown so fast.
\item If there is a long \textbf{...}, farmers will lose their crops.
\item If everyone practised \textbf{...}, we would save natural resources.
\item If the water is not boiled, children can catch a \textbf{...}.
\end{enumerate}
\end{exercice}

\begin{exercice}[Conjugaison à compléter]
Mets le verbe entre parenthèses à la forme correcte.
\begin{enumerate}
\item If the government \textbf{...} (invest) in hospitals, many lives \textbf{...} (be) saved.
\item If I \textbf{...} (know) about the danger, I \textbf{...} (not/go) there yesterday.
\item If we \textbf{...} (recycle) more, the city \textbf{...} (be) cleaner.
\item If she \textbf{...} (be) a doctor, she \textbf{...} (work) in a rural area.
\end{enumerate}
\end{exercice}

\courssub{Je m'entraîne}

\begin{exercice}[Traduction]
Traduis les phrases suivantes en anglais.
\begin{enumerate}
\item Si j'étais toi, je planterais des arbres.
\item Si nous avions recyclé, la ville aurait été plus propre.
\item S'il pleut demain, le match de football sera annulé.
\item Si les gens ne jetaient pas de plastique dans les rivières, les poissons ne mourraient pas.
\end{enumerate}
\end{exercice}

\begin{exercice}[Transformation : les trois conditionnels]
Réécris la phrase \eng{If people recycle, the city will be clean} :
\begin{enumerate}
\item au deuxième conditionnel ;
\item au troisième conditionnel ;
\item puis explique en français la différence de sens entre les trois versions.
\end{enumerate}
\end{exercice}

\begin{textebox}[A cleaner Buea]
If the town council (build) \textbf{...} more rubbish bins, the streets of Buea (be) \textbf{...} much cleaner. Last year, the situation was terrible: if the inhabitants (not/burn) \textbf{...} their waste in the open air, the air (be) \textbf{...} less polluted. Now things are changing. If every family (sort) \textbf{...} its rubbish next year, the town (become) \textbf{...} a model for the whole region. A local student says: “If I (be) \textbf{...} the mayor, I (create) \textbf{...} a recycling centre in every quarter.”
\end{textebox}

\begin{exercice}[Texte à trous]
Complète le texte « A cleaner Buea » ci-dessus en mettant les verbes entre parenthèses à la forme correcte. Identifie ensuite le type de conditionnel utilisé dans chaque phrase.
\end{exercice}

\begin{exercice}[Correction d'erreurs]
Chaque phrase contient une ou plusieurs fautes. Souligne les fautes et réécris la phrase correctement.
\begin{enumerate}
\item \eng{If it will rain tomorrow, we would stayed at home.} (trois fautes)
\item \eng{If I would be you, I would see a doctor.}
\item \eng{If she had went to the hospital earlier, she would have been saved.}
\item \eng{If they will protect the forest, the climate would improve.}
\end{enumerate}
\end{exercice}

\courssub{Je me prépare au Bac}

\begin{textebox}[The cost of dirty water]
In many villages of the Far North region of Cameroon, families still drink water from open wells. Nurse Amina, who works in a small health centre, explains: “If people boiled their water, we would have far fewer cases of cholera. Last year, if the authorities had repaired the broken pump earlier, the epidemic would not have spread so fast.” The government now promises new boreholes. “If they build one borehole in every village, life will change completely,” says the village chief. “Children will stop walking five kilometres every morning, and they will arrive at school on time.” But Amina remains cautious: “If I were the minister of health, I would also train local technicians. A pump that nobody can repair is useless.”
\end{textebox}

\begin{exercice}[Type Bac — Compréhension et langue]
Lis le texte « The cost of dirty water » ci-dessus, puis réponds aux questions en anglais, avec tes propres mots.\\
\textbf{A. Comprehension questions}
\begin{enumerate}
\item According to Nurse Amina, how can villagers avoid cholera?
\item What happened last year because the pump was not repaired in time?
\item Give two benefits of the new boreholes mentioned by the village chief.
\item Why does Amina want to train local technicians?
\end{enumerate}
\textbf{B. Language work}
\begin{enumerate}
\item Find in the text one example of each conditional (1st, 2nd, 3rd) and copy it.
\item Complete: If the authorities \textbf{...} (repair) the pump earlier, the epidemic \textbf{...} (not/spread).
\item Choose the correct answer: If I had known about the pollution, I (would come / would have come / will have come) earlier.
\end{enumerate}
\end{exercice}

\begin{exercice}[Type Bac — Traduction]
Traduis le passage suivant en anglais.\\
« Si le gouvernement investissait davantage dans les hôpitaux, beaucoup de vies seraient sauvées. L'année dernière, si nous avions protégé la forêt, les inondations n'auraient pas détruit les récoltes. Demain, si chaque citoyen plante un arbre, notre pays deviendra plus vert. »
\end{exercice}

\begin{exercice}[Type Bac — Expression écrite]
Write a paragraph of about \textbf{80 words} beginning with: \eng{If everyone protected the environment, ...}\\
Consignes : utilise au moins \textbf{deux conditionnels différents} (par exemple un 1st et un 2nd conditional, ou un 2nd et un 3rd) ; fais référence à la situation de ton école ou de ta ville (déchets, eau, arbres, pollution) ; souligne les deux structures conditionnelles utilisées.
\end{exercice}

\begin{methode}[Réussir le paragraphe au Bac]
\begin{enumerate}
\item \textbf{Planifie} : note deux ou trois idées concrètes (déchets dans ta cour d'école, eau potable, arbres plantés) et choisis tes deux conditionnels avant d'écrire.
\item \textbf{Rédige} : commence par la phrase imposée au 2nd conditional, puis enchaîne avec un 1st conditional (projet futur réaliste) ou un 3rd conditional (regret sur le passé).
\item \textbf{Vérifie} : pas de \eng{will} ni de \eng{would} dans la subordonnée en \eng{if} ; \eng{were} à toutes les personnes au 2nd conditional ; \eng{had} + participe passé au 3rd conditional.
\item \textbf{Compte} les mots (environ 80) et souligne les deux structures demandées.
\end{enumerate}
\end{methode}

\begin{exemplebox}[Modèle de paragraphe]
\eng{If everyone protected the environment, our town would be a much healthier place to live in. In my school, for example, students throw plastic bags on the ground every day. If we had learnt to sort our waste earlier, the schoolyard would not have become so dirty. But it is not too late: if every class plants two trees this year, the school will be greener and shadier in a few years. If I were the headmaster, I would also organise a cleaning day every month. Small actions can change our lives.}\\
Traduction : « Si tout le monde protégeait l'environnement, notre ville serait un endroit bien plus sain où vivre. Dans mon école, par exemple, les élèves jettent des sachets en plastique par terre tous les jours. Si nous avions appris à trier nos déchets plus tôt, la cour ne serait pas devenue si sale. Mais il n'est pas trop tard : si chaque classe plante deux arbres cette année, l'école sera plus verte et plus ombragée dans quelques années. Si j'étais le proviseur, j'organiserais aussi une journée de nettoyage chaque mois. De petites actions peuvent changer nos vies. »
\end{exemplebox}

\newpage

\coursec{Corrigés des exercices}

\begin{corrige}[Exercice 1 — QCM : le bon conditionnel]
\begin{enumerate}
\item \textbf{a) will stay} — 1st conditional : \eng{if} + present simple (\eng{rains}), \eng{will} + base verbale. Hypothèse réelle sur le futur. \eng{If it rains tomorrow, we will stay at home.} « S'il pleut demain, nous resterons à la maison. »
\item \textbf{b) would build} — 2nd conditional : \eng{if} + past simple (\eng{were}), \eng{would} + base verbale. Hypothèse irréelle au présent. \eng{If I were rich, I would build a hospital in my village.} « Si j'étais riche, je construirais un hôpital dans mon village. »
\item \textbf{c) would not have been} — 3rd conditional : \eng{if} + past perfect (\eng{had planted}), \eng{would have} + participe passé. Regret sur le passé. \eng{If they had planted trees, the soil would not have been so eroded.} « S'ils avaient planté des arbres, le sol n'aurait pas été aussi érodé. »
\item \textbf{a) boils} — zero conditional : vérité scientifique générale, present simple dans les deux propositions. \eng{If you heat water to 100°C, it boils.} « Si l'on chauffe l'eau à 100°C, elle bout. »
\end{enumerate}
\end{corrige}

\begin{corrige}[Exercice 2 — Vrai ou faux — justifie]
\begin{enumerate}
\item \textbf{F} — Dans la subordonnée en \eng{if}, on emploie le present simple ; \eng{will} se place dans la proposition principale : \eng{If it rains, we will stay at home.}
\item \textbf{V} — Au 2nd conditional, on utilise \eng{were} pour toutes les personnes (forme du subjonctif) : \eng{If I were you, I would see a doctor.}
\item \textbf{F} — Le 3rd conditional exprime une hypothèse irréelle sur le \textbf{passé} (un regret), jamais une possibilité future.
\item \textbf{V} — \eng{If she had studied, she would have passed} : la condition n'a pas été remplie dans le passé, d'où le regret. « Si elle avait étudié, elle aurait réussi. »
\item \textbf{V} — Structure : \eng{if} + past simple, \eng{would} + base verbale : \eng{If I had money, I would travel.}
\end{enumerate}
\end{corrige}

\begin{corrige}[Exercice 3 — Vocabulaire : l'environnement et la santé]
\begin{enumerate}
\item \eng{pollution} (la pollution) — \eng{If factories reduce their pollution, the air in Douala will be cleaner.}
\item \eng{waste} (les déchets) — \eng{If people had sorted their waste, the dump would not have grown so fast.}
\item \eng{drought} (la sécheresse) — \eng{If there is a long drought, farmers will lose their crops.}
\item \eng{recycling} (le recyclage) — \eng{If everyone practised recycling, we would save natural resources.}
\item \eng{disease} (la maladie) — \eng{If the water is not boiled, children can catch a disease.}
\end{enumerate}
\end{corrige}

\begin{corrige}[Exercice 4 — Conjugaison à compléter]
\begin{enumerate}
\item \eng{If the government \textbf{invests} in hospitals, many lives \textbf{will be} saved.} (1st conditional ; passif : \eng{will be} + participe passé)
\item \eng{If I \textbf{had known} about the danger, I \textbf{would not have gone} there yesterday.} (3rd conditional ; participe passé de \eng{go} : \eng{gone})
\item \eng{If we \textbf{recycled} more, the city \textbf{would be} cleaner.} (2nd conditional)
\item \eng{If she \textbf{were} a doctor, she \textbf{would work} in a rural area.} (2nd conditional ; \eng{were} à toutes les personnes)
\end{enumerate}
\end{corrige}

\begin{corrige}[Exercice 5 — Traduction]
\begin{enumerate}
\item \eng{If I were you, I would plant trees.}
\item \eng{If we had recycled, the city would have been cleaner.}
\item \eng{If it rains tomorrow, the football match will be cancelled.}
\item \eng{If people did not throw plastic into the rivers, the fish would not die.}
\end{enumerate}
Points de vigilance : jamais de \eng{would} dans la subordonnée en \eng{if} ; le futur français « sera annulé » se rend par le passif \eng{will be cancelled} ; « si j'étais toi » se traduit toujours \eng{if I were you}.
\end{corrige}

\begin{corrige}[Exercice 6 — Transformation : les trois conditionnels]
\begin{enumerate}
\item 2nd conditional : \eng{If people recycled, the city would be clean.}
\item 3rd conditional : \eng{If people had recycled, the city would have been clean.}
\item Différence de sens : au 1st conditional, l'hypothèse est réelle et porte sur le futur (les gens peuvent encore recycler). Au 2nd conditional, elle est irréelle au présent (les gens ne recyclent pas, et l'on imagine le contraire). Au 3rd conditional, elle est irréelle au passé (les gens n'ont pas recyclé : il est trop tard, c'est un regret).
\end{enumerate}
\end{corrige}

\begin{corrige}[Exercice 7 — Texte à trous]
\eng{If the town council \textbf{built} more rubbish bins, the streets of Buea \textbf{would be} much cleaner.} (2nd conditional : hypothèse irréelle au présent)

\eng{... if the inhabitants \textbf{had not burnt} their waste in the open air, the air \textbf{would have been} less polluted.} (3rd conditional : regret sur l'année passée)

\eng{If every family \textbf{sorts} its rubbish next year, the town \textbf{will become} a model for the whole region.} (1st conditional : projet futur réaliste)

\eng{“If I \textbf{were} the mayor, I \textbf{would create} a recycling centre in every quarter.”} (2nd conditional : \eng{were} à toutes les personnes)
\end{corrige}

\begin{corrige}[Exercice 8 — Correction d'erreurs]
\begin{enumerate}
\item \eng{If it \textbf{rains} tomorrow, we \textbf{will stay} at home.} Trois fautes : \eng{will rain} $\rightarrow$ \eng{rains} (pas de \eng{will} après \eng{if}) ; \eng{would} $\rightarrow$ \eng{will} (cohérence du 1st conditional) ; \eng{stayed} $\rightarrow$ \eng{stay} (après \eng{will}/\eng{would}, base verbale sans \eng{-ed}).
\item \eng{If I \textbf{were} you, I would see a doctor.} Faute : \eng{would be} $\rightarrow$ \eng{were} (jamais \eng{would} dans la subordonnée ; \eng{were} à toutes les personnes).
\item \eng{If she had \textbf{gone} to the hospital earlier, she would have been saved.} Faute : \eng{went} $\rightarrow$ \eng{gone} (après \eng{had}, il faut le participe passé ; \eng{go — went — gone}).
\item \eng{If they \textbf{protect} the forest, the climate \textbf{will} improve.} Fautes : \eng{will protect} $\rightarrow$ \eng{protect} ; \eng{would} $\rightarrow$ \eng{will} (cohérence du 1st conditional). Autre correction possible au 2nd conditional : \eng{If they protected the forest, the climate would improve.}
\end{enumerate}
\end{corrige}

\begin{corrige}[Exercice 9 — Type Bac — Compréhension et langue]
\textbf{A. Comprehension questions}
\begin{enumerate}
\item They can avoid cholera by boiling their water before drinking it.
\item The cholera epidemic spread very fast because the broken pump was not repaired earlier.
\item Children will stop walking five kilometres every morning, and they will arrive at school on time.
\item Because a pump that nobody can repair is useless: trained technicians will be able to maintain and repair the boreholes.
\end{enumerate}
\textbf{B. Language work}
\begin{enumerate}
\item 1st conditional : \eng{If they build one borehole in every village, life will change completely.} 2nd conditional : \eng{If people boiled their water, we would have far fewer cases of cholera.} (ou \eng{If I were the minister of health, I would also train local technicians.}) 3rd conditional : \eng{If the authorities had repaired the broken pump earlier, the epidemic would not have spread so fast.}
\item \eng{If the authorities \textbf{had repaired} the pump earlier, the epidemic \textbf{would not have spread}.}
\item \eng{would have come} (3rd conditional : \eng{would have} + participe passé).
\end{enumerate}
Barème indicatif (sur 10) : A. 4 $\times$ 1,5 pt = 6 pts (réponse correcte, en anglais, avec ses propres mots) ; B. question 1 : 1,5 pt (0,5 pt par conditionnel correctement recopié) ; question 2 : 1,5 pt ; question 3 : 1 pt.
\end{corrige}

\begin{corrige}[Exercice 10 — Type Bac — Traduction]
\eng{If the government invested more in hospitals, many lives would be saved. Last year, if we had protected the forest, the floods would not have destroyed the crops. Tomorrow, if every citizen plants a tree, our country will become greener.}

Points de vigilance : \eng{invest in} (investir dans) ; passif \eng{would be saved} ; \eng{had protected} (past perfect) et \eng{would not have destroyed} au 3rd conditional ; article défini \eng{the floods} ; pas de futur après \eng{if} dans la dernière phrase (\eng{plants}, pas \eng{will plant}).

Barème indicatif (sur 6) : 2 pts par phrase (1 pt pour la structure conditionnelle correcte, 1 pt pour le vocabulaire et l'orthographe).
\end{corrige}

\begin{corrige}[Exercice 11 — Type Bac — Expression écrite]
\textbf{Proposition de rédaction modèle} (environ 80 mots ; les deux structures conditionnelles sont soulignées dans la copie) :

\eng{If everyone protected the environment, our town would be a much healthier place to live in. In my school, for example, students throw plastic bags on the ground every day. If we had learnt to sort our waste earlier, the schoolyard would not have become so dirty. But it is not too late: if every class plants two trees this year, the school will be greener and shadier in a few years. If I were the headmaster, I would also organise a cleaning day every month. Small actions can change our lives.}

Traduction : « Si tout le monde protégeait l'environnement, notre ville serait un endroit bien plus sain où vivre. Dans mon école, par exemple, les élèves jettent des sachets en plastique par terre tous les jours. Si nous avions appris à trier nos déchets plus tôt, la cour ne serait pas devenue si sale. Mais il n'est pas trop tard : si chaque classe plante deux arbres cette année, l'école sera plus verte et plus ombragée dans quelques années. Si j'étais le proviseur, j'organiserais aussi une journée de nettoyage chaque mois. De petites actions peuvent changer nos vies. »

Points de vigilance : la phrase imposée est au 2nd conditional (\eng{protected} / \eng{would be}) ; le modèle ajoute un 3rd conditional (\eng{had learnt} / \eng{would not have become}) et un 1st conditional (\eng{plants} / \eng{will be}) ; vérifier l'absence de \eng{will} ou \eng{would} après \eng{if}.

Barème indicatif (sur 10) : respect de la consigne et de la longueur (2 pts) ; deux conditionnels différents correctement formés (4 pts) ; richesse du vocabulaire lié à l'environnement (2 pts) ; orthographe, grammaire et ponctuation (2 pts).
\end{corrige}

\newpage

\coursec{Fiche bilan}

\begin{popfigure}
\begin{center}
\begin{tikzpicture}[
  mindmap,
  grow cyclic,
  every node/.style={concept, font=\small},
  root concept/.append style={concept color=popblue, font=\bfseries\small, minimum size=3.2cm},
  level 1 concept/.append style={level distance=3.6cm, sibling angle=60, font=\scriptsize},
  text=white
]
\node[root concept, align=center]{Conditionals\\ 1st, 2nd, 3rd}
  child[concept color=popgreen]{ node[align=center]{1st : réel\\ \eng{If + present,}\\ \eng{will + BV}} }
  child[concept color=poporange]{ node[align=center]{2nd : irréel\\ \eng{If + past,}\\ \eng{would + BV}} }
  child[concept color=poppurple]{ node[align=center]{3rd : regret\\ \eng{If + past perfect,}\\ \eng{would have + V-en}} }
  child[concept color=poppink]{ node[align=center]{Variantes\\ \eng{unless,}\\ \eng{provided that,}\\ \eng{in case}} }
  child[concept color=popgold]{ node[align=center]{Pièges\\ jamais \eng{will}\\ après \eng{if}} }
  child[concept color=popblue!70]{ node[align=center]{Emplois\\ conseil, rêve,\\ reproche,\\ avertissement} };
\end{tikzpicture}
\end{center}
\legende{Carte mentale : les trois conditionnels en un coup d'œil}
\end{popfigure}

\begin{bilan}
\textbf{1. Lexique clé (English / Français)}

\begin{center}
\begin{tabularx}{\linewidth}{|X|X|}
\hline
\rowcolor{popblueL} \textbf{English} & \textbf{Français} \\
\hline
\eng{unless} & à moins que / sauf si \\
\hline
\eng{provided (that) / as long as} & à condition que / pourvu que \\
\hline
\eng{in case} & au cas où \\
\hline
\eng{otherwise} & sinon \\
\hline
\eng{to pollute / pollution} & polluer / la pollution \\
\hline
\eng{to protect the environment} & protéger l'environnement \\
\hline
\eng{to recycle waste} & recycler les déchets \\
\hline
\eng{to regret} & regretter \\
\hline
\eng{if I were you} & si j'étais toi (à ta place) \\
\hline
\end{tabularx}
\end{center}

\textbf{2. Règles à connaître par cœur}

\begin{center}
\begin{tabularx}{\linewidth}{|l|X|X|}
\hline
\rowcolor{popblueL} \textbf{Type} & \textbf{Structure} & \textbf{Exemple} \\
\hline
1st (réel) & \eng{If + present simple, will + BV} & \eng{If it rains, we will stay at home.} \\
\hline
2nd (irréel présent) & \eng{If + past simple, would + BV} & \eng{If I had money, I would plant trees.} \\
\hline
3rd (irréel passé) & \eng{If + past perfect, would have + participe passé} & \eng{If we had recycled, the city would have been cleaner.} \\
\hline
\end{tabularx}
\end{center}

\begin{itemize}
  \item Après \cle{if} (condition), on ne met \textbf{jamais} \eng{will} ni \eng{would} : on dit \eng{If it rains}, pas \eng{If it will rain}.
  \item Au 2nd conditional, \eng{were} s'emploie pour toutes les personnes : \eng{If I were rich...}, \eng{If she were here...}
  \item Les deux propositions peuvent s'inverser ; la virgule disparaît quand \eng{if} est au milieu : \eng{We will stay at home if it rains.}
  \item \eng{Unless} = \eng{if... not} : \eng{Unless you hurry, you will miss the bus.} (= \eng{If you don't hurry...})
  \item \eng{Would} se contracte en \eng{'d} : \eng{I'd go.} ; attention à ne pas le confondre avec \eng{had} (\eng{I'd been} = \eng{I had been}).
\end{itemize}

\textbf{3. Méthode express}
\begin{enumerate}
  \item Repérer le temps après \eng{if} : present $\rightarrow$ 1st ; past $\rightarrow$ 2nd ; past perfect $\rightarrow$ 3rd.
  \item Vérifier le sens : hypothèse possible (1st), imaginaire (2nd), regret sur le passé (3rd).
  \item Construire la principale : \eng{will + BV} / \eng{would + BV} / \eng{would have + V-en}.
  \item Relire : pas de \eng{will}/\eng{would} dans la subordonnée, participe passé correct au 3rd.
\end{enumerate}
\end{bilan}

\begin{aretenir}[Checklist avant le Bac]
\begin{itemize}
  \item[$\square$] Je connais la structure du 1st conditional : \eng{If + present, will + BV}.
  \item[$\square$] Je connais la structure du 2nd conditional : \eng{If + past, would + BV}.
  \item[$\square$] Je connais la structure du 3rd conditional : \eng{If + past perfect, would have + participe passé}.
  \item[$\square$] Je ne mets jamais \eng{will} ou \eng{would} juste après \eng{if}.
  \item[$\square$] J'utilise \eng{were} avec toutes les personnes au 2nd conditional (\eng{If I were you...}).
  \item[$\square$] Je sais transformer \eng{unless} en \eng{if... not} et inversement.
  \item[$\square$] Je place la virgule correctement quand la phrase commence par \eng{if}.
  \item[$\square$] Je connais les participes passés irréguliers utiles (\eng{been, done, gone, had, made, taken}).
  \item[$\square$] Je sais exprimer un conseil avec \eng{If I were you, I would...}.
  \item[$\square$] Je sais exprimer un regret avec le 3rd conditional (\eng{If I had studied harder, I would have passed.}).
  \item[$\square$] Je peux écrire un court paragraphe sur l'environnement en utilisant au moins deux conditionnels différents.
  \item[$\square$] Je distingue \eng{'d} = \eng{would} et \eng{'d} = \eng{had} selon le verbe qui suit.
\end{itemize}
\end{aretenir}

\findecours

\end{document}
